%%%%%%%%%%%%%%%%%%%%%%%%%%%%%%%%%%%%%%%%%%%%%%%%%%%%%%%%%%%%%%%%%%%%%
%% Submission template for the Journal of Computational Physics
%% (Elsevier). Built on the els-cas-templates "cas-sc" single-column
%% class used by Elsevier's CAS workflow.
%%%%%%%%%%%%%%%%%%%%%%%%%%%%%%%%%%%%%%%%%%%%%%%%%%%%%%%%%%%%%%%%%%%%%
\documentclass[a4paper,fleqn]{cas-sc}

% Use this option instead if the frontmatter overflows one page:
%\documentclass[a4paper,fleqn,longmktitle]{cas-sc}

%%%%%%%%%%%%%%%%%%%%%%%%%%%%%%%%%%%%%%%%%%%%%%%%%%%%%%%%%%%%%%%%%%%%%
%% Bibliography handling -- JCP uses numbered references.
%%%%%%%%%%%%%%%%%%%%%%%%%%%%%%%%%%%%%%%%%%%%%%%%%%%%%%%%%%%%%%%%%%%%%
\usepackage[numbers,sort&compress]{natbib}

%%%%%%%%%%%%%%%%%%%%%%%%%%%%%%%%%%%%%%%%%%%%%%%%%%%%%%%%%%%%%%%%%%%%%
%% Math, algorithms, and project-specific packages
%%%%%%%%%%%%%%%%%%%%%%%%%%%%%%%%%%%%%%%%%%%%%%%%%%%%%%%%%%%%%%%%%%%%%
\usepackage{lyu}
\usepackage{algorithm}
\usepackage{algpseudocode}%[noend]
\usepackage{amsthm}
% \mathchardef\mhyphen="2D
\newcommand{\mhyphen}{\text{-}}

\newtheorem{theorem}{Theorem}
\newtheorem{lemma}[theorem]{Lemma}
\theoremstyle{definition}
\newtheorem{definition}[theorem]{Definition}
\theoremstyle{remark}
\newtheorem{remark}[theorem]{Remark}
\newtheorem{proposition}{Proposition}
\newtheorem{assumption}{Assumption}
\newtheorem{corollary}{Corollary}

\begin{document}
\let\WriteBookmarks\relax
\def\floatpagepagefraction{1}
\def\textpagefraction{.001}

\shorttitle{Stochastic branching particle method}
\shortauthors{L.~Lyu and H.~Lei}

\title[mode=title]{A stochastic branching particle method for solving non-conservative reaction--diffusion equations}

\author[1,2]{Liyao Lyu}
\ead{lyuliyao@math.ucla.edu}

\author[2,3]{Huan Lei}
\ead{leihuan@msu.edu}

\affiliation[1]{organization={Department of Mathematics, University of California, Los Angeles},
            addressline={520 Portola Plaza},
            city={Los Angeles},
            state={CA},
            country={USA}}

\affiliation[2]{organization={Department of Computational Mathematics, Science \& Engineering, Michigan State University},
            addressline={428 S Shaw Ln},
            city={East Lansing},
            state={MI},
            country={USA}}

\affiliation[3]{organization={Department of Statistics \& Probability, Michigan State University},
            addressline={619 Red Cedar Road},
            city={East Lansing},
            state={MI},
            country={USA}}

\begin{abstract}
We propose a stochastic branching particle-based method for solving nonlinear non-conservative advection–diffusion–reaction equations. 
The method splits the evolution into an advection–diffusion step, based on a linearized Kolmogorov forward equation and approximated by stochastic particle transport, and a reaction step implemented through a branching birth–death process that provides a consistent temporal discretization of the underlying reaction dynamics. This construction yields a mesh-free, nonnegativity-preserving scheme that naturally accommodates non-conservative systems and remains robust in the presence of singularities or blow-up. 
We validate the method on two representative two-dimensional systems: the Allen–Cahn equation and the Keller–Segel chemotaxis model. In both cases, the present method accurately captures nonlinear behaviors such as phase separation and aggregation, and achieves reliable performance without the need for adaptive mesh refinement.

\end{abstract}

\begin{keywords}
particle methods \sep
stochastic branching \sep
reaction-diffusion \sep
chemotaxis \sep
aggregation
\end{keywords}

\maketitle

\section{Introduction}
The advection--diffusion--reaction (ADR) type partial differential equations (PDEs) have attracted much interest in physics, biology, and engineering for modeling various phenomena relevant to chemical reactions, transport, and collective behaviors in biological systems \cite{hundsdorfer2013numerical}. Traditional grid-based methods, such as finite difference~\cite{hoff1978stability}, finite element~\cite{codina1998comparison}, and discontinuous Galerkin schemes~\cite{ayuso2009discontinuous,xia2009application}, have achieved great success in solving many of these problems. However, from a numerical perspective, an essential challenge in solving such PDEs lies in handling nonlinearities and singular structures in certain cases, e.g., the aggregation phenomena of the Keller--Segel (KS) model \cite{keller1970initiation} and the interfacial dynamics of the Allen--Cahn (AC) model \cite{Allen_Cahn_1979}.  
Mesh-based discretizations approximate solutions in function spaces that may not be well-suited for capturing the blow-up or highly localized behavior, often resulting in reduced accuracy or numerical instabilities, and requiring additional numerical treatments.


Particle-based approaches offer a natural alternative by representing the solutions as empirical measures, which are particularly effective for problems exhibiting concentration or blow-up phenomena. 
Existing particle methods, such as the particle-in-cell method~\cite{birdsall2018plasma, hockney2021computer} for solving the Vlasov--Poisson equation, are often based on well-established  underlying physical laws. In this setting, particle dynamics can be derived from the explicit micro-scale governing equations. In contrast, the explicit physical laws are often unknown for general non-conservative ADR type PDEs, where effective interactions may be highly nonlinear, system-specific, or even not well characterized. For example, in many biological and chemical systems such as chemotaxis, the micro-scale interactions remain poorly understood, which limits the applicability of traditional particle-based approaches.
To extend particle-based methods beyond physically motivated settings, blob methods~\cite{craig2016blob,carrillo2019blob,huang2017error} and semi-Lagrangian particle schemes~\cite{qiu2010conservative,qiu2011conservative,qiu2011positivity} have been developed for conservative PDEs.
These approaches are based on the mean-field perspective, in which the empirical measures of interacting particle systems converge to the solution of the original PDE. This mathematical framework establishes a rigorous link between the particle dynamics and the continuum equations and does not rely on explicit microscale physical laws. Despite these advances, there is still no unified particle-based framework for systematically discretizing nonlinear and non-conservative PDEs, especially in the presence of singularities, sharp interface or blow-up phenomena.

Machine learning-based numerical PDE algorithms, such as Deep Galerkin Methods~\cite{sirignano2018dgm}, Deep Ritz method~\cite{yu2018deep}, and Mixed Residual method~\cite{lyu2022mim}, have attracted significant attention in recent years. By leveraging the strong representability of neural networks, these methods are particularly promising for solving high-dimensional problems. In the context of particle-based approaches, coupling neural networks with the Feynman-Kac formula~\cite{han2024deep, hutzenthaler2021speed, nguwi2024deep,beck2021deep} or backward stochastic differential equations (SDEs) ~\cite{han2017deep, han2018solving, hure2020deep} has also emerged as a powerful strategy for solving a broad class of PDEs. Nevertheless, the approximation capability of neural networks is fundamentally grounded in the universal approximation theorem that targets continuous solutions on compact sets.
As a result, neural-network–based solvers may face intrinsic difficulties when dealing with nonlinear, non-conservative PDEs that develop singularities or blow-up phenomena, where discontinuities or concentration effects play a central role.

In this work, we propose a \textit{stochastic branching particle method} to address these challenges. 
At each time step, the original PDE is locally linearized and approximated by the Kolmogorov forward equation, which establishes a direct connection between the PDE and the stochastic particle dynamics.  Unlike approaches that rely on explicit physical laws or mean-field limit arguments, this framework establishes a systematic particle-based discretization applicable to general nonlinear and non-conservative PDEs. 
%
A key feature of the method is a particle-level birth–death mechanism, introduced to represent non-conservative reaction terms by allowing the number of particles to evolve consistently with the underlying reaction dynamics. The only step that invokes a continuity assumption is the interpolation of the empirical measure, which can be replaced by blob- or kernel-based representations. Consequently, the developed method does not fundamentally rely on continuity and is well-suited for problems with singularities or blow-up phenomena. 
%
From a numerical perspective, the proposed method is broadly compatible with existing stochastic interpretations of reaction–diffusion systems. For example, in the case of chemotaxis, it aligns with the mean-field derivation of the KS equation from moderately interacting particles developed by Stevens~\cite{Stevens_SIAM_2000}. However, our formulation is not restricted to this setting and applies generally to nonlinear ADR systems with non-conservative terms.

% From a stochastic particle perspective, the algorithm is naturally consistent with the physical intuition behind Keller–Segel dynamics, in line with Stevens' mean-field derivation of the Keller–Segel equations from moderately interacting particle systems \cite{Stevens_SIAM_2000} and can be applied naturally to a broader class of advection–diffusion–reaction systems.
% formulation is closely related to Stevens’ mean-field derivation of Keller--Segel equations from moderately interacting particle systems \cite{Stevens_SIAM_2000}. Unlike  \cite{Stevens_SIAM_2000}, we develop a finite-$N$ particle discretization based on local Kolmogorov linearization and a particle-level birth-death update, which directly handles the nonlinear and nonconservative terms and extends to a class of advection–diffusion–reaction systems.  


To illustrate the effectiveness of the proposed method, we consider two representative non-conservative reaction–diffusion systems: the AC equation for the phase separation with interfacial motion and nonlinear reactions,  
and the KS chemotaxis model for the aggregation-driven blow-up dynamics~\cite{keller1970initiation, patlak1953random}. 
While the AC equation serves to demonstrate the generality of our method, the KS model provides a more challenging test problem and is the primary focus of the present study. In particular, we consider the two-dimensional KS chemotaxis model as the theme problem, which takes the form 
\begin{equation}\label{equ:KS}
\begin{aligned}
\partial_t u - \nabla \cdot ( \nabla u - \chi(u) \nabla v) &= f_u(u,v), & \bx\in \mathbb{T}^2 , t\in(0,T), \\
\partial_t v - \Delta v &= f_v(u,v), & \bx\in \mathbb{T}^2 , t\in(0,T),
\end{aligned}
\end{equation}
where $u$ and $v$ represent the cell density and chemical concentration, respectively. $\chi$ represents the chemotactic sensitivity, and $f_u$ and $f_v$ represent the production or decay rate.  For simplicity, we focus on periodic boundary conditions or unbounded domains, though the framework can be extended to more general boundary settings in future research.  
%The KS system~\eqref{equ:KS} has been extensively studied over the past several decades, with numerous works addressing both analytical properties and numerical methods across model variants and dimensions. 
In two dimensions, it is well understood that both parabolic–elliptic and fully parabolic cases can exhibit finite-time blow-up to a Dirac-delta profile when the initial mass is sufficiently large~\cite{herrero1997blow, perthame2007transport}. Furthermore, in the three-dimensional parabolic–elliptic case, an alternative self-similar blow-up profile of the form $ C(T-t+\vert \mb x\vert ^2)^{-1}$ has also been identified in the literature~\cite{giga2011asymptotic,souplet2019blow}. 

On the numerical side, many notable methods have been developed for the KS systems. The finite-volume method~\cite{filbet2006finite,chertock2008second,chertock2018high} achieves efficient and positivity-preserving simulations for chemotaxis and haptotaxis models.  
%
Shen et al.~\cite{shen2020unconditionally} introduced an energy-dissipative and bound-preserving formulation, subsequently extended by Chen et al.~\cite{chen2022error} to a fully discrete finite element scheme for the 2D parabolic–elliptic KS, which reproduces finite-time blow-up under assumptions analogous to the continuous setting. 
Liu and Wang~\cite{liu2018positivity} reformulated the equations via the Le Ch\^atelier principle to construct positivity-preserving schemes. Discontinuous Galerkin methods have also been applied to KS systems~\cite{zhong2024direct, li2017local,epshteyn2009new}, in addition to various mesh-based finite element, finite volume, and domain-decomposition techniques for 3D fully parabolic models~\cite{epshteyn2019efficient, strehl2013positivity}. For particle-based approaches, the most relevant work is the recently proposed stochastic interacting particle–field algorithm~\cite{wang2025novel,hu2024stochastic,chertock2025hybrid}. However, this approach does not deal with the non-conservative equations at the particle level, limiting its applicability to the broader class of nonlinear non-conservative PDEs. 

We address this challenge by explicitly coupling a SDE-based advection--diffusion process with a particle-level birth--death mechanism for reactions, which overcomes the limitations of traditional methods and yields a particle measure discretization that remains robust near concentration and blow-up. We validate the method on two-dimensional KS benchmarks, including near blow-up and nonlinear growth, and compare with the standard grid-based method, observing the expected $1/\sqrt{N}$ scaling with the particle number.
%
The remainder of the paper is organized as follows: Section 2  and 3 present the stochastic branching method for a non-conservative ADR equation and the convergence analysis; Section 4 extends the approach for the KS systems; Section 5 reports numerical results and Section 6 discusses extensions and outlook.

% This work makes three contributions. (\Rmnum{1}) We introduce a stochastic branching framework that couples SDE-based advection–diffusion with a birth–death mechanism for reactions, yielding a particle–measure discretization that is not tied to smooth function spaces and remains effective near concentration and blow-up. (\Rmnum{1}) We derive the scheme via local Kolmogorov linearization and an operator splitting, giving a transparent PDE–SDE correspondence and an implementable algorithm with Fourier/kernel reconstructions of empirical measures. (\Rmnum{1}) We validate the method on two-dimensional KS benchmarks—including near-blow-up and nonlinear growth—against standard grid-based baselines, observing the expected $1/\sqrt{N}$ scaling with particle number. We also discuss extensions to other reaction–diffusion systems (Gray–Scott, Allen–Cahn, Schnakenberg) and paths toward higher-dimensional settings, with details deferred to the Discussion


% In summary, we propose a novel stochastic branching method for solving the Keller–Segel chemotaxis model, which overcomes the limitations of traditional numerical methods by efficiently handling singularities and nonlinearities at particle level. A heuristic derivation of the algorithm is presented in the current version and will be discussed in detail in future work. Although our focus here is on the two-dimensional Keller–Segel model, the proposed framework is general and can be applied to a broad class of reaction–diffusion systems, such as the Gray–Scott model, the Allen–Cahn equation, and Schnakenberg pattern formation. Furthermore, by leveraging recent advances in density estimation techniques~\cite{novikov2021tensor, tang2025wavelet, sheng2025approximation, tang2024solving}, our approach can be naturally extended to higher-dimensional systems. The effectiveness of the method is demonstrated through several numerical examples, highlighting both its robustness and potential for generalization.

\section{Stochastic branching method}
Let us start with the following scalar non-conservative advection-diffusion–reaction system on the $d$-dimensional torus $\mathbb{T}^d$ for a time $t \in (0, T]$:
\begin{subequations}\label{equ:reaction-diffusion}
\begin{align}
\partial_t u &= \nabla \!\cdot\! \big( a(u) u \big) + D \Delta u + r(u) u, & \text{for } (x, t) \in \mathbb{T}^d \times (0, T], \label{equ:pde} \\
u(x, 0) &= u_0(x), & \text{for } x \in \mathbb{T}^d. \label{equ:ic}
\end{align}
\end{subequations}
Here, $u(x,t)$ is the scalar field, $a(u): \mathbb{R}_{+} \to \mathbb{R}^d$ is the advective velocity, $D > 0$ is the diffusion coefficient, and $r(u): \mathbb{R}_{+} \to \mathbb{R}$ is the reaction rate.
Let $S(N,t)=\{\bX^N_i\}_{i\in\mathcal I(N,t)}$ denote a set of particles at time $t$ whose initial population is $N$, where the particles' indices are numbered consecutively in the index set $\mathcal I(N,t)$ and $\bX^N_i(t)$ denote the position of the $i\mhyphen$th particle. The associated empirical (finite) measure is
\begin{equation}
    t \to \mu^N(t,\intd \bx) =  \frac{1}{N}\sum_{i\in \mathcal I(N,t)}\delta_{\bX^N_i(t)},
\end{equation}
where $\delta_\bx $ denotes the Dirac measure at $\bx\in \mathbb{T}^d$. The core idea of particle-based methods is to approximate the PDE solution by the limit of the empirical measure
\[
\mu(t,\mathrm{d}\bx) = \lim_{N\to\infty} \mu^N(t,\mathrm{d}\bx).
\]
Since Eq. \eqref{equ:reaction-diffusion} is non-conservative, the total mass of the system is not preserved. Consequently, both $\mu^N$ and its limit $\mu$ should be viewed as finite measures rather than probability measures.


\subsection{Initial condition and Normalization}
To connect the solution $u(t,\bx)$ to the measure $\mu(t,\mathrm{d}\bx)=\rho(t,\bx) \intd \bx$, we introduce a linear transformation between them. 
Without loss of generality, we assume that the initial measure $\mu(0,\cdot)$ is a probability measure with total mass $1$, and define
\begin{equation*}
    u(t,\bx)\,\mathrm{d}\bx = Z \,\mu(t,\mathrm{d}\bx), \quad 
Z = \int_{\Omega} u(0,\bx)\,\mathrm{d}\bx,
\end{equation*}
where $Z$ denotes the total initial mass.
Thus $\rho(0,\cdot)$ is a probability density
\[
\rho(0,\mathbf{x}) \;=\; \frac{u(0,\mathbf{x})}{Z},
\]
and we can sample the initial particle positions independently from $\rho(0,\cdot)$ (e.g., via Metropolis–Hastings) to form $S(N,0)=\{\mathbf{X}^i_0\}_{i=1}^N$. 
Substituting $u=Z\rho$ into~\eqref{equ:reaction-diffusion}  yields
\begin{equation}\label{equ:normalized_diffusion_reaction}
\partial_t \rho 
= \nabla \cdot \left( \tilde a(\rho)\rho \right)+ D \Delta\rho 
  + \tilde r(\rho)\rho,
\end{equation}
with rescaled coefficients
\[
\tilde a(\rho) = a(Z\rho),  
\qquad \tilde r(\rho) = r(Z\rho).
\]
Since reactions generally change the total mass, $\int_{\Omega}\rho(t,\mathbf{x})\,\mathrm{d}\mathbf{x}$ may not equal $1$ for $t>0$. Rather, it tracks the mass of $u$ via $\int_{\Omega} u(t,\mathbf{x})\,\mathrm{d}\mathbf{x} = Z \int_{\Omega} \rho(t,\mathbf{x})\,\mathrm{d}\mathbf{x}$.

\subsection{Operator Splitting}
To approximate the solution of the PDE, we discretize the time interval $[0,T]$ with step size $\tau>0$ and grid points
\[
0 = t_0 < t_1 < \cdots < t_M = T,\qquad \tau=t_{n+1}-t_n.
\]
Then an operator-splitting strategy is used to separate the dynamics according to their physical mechanisms. 
Equation~\eqref{equ:normalized_diffusion_reaction} can be written as
\begin{equation}\label{eq:rho-pde}
\partial_t \rho = \mathcal{A}(\rho) + \mathcal{B}(\rho) ,
\end{equation}
where the operators $\mathcal{A}$ and $\mathcal{B}$ correspond to
\begin{align}
\mathcal{A}(\rho) =& \nabla \cdot \left( \tilde a(\rho)\rho \right) + D\Delta\rho  ,\label{eq:operator-A} \\
\mathcal{B}(\rho) =& \tilde r(\rho)\rho \label{eq:operator-B}.
\end{align}
Here, $\mathcal{A}$ describes the advection--diffusion dynamics, while $\mathcal{B}$ accounts for the reaction mechanism (e.g., proliferation or death).  The evolution of the exact solution over $t$ can be expressed formally as
\begin{equation}\label{}
\rho(t,\bx) = \exp\left( t (\mathcal{A} + \mathcal{B}) \right)\rho(0,\bx).
\end{equation}
In practice, the operator $\exp\!\big(\tau(\mathcal{A}+\mathcal{B})\big)$ cannot be evaluated explicitly, 
we adopt the first-order Lie-Trotter splitting \cite{trotter1959product}, which is justified by the Trotter product formula:
\begin{equation}
\exp\!\left( t (\mathcal{A} + \mathcal{B}) \right) 
= \lim_{M\to\infty} \Big( \exp\!\left(\tfrac{t}{M} \mathcal{B}\right) 
                          \exp\!\left(\tfrac{t}{M} \mathcal{A}\right) \Big)^M .
\end{equation}
At the discrete level, we approximate $\rho(t_{n+1},\cdot)$ by
\begin{equation}
\rho_{n+1} = \exp(\tau \mathcal{B}) \exp(\tau \mathcal{A}) \rho_n ,
\end{equation}
for $n=0,\cdots,M-1$.
This scheme provides a simple and efficient first-order approximation; 
higher-order splittings, such as Strang splitting \cite{strang1968construction}, may be employed for improved accuracy.
In practice, this decomposition allows us to design particle-based updates for each subprocess separately:
\begin{equation}\label{eq:lie-trotter}
\rho_n^\ast = \exp(\tau \mathcal{A})\rho_n, 
\qquad 
\rho_{n+1} = \exp(\tau \mathcal{B})\rho_n^\ast.
\end{equation}

\subsection{Advection Diffusion Step}
\label{sec:adevection_diffusion}
In this step, we propagate particles under the advection–diffusion dynamics, where each particle follows a stochastic path determined by both a drift and a diffusion term. At time step $n$, let $S^N_n = \{\bX^i_{N,n}\}_{i\in \mathcal I_n}$ denote  a set of particles  whose initial population is $N$, which induce the empirical measure 
\[
\mu^N_{n}= \frac{1}{N}\sum_{i\in \mathcal I_n}\delta_{\bX_{N,n}^i}
\] 
approximates the solution at step $n$, i.e.,  $\mu^N_{n} (\intd \bx) \approx \mu_n(\intd \bx) = \rho_n(\bx)\intd \bx$. 
Our goal is to evolve these particles according to the transition probability $p(\bx,s;\by,t)$, which leads to an updated empirical measure
\[\mu_n^{N,*} = \int p(\bx,t_n;\by,t_n+\tau)  \mu_n^N(\intd \bx) =   \frac{1}{N}\sum_{i\in \mathcal I_n}\delta_{\bX_{N,n}^{i,*}}
\]  
reflecting the action of the advection--diffusion operator on the density.
By the Kolmogorov forward equation, the transition density $p(\bx,s;\by,t)$ satisfies 
\begin{equation}
    \partial_t p = \tilde {\mathcal A }p := - \nabla \cdot(b  p) + \frac{1}{2}\nabla^2:(\sigma \sigma^\top p) 
\end{equation}
with initial condition $\lim_{t\to s^+} p(\bx,s;\cdot,t) = \delta_\bx$. By matching the differential operator $\mathcal A$ with the target PDE \eqref{equ:reaction-diffusion}, we have 
\begin{equation}
    b(x) = -\tilde a(\rho^N_n(x)), \quad \frac{1}{2}\sigma \sigma^\top = D,
\end{equation}
where $\rho^N_n(x)$ represents the interpolated probability density of $\mu_n^N$ and will be discussed in Sec. \ref{sec:interpolation}.
%
Accordingly, the evolution of each particle $\bX^i_n$ as a stochastic process satisfying the It\^o stochastic differential equation
\begin{equation}\label{equ:overd_langevin}
    \intd \bX_{N,n}^i = b(\bX_{N,n}^i) \intd t  + \sigma \intd \bW^N_{i}(t),
\end{equation}
where $\bW^i_{N}(t)$ is a $2\mhyphen$dimensional Brownian motion. 
In practice, we approximate Eq. \eqref{equ:overd_langevin} using an Euler-Maruyama scheme with step size $\tau$, 
\begin{equation}
 \bX_{N,n}^{i,*} = \mathbf{X}_{N,n}^{i} + b(\bX_{N,n}^i) \tau  + \sigma \sqrt{\tau} \boldsymbol{\xi}^i_{N,n},
 \label{eq:euler_maruyama}
\end{equation}
where $\boldsymbol{\xi}^i_{N,n}$ represent identically independent distributed normal random variables; see Algorithm \ref{alg:sde} for the summary of the particle propagation for the advection--diffusion step. 

\begin{algorithm}
\caption{SDE propagation of the advection--diffusion step}\label{alg:sde}
\begin{algorithmic}
\Require Particle set: $S=\{\mathbf{X}_i\}_{i\in\mathcal I}$; time step $\tau>0$
\Require Drift term $b(\cdot)$
\Require Diffusion strength $\sigma$
\Ensure Propagated set $S_*=\{\mathbf{X}^{i,*}\}_{i\in\mathcal I}$
\Procedure{SDE}{$S,\,\tau,\,b,\,\sigma$}
  \State $S_{*}\gets [\ ]$
  \For{$i\in\mathcal I$}
    \State Sample $\boldsymbol{\xi}_i \sim \mathcal N(0,I_d)$ independently
    \State $\mathbf{X}^*_i \gets \mathbf{X}_i \;+\; b(\mathbf{X}_i)\,\tau \;+\; \sqrt{\tau}\,\sigma\,\boldsymbol{\xi}_i$
    \State Append $\mathbf{X}^{*}_i$ to $S_{*}$
  \EndFor
  \State \Return $S_{*}$
\EndProcedure
\end{algorithmic}
\end{algorithm}

%======================================================================
%  Particle-population growth: exact finite-N identity + mean-field limit
%======================================================================
For a quantitative description of the particle population, let
$N(t)=\lvert\mathcal I(N,t)\rvert$ denote the number of particles alive at
time $t$, and let
\[
   M^N_n := \big\langle 1,\,\mu^N_n\big\rangle
   =\frac{N(t_n)}{N}
\]
be the corresponding empirical mass. Let $\mathcal F_n$ be the
$\sigma$-algebra of all information available up to time $t_n$, i.e.\ before the
step-$n$ advection--diffusion update; the post-transport information, which
additionally fixes $\mu^{N,*}_n$, is the larger $\sigma$-algebra $\mathcal G_n$
introduced in the error analysis (eq.~\eqref{eq:filtrations}).

\subsection{Branching Step}
\label{sec:branching}
This step modifies the mass of the empirical measure and change the particle number. Given the intermediate empirical measure from the advection--diffusion step $\mu_{N,n}^*$, we aim to evolve this measure under the action of the reaction operator $\mathcal  B$ over a time step $\tau$ following
\[
\partial_t \rho = c_n(\mathbf{x})\,\rho,
\qquad c_n(\mathbf{x}) := \tilde r\!\big(\rho_n^\ast(\mathbf{x})\big),
\]
such that the resulting empirical measure
\begin{equation}
    \mu^N_{n+1} = \frac{1}{N} \sum_{i\in \mathcal I (N,(n+1)\tau)} \delta_{\bX^N_{i,n+1}}
\end{equation}
approximates the updated density $\mu^N_{n+1}(\intd \bx ) = \rho^N_{n+1}(\bx) \intd \bx \approx {\rm e}^{c_n(\mathbf{x}) \tau} \rho_{n}^{N,*}(\bx) \intd \bx$. Therefore, the particles located at $\bx$ should birth and death with probability $c(\bx)=\tilde r(\rho^*(\bx))$ depending on the positiveness of $\tilde r$.
To implement this step, we adopt a branching process. For each particle $\bX^{N,\ast}_{i,n}$, we associate a branching rule determined by $c_i:=c_n(\mathbf X^{N,\ast}_{i,n})$. Specifically, over the interval $[t_n,t_{n+1}]$, each particle undergoes one of the following stochastic events:
\begin{itemize}
  \item \emph{If \(c_i\ge 0\) (birth):} keep the parent and draw \(K_i\sim \mathrm{Poisson}\!\big(e^{c_i\tau}-1\big)\); output \(1+K_i\) identical copies at \(\mathbf X^{N,\ast}_{i,n}\).
  \item \emph{If \(c_i<0\) (death):} keep the parent with probability \(p_i=e^{c_i\tau}\), otherwise remove it.
\end{itemize}
% \begin{itemize}
%     \item Survival (no branching): particle $\bX^{i,*}_{N,n}$ survives with probability approximately $\exp(-| c(\bX^{i,*}_{N,n})|\tau)$.
%     \item Death (removal): if $c<0 $, particle $\bX^{i,*}_{N,n}$ dead  with the probability $1-\exp( c(\bX^{i,*}_{N,n})\tau)$.
%     \item Birth (duplication ): if $ c>0 $, particle $\bX^{i,*}_{N,n}$ duplication with the probability $1-\exp(- c (\bX^{i,*}_{N,n})\tau)$.
% \end{itemize}

Let the resulting set be \(S_{n+1}^N=\{\mathbf X^N_{i,n+1}\}_{i\in\mathcal I_{n+1}^N}\) and
\[
\mu_{n+1}^{N}(\mathrm{d}\mathbf{x})
\;=\;
\frac{1}{N}\sum_{i\in\mathcal I_{n+1}^N}\delta_{\mathbf X^N_{i,n+1}}(\mathrm{d}\mathbf{x}).
\]


\begin{proposition}
    Let $S=\{\bX_i\}_{i\in \mathcal I}$ be a set of particles, whose empirical measure is defined as $\mu (\intd \bx) = \frac{1}{N}\sum_{i\in \mathcal I }\delta_{\bX_i}$. If the particles are updated according to the above rules then the updated particle set $S_+=\{ \bX_{i,+}\}_{i\in {\mathcal I}_+}$ and corresponding  measure $ \mu_+ (\intd \bx) = \frac{1}{N}\sum_{i\in  {\mathcal  I}_+}\delta_{ \bX_{i,+}}$ should satisfies 
    \begin{equation}
        \mathbb E [\left<f, \mu_+\right>] =   \left<f,\exp( c\tau)\mu\right>,
    \end{equation}
    for all $f\in C^2_b(\mathbb R)$.
\label{prop:stochastic_branch}
\end{proposition}

\begin{proof} Let $\lambda_i = {\rm e}^{c(\mb X^i) \tau} -1 $ denote the  rate of the Poisson process for the birth branch with $c(\mb x^i) \ge 0$. For any test function $f\in C^2_b(\mathbb R^+, \mathbb R^2)$, we have
\begin{equation}
\begin{aligned}
     & \mathbb E\left[\left<f,  \mu_+\right>\right]
    \\=& \frac{1}{N} \sum_{i\in \mathcal I ,c(\bX_i)>0} \sum_{j=0}^{\infty} 
    \left[(j+1) \frac{e^{-\lambda_i} \lambda_i^j}{j!}\right] f(\bX_i)
     + \frac{1}{N} \sum_{i\in \mathcal I , c(\bX_i)<0}
    \exp( c(\bX_i)\tau)f(\bX_i) \\
    % \\=& \frac{1}{N} \sum_{i\in \mathcal I ,c(\bX_i)>0} 
    % \left[\exp(- c(\bX_i)\tau)f(\bX_i) + (1-\exp( -c(\bX_i)\tau))2f(\bX_i)\right]\\
    % & + \frac{1}{N} \sum_{i\in \mathcal I , c(\bX_i)<0}
    % \exp( c(\bX_i)\tau)f(\bX_i) \\
    =& \frac{1}{N}\sum_{i\in \mathcal I ,c(\bX_i)>0} 
    \left(
    \exp(c(\bX_i)\tau)f(\bX_i)
    \right) 
     + \frac{1}{N}\sum_{i\in \mathcal I , c(\bX_i)<0}\left(
    \exp( c(\bX_i)\tau)f(\bX_i) \right)\\
    = & \frac{1}{N}\sum_{i\in \mathcal I}\exp( c(\bX_i)\tau)f(\bX_i) 
    = \left<f,\exp( c\tau)\mu\right>,
\end{aligned}
\end{equation}
where the second equality follows from the fact that the expectation of a Poisson process equal to the rate.
% where the first equality comes from the expectation number of population after birth is $2$ and death if $0$. 
\end{proof}
By choosing $c(\mb x) = \tilde{r}(\rho^{\ast}(\mb x))$, this stochastic branching process realizes the operator $\exp(\mathcal{B}\tau)$ at the particle level. We summarize the process in Algorithm \ref{alg:birthdeath}.



\begin{algorithm}
\caption{Birth--Death Step}\label{alg:birthdeath}
\begin{algorithmic}
\Require Propagated particles $S=\{\mathbf X_{i}\}_{i\in\mathcal I}$, time step $\tau>0$, rate function $r:\mathbb R^2\to\mathbb R$
\Ensure Post-branching set $S_{+}=\{\mathbf X_{i}\}_{i\in\mathcal I_{+}}$
\Procedure{Birth\_Death}{$S,\,\tau,\,r(\cdot)$}
    \State $S_{n+1}\gets [\ ]$
    \For{$i\in\mathcal I$}
        \State $r_i \gets r(\mathbf X_{i})$
        \If{$r_i>0$} \Comment{birth}
            \State Append $\mathbf X_{i}$ to $S_{+}$ \Comment{keep parent}
            \State $p_{\mathrm{birth}} \gets \exp(r_i\,\tau) - 1$
            %\State With probability $p_{\mathrm{birth}}$, append a duplicate of $\mathbf X^{i}$ to $S_{+}$
            \State Draw $K_i \sim \mathrm{Poisson}\!\big(p_{\mathrm{birth}}\big)$, append $K_i$ duplicates of $\mathbf X_{i}$ to $S_{+}$
        \ElsIf{$r_i<0$} \Comment{death}
            \State $p_{\mathrm{survive}}\gets \exp(r_i\,\tau)$ 
            \State With probability $p_{\mathrm{survive}}$, append $\mathbf X_{i}$ to $S_{+}$
            \State \quad (otherwise the particle is removed)
        \Else \Comment{$r_i=0$}
            \State Append $\mathbf X_{i}$ to $S_{+}$
        \EndIf
    \EndFor
    \State \Return $S_{+}$
\EndProcedure
\end{algorithmic}
\end{algorithm}




\subsection{Nonlinear Interpolation}
\label{sec:interpolation}
To evaluate the functions $\tilde a(\rho)$ , and $\tilde r(\rho)$ , that govern the drift and reaction dynamics in our particle-based scheme, we need to reconstruct the continuous density 
$\rho$ from its empirical measure. Specifically, given an empirical measure
\[
\mu^N_n= \frac{1}{N}\sum_{i\in \mathcal I_n}\delta_{\bX_{i,n}^N},
\] 
we seek a smooth approximation of $\rho^N_n$ that allows pointwise evaluation and differentiation as required by the operator $\mathcal{A}$ and $\mathcal{B}$.

In this work, we choose to interpolate the empirical distribution in the Fourier space. Specifically, we project the density onto a finite set of Fourier basis functions, allowing efficient representation and differentiation. 

The Fourier coefficients are evaluated from the particle locations $\bX_{N,n}^i$ under periodic boundary conditions. Without loss of generality, we take the periodic domain to be $\Omega=[-\pi,\pi]^d$ , from which the reconstructed density $\rho_n(t,x)$. To be more specific, we define a Fourier interpolation operator $\mathcal P_K$
that maps the empirical particle distribution into a truncated Fourier representation 
\[
\rho^N_n(\bx) = \mathcal{P}_K(\mu^N_n)(\bx) 
          = \sum_{\vert k_1\vert ,\vert k_2\vert \leq K} \hat{\rho}_n(k_1,k_2,\cdots,k_d)\,\psi_{k_1}(\bx_1)\psi_{k_2}(\bx_2)\cdots\psi_{k_d}(\bx_d),
\]
where the basis functions are chosen as
\[
\psi_k(x) =
\begin{cases}
\frac{1}{\sqrt{2\pi}}, & k=0, \\
\frac{1}{\sqrt{2\pi}} \cos(k  x), & k \in \mathbb{Z}_+, \\
\frac{1}{\sqrt{2\pi}}\sin(k  x), & k \in \mathbb{Z}_- .
\end{cases}
\]

The Fourier coefficients are computed directly from the empirical measure:
\begin{equation}
\begin{aligned}
    \hat{\rho}_n(k_1,k_2,\cdots,k_d) &= \int \psi_{k_1}(y_1)\psi_{k_2}(y_2)\cdots \psi_{k_d}(y_d)\,\mu^N_n(\intd\by)
% &= \frac{1}{N}\sum_{i\in \mathcal I^N_n} \psi_{k_1}(\bX_{N,n,1}^{i})\psi_{k_2}(\bX_{N,n,2}^{i})\cdots \psi_{k_d}(\bX_{N,n,d}^{i}).
\end{aligned}
\label{eq:Fourier_mode}
\end{equation}
The process is shown in Algorithm \ref{alg:interpolate}. 
This choice is motivated by the simplicity and spectral accuracy of Fourier methods for smooth periodic functions. However, we note that other functional spaces, such as polynomial bases, wavelets, or kernel-based representations, could offer better adaptability in non-periodic or highly localized settings. We leave the exploration of such alternatives to future work.



\begin{algorithm}
\caption{Interpolation Step (fixed-$N$ normalization)}\label{alg:interpolate}
\begin{algorithmic}
\Require Particle set: $S=\{ \mathbf{X}^i \}_{i\in \mathcal{I}}$; fixed initial population $N$; truncation index set $\mathcal{M}_K$; domain descriptor $\Omega$; basis $\{\psi_m\}_{m\in\mathcal{M}_K}$ with (optionally) known gradients $\{\nabla\psi_m\}_{m\in\mathcal{M}_K}$
\Ensure Truncated field $\tilde{\rho}_n(\mathbf{x})$ (and, if needed, $\nabla\tilde{\rho}_n(\mathbf{x})$)
\Procedure{Interp}{$S_n,\,N,\,\mathcal{M}_K,\,\Omega,\,\{\psi_m\}$}
    \State \textbf{(Coefficient projection with fixed-$N$ normalization)}
    \State For each $m\in\mathcal{M}_K$, set
    \[
      a_m \;\gets\; \frac{1}{N}\sum_{i\in\mathcal{I}} \psi_m\!\big(\mathbf{X}^i\big).
    \]
    \State \textbf{(Series assembly)}
    \[
      \tilde{\rho}_n(\mathbf{x}) \;\gets\; \sum_{m\in\mathcal{M}_K} a_m\,\psi_m(\mathbf{x}).
    \]
    \State \textbf{(Optional gradients for drift/evaluation)}
    \If{gradients are required and $\{\nabla\psi_m\}$ are available}
      \[
        \nabla\tilde{\rho}_n(\mathbf{x}) \;\gets\; \sum_{m\in\mathcal{M}_K} a_m\,\nabla\psi_m(\mathbf{x}).
      \]
    \EndIf
    \State \Return $\tilde{\rho}_n$ (and $\nabla\tilde{\rho}_n$ if computed)
\EndProcedure
\end{algorithmic}
\end{algorithm}


\section{Error Estimations}
In this section, we analyze the error between the numerical and exact solutions for Equation (\ref{equ:reaction-diffusion}). 
We decompose the total error into three parts:
\begin{equation}\label{eq:err-decomp}
u(t_n)-\mu_n^N
=\underbrace{\bigl(u(t_n)-u_K(t_n)\bigr)}_{\text{coefficient projection}}
+\underbrace{\bigl(u_K(t_n)-\mu^N(t_n)\bigr)}_{\text{particle discretization}}
+\underbrace{\bigl(\mu^N(t_n)-\mu_n^N\bigr)}_{\text{time discretization}}.
\end{equation}
The first term is the consistency error between the exact PDE and the
\emph{coefficient-projected} PDE
\begin{equation}\label{equ:projected}
    \partial_t u_K  = \nabla \cdot \left( a (\mathcal P_K u_K) u_K\right) + D \Delta u_K  + r\left(\mathcal P_K u_K\right)  u_K,
\end{equation}
where $\mathcal P_K$ is the projection onto the space of trigonometric
polynomials of degree $\le K$,
\[
  \mathcal T_K
  := \operatorname{span}_{\mathbb C}
     \bigl\{\, e^{i\mathbf k\cdot\mathbf x} : \mathbf k\in\mathbb Z^d,\ |\mathbf k|_\infty\le K \,\bigr\},
  \qquad
  \mathcal P_K u = \sum_{|\mathbf k|_\infty\le K}\hat u_{\mathbf k}\,e^{i\mathbf k\cdot\mathbf x},
\]
with Fourier coefficients
$\hat u_{\mathbf k}=(2\pi)^{-d}\int_{\mathbb T^d}u(\mathbf x)\,e^{-i\mathbf k\cdot\mathbf x}\,\mathrm d\mathbf x$.
Notice that $u_K$ only replaces the coefficients $a(u),r(u)$ in
\eqref{equ:reaction-diffusion} by $a(\mathcal P_K u_K),r(\mathcal P_K u_K)$, so it
is \emph{not} a spectral Galerkin projection of the equation
($u_K\notin\mathcal T_K$ in general); this term is estimated in
Proposition~\ref{prop:u_u_K}.
The second term is the continuous-time particle approximation error, in which the
solution $u_K$ of the coefficient-projected PDE \eqref{equ:projected} is
approximated by the empirical measure $\mu^N$ of the $N\mhyphen$particle system
(Proposition~\ref{prop:u_K_mu}).
The third term is the time-discretization error between the continuous-time
branching particle process $\mu^N$ and the fully discrete algorithm $\mu^N_n$,
controlled for fixed $K$ in Proposition~\ref{prop:time_discrete_particle}.
For notational simplicity, we assume throughout the analysis that the initial
mass is normalized, $Z:=\int_\Omega u(0,\bx)\,\intd\bx=1$. The rescaled
coefficients then coincide with the originals (e.g.\ $a(Z\rho)=a(\rho)$), so we
write $a$ and $r$ throughout this section.

First, we estimate the error between the coefficient-projected PDE \eqref{equ:projected} and the original PDE \eqref{equ:reaction-diffusion}.
\begin{proposition}[Consistency of the coefficient-projected PDE]
\label{prop:u_u_K}
Let $\mathbb{T}^d = (\mathbb{R}/2\pi\mathbb{Z})^d$, and assume that the coefficients $a$ and $r$ are bounded and Lipschitz continuous with a uniform constant $C_{b,lip}$, i.e. for any $x, y\in \mathbb R$, we have
    \[
    |a(x)|,|r(x)| < C_{b,lip}, \qquad |a(x)-a(y)|,|r(x)-r(y)| < C_{b,lip} |x-y|.
    \]
    Suppose that for a finite time horizon $[0,T)$: $u$ is a sufficiently smooth solution of the nonlinear PDE~\eqref{equ:reaction-diffusion} on $\Omega$ with periodic boundary conditions, and $u_K$ is the solution of the coefficient-regularized PDE~\eqref{equ:projected}.
If $u\in L^\infty\bigl(0,T;\mathbb H^m(\Omega)\bigr)$ for some $m>d/2+1$, so that $\mathbb H^m\hookrightarrow W^{1,\infty}\hookrightarrow L^\infty$, then the error $e_K=u- u_K$ satisfies
\[
\|e_K(t)\|_{L^2}^2 \;\le\; C_T\Big(\|e_K(0)\|_{L^2}^2 + K^{-2m}\|u\|_{L^2(0,t;\mathbb H^m)}^2\Big),
\]
for some constant $C_T>0$ depending on $T$, $D$,
but independent of $K$.
\end{proposition}


\begin{proof}
    Let $\eta(t) = \frac{1}{2}\|e_K(t,\cdot)\|^2_{L^2(\Omega)}  = \frac{1}{2}\|u(t,\cdot)- u_K(t,\cdot)\|^2_{L^2(\Omega)} $, then we have 
    \begin{equation}
        \begin{aligned}
    \frac{\partial}{\partial t} \eta(t)
        = & -\int_\Omega \nabla (u-u_K)\cdot (a(u)u - a(\mathcal P_K u_K )u_K )\intd \bx \\
        & + D\int_\Omega (u-u_K)  \Delta (u -  u_K  )\intd \bx \\
        & + \int_\Omega (u-u_K)  (r(u)u - r(\mathcal P_K  u_K ) u_K) \intd \bx  \\
        = &  E_1 + E_2 + E_3
        \end{aligned}
    \end{equation}
Notice that for $E_1$, we have
\[
\begin{aligned}
    E_1 =& -\int_\Omega \nabla(u-u_K) \cdot \left(\left(a(u) - a(\mathcal P_K u_K )\right)u\right) \intd \bx\\
    &- \int_\Omega \nabla(u-u_K) \cdot \left(a(\mathcal P_K u_K )(u  - u_K )\right)\intd \bx = E_{11}+ E_{12}.\\
\end{aligned}
\]
For $E_{11}$, we have 
\[
\begin{aligned}
E_{11} \leq & \|\nabla e_K \|_{L^2} \|\left(a(u) - a(\mathcal P_K u_K )\right) u \|_{L^2}\\
 \leq& \|\nabla e_K \|_{L^2} C_{b,lip}\left\|u\|_{L^\infty}\right\|u - \mathcal P_K u_K \|_{L^2}\\
\leq&  \frac{\epsilon C_{b,lip} \left\|u\right\|_{L^\infty} }{2} \|\nabla e_K \|_{L^2}^2 
+\frac{C_{b,lip}\|u\|_{L^\infty}}{2\epsilon} \|u - \mathcal P_K u_K\|_{L^2}^2.
\end{aligned}
\]
Since $m>d/2+1$, the embedding $\mathbb H^m\hookrightarrow L^\infty$ gives
$\sup_{t\le T}\|u(t)\|_{L^\infty}\le C\,\|u\|_{L^\infty(0,T;\mathbb H^m)}=:C_u'<\infty$. By the spectral accuracy of $\mathcal P_K$,
\[
\begin{aligned}
\|u - \mathcal P_K u_K\|_{L^2} \leq& \|u - \mathcal P_K u\|_{L^2} + \|\mathcal P_K u - \mathcal P_K u_K\|_{L^2}\\
=& \|(I - \mathcal P_K) u\|_{L^2} + \|\mathcal  P_K (u - u_K)\|_{L^2} \\
\leq &\|(I - \mathcal P_K) u\|_{L^2} + \|e_K\|_{L^2}\\
\leq & C K^{-m} \|u\|_{H^m} + \|e_K\|_{L^2}.
\end{aligned}
\]
For $E_{12}$, we have 
\[
\begin{aligned}
    E_{12} \leq& \left| \int_\Omega \nabla(u-u_K) \cdot \left(a(\mathcal P_K u_K )(u  - u_K )\right)\intd \bx \right|\\
    \leq& \|\nabla e_K\|_{L^2}\left\|a(\mathcal P_K u_K )(u  - u_K ) \right\|_{L^2} \\
    \leq& C_{b,lip} \|\nabla e_K\|_{L^2}\left\|e_K \right\|_{L^2} \\
    \leq& \frac{\epsilon C_{b,lip}}{2}\|\nabla e_K\|_{L^2}^2 +\frac{C_{b,lip}}{2\epsilon}\|e_K\|_{L^2}^2.
\end{aligned}
\]
Add two term together we have 
\[
E_1 \leq \frac{\epsilon C_1}{2}\|\nabla e_K\|_{L^2}^2 +\frac{C_1}{2\epsilon}\| e_K\|_{L^2}^2 + C_1 K^{-2m}\|u\|_{\mathbb H^m}^2,
\]
where $C_1$  depends on $ \|a\|_{L^\infty},\left\|a'\right\|_{L^\infty},\|u\|_{L^\infty}$.
For $E_2$, we have 
\[
E_2 \leq - D\int_\Omega (\nabla e_K )^2\intd \bx \leq -D\|\nabla e_K\|^2_{L^2}.
\]
For $E_3$, we have 
\[
\begin{aligned}
    E_3 = &   \int_\Omega (u-u_K)  (r(u) - r(\mathcal P_K  u_K )) u \intd \bx 
     +  \int_\Omega (u-u_K)  r(\mathcal P_K  u_K ) (u-u_K) \intd \bx \\
    \leq & C_{b,lip}\,C_u'\,\|e_K\|_{L^2}\,\|u-\mathcal P_K u_K\|_{L^2 }
     +C_{b,lip} \|e_K\|_{L^2}^2\\
    \leq & \tfrac{1}{2}C_{b,lip}C_u' \|e_K\|_{L^2} ^2 + \tfrac{1}{2}C_{b,lip}C_u'\bigl( C K^{-2m} \|u\|_{H^m}^2 + \|e_K\|_{L^2}^2\bigr)
     +C_{b,lip} \|e_K\|_{L^2}^2\\
    \leq & C_3 \|e_K\|_{L^2}^2+ C_3 K^{-2m} \|u\|_{H^m}^2.
\end{aligned}
\]
where $C_3$ depends on $\|r\|_{L^\infty},\|r'\|_{L^\infty}$ and $C_u'=\sup_{t\le T}\|u(t)\|_{L^\infty}$.

By choose $\epsilon$ small enough compared with $D$, we can have
\[
\frac{\partial }{\partial t }\eta \leq(C_1+C_2)\eta(t) + (C_1+C_3) K^{-2m} \|u\|_{H^m}^2.
\]
Applying the Gr\"onwall inequility yields
$$\eta(t) \le \eta(0)e^{(C_1+C_2)t} + \int_0^t e^{(C_1+C_2)(t-s)} (C_1+C_3) K^{-2m} \|u(s)\|_{H^m}^2 \intd s.$$
Therefore, we have
$$\|e_K(t)\|_{L^2}^2 \le C_T \Big( \|e_K(0)\|_{L^2}^2 + K^{-2m} \|u\|_{L^2(0,t;H^m)}^2 \Big).$$
\end{proof}

Next, we analyze the particle discretization error in continuous time.
Let $S(N,t)=\{\bX^N_{i,K}(t)\}_{i\in\mathcal I(N,t)}$ denote the set of particles
alive at time $t$, where $\mathcal I(N,t)$ is the corresponding index set, and
define the associated empirical measure
\[
  \mu^N(t,\intd\bx)
  = \frac{1}{N}\sum_{i\in\mathcal I(N,t)}\delta_{\bX^N_{i,K}(t)}(\intd\bx),
\]
the random finite measure obtained from the particle configuration at time $t$.
In weak form, for every test function $f\in C^2(\mathbb T^d)$,
 \begin{equation}\label{eq:weak-form}
     \begin{aligned}
\langle f, \mu^N(t)\rangle
&= \langle f, \mu^N(0)\rangle
 - \int_0^t
   \!\!\langle\nabla f,\,a(\mathcal P_K\mu^N)\mu^N\rangle\,\intd\tau
 + \int_0^t
   \!\!\langle\Delta f,\,D\,\mu^N\rangle\,\intd\tau \\
&\quad + \int_0^t \langle f, r(\mathcal P_K \mu^N)\mu^N\rangle\,\intd\tau 
 + \frac{1}{N}\int_0^t\sum_{i\in\mathcal I(N,\tau)}
   \sqrt{2D}\,\nabla f(\bX_{i,K}^N(\tau))\!\cdot\! \intd \bW^N_i(\tau)\\
 &\quad + \frac{1}{N}\int_0^t\!\sum_{i\in\mathcal I(N,\tau)}\!
   f(\bX_{i,K}^N(\tau))\, \intd \widetilde{\mathcal R}^{+,N}_i(\tau)
 - \frac{1}{N}\int_0^t\!\sum_{i\in\mathcal I(N,\tau)}\!
   f(\bX_{i,K}^N(\tau))\, \intd \widetilde{\mathcal R}^{-,N}_i(\tau),
\end{aligned}
 \end{equation}
where the last two terms are \emph{compensated} Poisson martingales. Let
$\{Q^{+,N}_i,Q^{-,N}_i\}_{i\ge1}$ be standard (unit-rate) Poisson processes,
mutually independent and independent of $\{\bW^N_i\}$, so that
\[
\mathbb E[Q^{\pm,N}_i(t)]=t,\qquad
Q^{\pm,N}_i(t)-t \text{ is a martingale},\qquad
\mathbb E[(Q^{\pm,N}_i(t)-t)^2]=t
\]
With $r^\pm:=\max(\pm r,0)$ and the birth/death rates
$c^\pm(\bx):=r^\pm\!\bigl((\mathcal P_K\mu^N)(\bx)\bigr)$, the counts are the
time-changed processes $\mathcal R^{\pm,N}_i=Q^{\pm,N}_i\!\circ\Lambda^{\pm,N}_i$,
\[
\mathcal R^{\pm,N}_i(t)=Q^{\pm,N}_i\!\bigl(\Lambda^{\pm,N}_i(t)\bigr),
\qquad
\Lambda^{\pm,N}_i(t):=\int_0^t \mathbf 1_{\mathcal I(N,\tau)}(i)\,
   c^\pm\bigl(\bX^N_{i,K}(\tau)\bigr)\,\intd\tau .
\]
Subtracting the compensator $\Lambda^{\pm,N}_i$ (the conditional mean) gives the
martingale $\widetilde{\mathcal R}^{\pm,N}_i:=\mathcal R^{\pm,N}_i-\Lambda^{\pm,N}_i$,
and the time change carries $t\mapsto\Lambda^{\pm,N}_i(t)$ into the quadratic
variation:
\begin{equation}\label{eq:R-qv}
\mathbb E\bigl[(\widetilde{\mathcal R}^{\pm,N}_i(t))^2\bigr]
=\Lambda^{\pm,N}_i(t).
\end{equation}
Independence across $\pm,i$ and from $\{\bW^N_i\}$ makes all these martingales
mutually orthogonal, so cross terms vanish in the $L^2$ estimates below.

The following proposition bounds the second error contribution in \eqref{eq:err-decomp}, the continuous-time particle approximation error between the projected PDE solution $u_K$ and the empirical measure $\mu_N$.
\begin{proposition}[Continuous-time particle approximation]
\label{prop:u_K_mu}
Let $\mathbb T^d=(\mathbb R/2\pi\mathbb Z)^d$, let $s>d/2+1$, and fix the spectral
cutoff $K$. Assume the coefficients $a,r$ are bounded and Lipschitz continuous
with a uniform constant $C_{b,\mathrm{lip}}$,
\[
  |a(x)|,|r(x)|\le C_{b,\mathrm{lip}},\qquad
  |a(x)-a(y)|,|r(x)-r(y)|\le C_{b,\mathrm{lip}}\,|x-y|
  \quad(x,y\in\mathbb R),
\]
and suppose that on a finite horizon $[0,T)$:
\begin{enumerate}
\item[(i)] the active set is uniformly bounded, $|\mathcal I(N,t)|\le C_N N$;
\item[(ii)] the projected PDE \eqref{equ:projected} admits a unique solution $u_K$;
\item[(iii)] the band-limited coefficients and the solution are uniformly
$W^{\lceil s\rceil,\infty}$-regular,
\begin{equation}\label{eq:prop2-coeff}
\sup_{t\le T}\Bigl(\bigl\|a(\mathcal P_K\mu^N(t))\bigr\|_{W^{\lceil s\rceil,\infty}}
+\bigl\|r(\mathcal P_K\mu^N(t))\bigr\|_{W^{\lceil s\rceil,\infty}}
+\|u_K(t)\|_{W^{\lceil s\rceil,\infty}}\Bigr)\le C_{T,K}.
\end{equation}
\end{enumerate}
Then the $H^{-s}$ error $E_K(t):=\tfrac12\|u_K(t)-\mu^N(t)\|_{H^{-s}}^2$ satisfies
\[
\frac{\intd}{\intd t}\,\mathbb E E_K(t)\le C\,\mathbb E E_K(t)+\frac{C}{N},
\]
with a constant $C=C_{T,K}$ depending on $C_N,C_{b,\mathrm{lip}},D,s,d$ and on $K$
through \eqref{eq:prop2-coeff}, but \emph{independent of $N$}. Consequently, by
Gr\"onwall's inequality,
\[
\mathbb E E_K(t)\le e^{C_{T,K}t}\,\mathbb E E_K(0)+\frac{C'_{T,K}}{N},
\qquad C'_{T,K}:=e^{C_{T,K}T}-1.
\]
\end{proposition}
%\end{theorem}
 \begin{proof}
Write $\widehat e_K:=u_K-\mu^N$ and recall $R:=(I-\Delta)^{-s}$,
$\phi_K:=R\,\widehat e_K$, so that
\[
E_K=\tfrac12\|\widehat e_K\|_{H^{-s}}^2
=\tfrac12\langle\widehat e_K,R\,\widehat e_K\rangle
=\tfrac12\langle\widehat e_K,\phi_K\rangle .
\]
The time derivative of the expectation $\mathbb{E}E_K(t)$ is given by Ito formula
\begin{equation}
 \begin{aligned}
 \frac{\intd}{\intd t} \mathbb E  E_K(t)  = &   \quad \mathbb E \left[\left\langle \nabla  \phi_K , a(\mathcal P_K \mu^N) \mu^N-  a(\mathcal P_K u_K) u_K   \right\rangle \right] \\
  & +  \mathbb E \left[\left\langle \Delta  \phi_K , D (u_K - \mu^N) \right\rangle \right] \\
  & + \mathbb E \left[\left\langle   \phi_K , r(\mathcal P_K u_K)u_K - r(\mathcal P_K \mu^N)\mu^N \right\rangle \right]  \\
  & +  \mathbb E \left[\frac{D}{N^2} \sum_{i\in \mathcal I (N,t)}  \bigl\|\nabla (I-\Delta)^{-s/2}\delta_{\bX_N^i(t)}\bigr\|_{L^2}^2 \right] \\
  & + \mathbb E \left[\frac{1}{2N^2} \sum_{i\in \mathcal I (N,t)} \bigl\|\delta_{\bX_N^i(t)}\bigr\|_{H^{-s}}^2 \left|r(\mathcal{P}_K \mu^N(\bX_N^i(t)))\right|\right]\\
  & = \mathbb E \left[I_1\right]+\mathbb E \left[I_2\right]+\mathbb E \left[I_3\right]+\mathbb E \left[I_4\right]+
  \mathbb E \left[I_5\right].
 \end{aligned}
\end{equation}
Before bounding $I_1$--$I_3$ we record two facts, both at fixed $K$ and both
with constants absorbed into $C_{T,K}$. Throughout the proof, $C_{T,K}$ denotes a positive constant depending on $C_N,C_{b,\mathrm{lip}},D,s,d,T$ and on $K$ through \eqref{eq:prop2-coeff}, but \emph{not} on $N$, whose value may change from line to line.


\emph{(M) Multiplier bound.} A $W^{\lceil s\rceil,\infty}$ function multiplies
$H^\sigma$ for every $|\sigma|\le s$,
$\|fg\|_{H^\sigma}\le C_{s,d}\|f\|_{W^{\lceil s\rceil,\infty}}\|g\|_{H^\sigma}$
(by duality for $\sigma<0$). With \eqref{eq:prop2-coeff} this gives, for
$|\sigma|\le s$,
\[
\|a(\mathcal P_K\mu^N)g\|_{H^\sigma}
+\|r(\mathcal P_K\mu^N)g\|_{H^\sigma}
+\|u_K g\|_{H^\sigma}\le C_{T,K}\|g\|_{H^\sigma}.
\]
An $L^\infty$ bound alone is insufficient once $s>1$, since $L^\infty$ is not an
$H^{s-1}$ multiplier --- this is what forces the $W^{\lceil s\rceil,\infty}$
regularity in \eqref{eq:prop2-coeff}.

\emph{(I) Inverse estimate.} For $h\in\mathcal T_K$ and any $\alpha\le\beta$,
$\|h\|_{H^\beta}\le C_{T,K}\|h\|_{H^\alpha}$. Applied to
$h=\mathcal P_K\hat e_K\in\mathcal T_K$ with $(\alpha,\beta)=(-s,0)$, together with
the $H^{-s}$-contraction $\|\mathcal P_K\hat e_K\|_{H^{-s}}\le\|\hat e_K\|_{H^{-s}}$,
\[
\|\mathcal P_K\hat e_K\|_{L^2}
\le C_{T,K}\|\mathcal P_K\hat e_K\|_{H^{-s}}
\le C_{T,K}\|\hat e_K\|_{H^{-s}}.
\]
\textbf{Term $I_1$ (Advection)}: We split this term: 
\begin{equation}
    \begin{aligned}
        I_1 
        = &  
         \left\langle \nabla  \phi_K,  a(\mathcal P_K \mu^N ) ( \mu^N - u_K )\right\rangle 
        +
        \left\langle \nabla  \phi_K,  \left( a(\mathcal P_K \mu^N) - a(\mathcal P_K u_K) \right) u_K \right\rangle \\
        =& I_{11} + I_{12}.
    \end{aligned}
\end{equation}
For $I_{11}$, using the duality $\langle H^{s-1}, H^{-s+1}\rangle$, the multiplier
bound (M) with $\sigma=s-1$, and Young's inequality:
\[
\begin{aligned}
    I_{11}
        \leq & 
        \left|\left\langle a(\mathcal P_K \mu^N)  \nabla  \phi_K,   (u_K - \mu^N )\right\rangle\right|\\
        \leq& \left\|a(\mathcal P_K \mu^N)  \nabla  \phi_K\right\|_{H^{s-1}} \left\|u_K - \mu^N \right\|_{H^{-s+1}} \\
        \leq& C_{T,K} \left\|\nabla  \phi_K\right\|_{H^{s-1}} \left\|u_K - \mu^N \right\|_{H^{-s+1}} \\
        \leq& \frac{C_{T,K}}{2\epsilon}\left\|\nabla  \phi_K\right\|_{H^{s-1}}^2 + \frac{C_{T,K}\epsilon}{2}\|\hat{e}_K\|^2_{H^{-s+1}}\\
         \leq& \frac{C_{T,K}}{2\epsilon}\left\|  \phi_K\right\|_{H^{s}}^2 + \frac{C_{T,K}\epsilon}{2}\|\hat{e}_K\|^2_{H^{-s+1}}
        = \frac{C_{T,K}}{\epsilon} E_K(t)+ \frac{C_{T,K}\epsilon}{2}\|\hat{e}_K\|^2_{H^{-s+1}}.
\end{aligned}
\]
For $I_{12}$, using the multiplier bound (M) for $u_K$, the Lipschitz continuity of
$a$, and the inverse estimate (I):
\[
\begin{aligned}
     I_{12} \leq & \left\|u_K \nabla \phi_K\right\|_{H^{s-1}} \bigl\| a(\mathcal  P_K \mu^N) - a(\mathcal P_K u_K)\bigr\|_{H^{-s+1}}\\
     \leq&  C_{T,K}\left\|\nabla  \phi_K\right\|_{H^{s-1}} \bigl\| a(\mathcal  P_K \mu^N) - a(\mathcal P_K u_K)\bigr\|_{H^{-s+1}}\\
     \leq&  C_{T,K}\left\|\nabla  \phi_K\right\|_{H^{s-1}} \,C_{b,lip}\left\|\mathcal P_K\hat{e}_K\right\|_{L^2} \\
     \leq & C_{T,K}\left\| \phi_K\right\|_{H^{s}} \left\|\hat{e}_K \right\|_{H^{-s+1}}
     \leq \frac{C_{T,K}}{2\epsilon}E_K(t)
        + \frac{C_{T,K}\epsilon}{2}\|\hat{e}_K\|^2_{H^{-s+1}}.
\end{aligned}
\]
In total, we have
\[
I_1\leq \frac{C_{T,K}}{\epsilon}E_K(t)
        + \frac{C_{T,K}\epsilon}{2}\|\hat{e}_K\|^2_{H^{-s+1}}.
\]
\textbf{Term $I_2$ (Diffusion): }
\[
\begin{aligned}
I_{2} = & D\left\langle (I-(I-\Delta))  \phi ,  \hat{e}_K \right\rangle\\
   =& D\left\langle  \phi_K ,  \hat{e}_K \right\rangle -  D\left\langle (I-\Delta)  \phi_K ,  \hat{e}_K \right\rangle\\
   = & 2D E_K(t) - D\left\langle (I-\Delta)^{-s+1}  \hat{e}_K ,  \hat{e}_K \right\rangle\\
   = & 2D E_K(t) - D\|\hat{e}_K\|_{H^{-s+1}}^2.
\end{aligned}
\]
\textbf{Term $I_3$ (Reaction):} This follows the same logic as $I_1$, using the
multiplier bound (M) for $u_K$ and $r(\mathcal P_K\mu^N)$ on $H^{\pm s}$, the Lipschitz
continuity of $r$, and the inverse estimate (I).
\[
\begin{aligned}
I_3 
=& \left\langle   \phi_K , \left(r\left(\mathcal P_K u_K\right)-r\left(\mathcal P_K \mu^N\right)\right)u_K\right\rangle 
+ \left\langle   \phi_K , r(\mathcal P_K \mu^N)\left(u_K-\mu^N\right) \right\rangle \\
\leq & \|\phi_K\|_{H^s}\bigl\|\bigl(r(\mathcal P_K u_K)-r(\mathcal P_K\mu^N)\bigr)u_K\bigr\|_{H^{-s}}
+ \bigl\|r(\mathcal P_K\mu^N)\,\phi_K\bigr\|_{H^s}\|\hat e_K\|_{H^{-s}} \\
\leq & C_{T,K}\,\|\phi_K\|_{H^s}\,\bigl\|r(\mathcal P_K u_K)-r(\mathcal P_K\mu^N)\bigr\|_{H^{-s}}
+ C_{T,K}\,\|\phi_K\|_{H^s}\|\hat e_K\|_{H^{-s}} \\
\leq & C_{T,K}\,\|\phi_K\|_{H^s}\,C_{b,lip}\|\mathcal P_K\hat e_K\|_{L^2}
+ C_{T,K}\,\|\phi_K\|_{H^s}\|\hat e_K\|_{H^{-s}}
\leq  C_{T,K}\,\|\phi_K\|_{H^s}\|\hat e_K\|_{H^{-s}}
\le  C_{T,K}\, E_K(t).
\end{aligned}
\]
\textbf{Term $I_4$ (Diffusion QV):}
This term is an \emph{additive} sampling-noise forcing, not proportional to the
error. Since $(I-\Delta)^{-s/2}\delta_x$ has Fourier coefficients
$(1+|\mathbf k|^2)^{-s/2}e^{-i\mathbf k\cdot x}$,
\[
\bigl\|\nabla(I-\Delta)^{-s/2}\delta_x\bigr\|_{L^2}^2
=\sum_{\mathbf k}|\mathbf k|^2(1+|\mathbf k|^2)^{-s}=:C_s',
\]
which is finite \emph{precisely because} $s>d/2+1$ (the summand decays like
$|\mathbf k|^{2-2s}$) and is independent of $x$. Hence, by $|\mathcal I(N,t)|\le C_NN$,
\begin{equation}
    I_4 = \frac{D}{N^2}\sum_{i\in \mathcal I (N,t)}C_s'
    \le \frac{D C_N C_s'}{N}=:\frac{C_4}{N}.
\end{equation}
\textbf{Term $I_5$ (Branching QV):}
This term is likewise additive. Since
\[
\|\delta_x\|_{H^{-s}}^2=\sum_{\mathbf k}(1+|\mathbf k|^2)^{-s}=:C_s<\infty
\qquad(s>d/2),
\]
is independent of $x$ and $\|r\|_{L^\infty}\le C_{b,lip}$,
\[
    I_5 = \frac{1}{2N^2}\sum_{i\in \mathcal I (N,t)}
    \bigl\|\delta_{\bX_N^i}\bigr\|_{H^{-s}}^2\,\bigl|r(\mathcal P_K\mu^N(\bX_N^i))\bigr|
    \le \frac{C_N C_{b,lip} C_s}{2N}=:\frac{C_5}{N}.
\]
\textbf{Conclusion:}
Combining all the terms, gives us
\[
\frac{\intd}{\intd t}\mathbb{E}E_{K}(t) \le \left( \frac{C_{T,K}}{\epsilon} + 2D + C_{T,K} \right) \mathbb{E}E_K(t) + \frac{C_4+C_5}{N} + \left( \frac{C_{T,K} \epsilon}{2} - D \right) \mathbb{E}\|\hat{e}_{K}\|_{H^{-s+1}}^{2}.
\]
Since $D > 0$, choosing $\epsilon$ small (e.g.\ $\epsilon = D/C_{T,K}$) makes the last
term nonpositive. Collecting the $E_K$-coefficient $\tfrac{C_{T,K}}{\epsilon}+2D+C_{T,K}$ at
$\epsilon=D/C_{T,K}$ together with the $K$-independent forcing constants $C_4,C_5$,
\[
    \frac{\intd}{\intd t} \mathbb E  E_K(t) \leq C_{T,K}\, \mathbb E  E_K(t) + \frac{C_{T,K}}{N}. 
\]
By Gr\"onwall with constant forcing, using $t\le T$,
\[
\mathbb E E_K(t) \leq e^{C_{T,K}t}\,\mathbb E E_K(0) + \frac{e^{C_{T,K}t}-1}{N}
\le e^{C_{T,K}T}\,\mathbb E E_K(0) + \frac{C'_{T,K}}{N},
\qquad C'_{T,K}:=e^{C_{T,K}T}-1.
\]
In particular, if $\mu^N(0)$ is an unbiased sample of $u_K(0)$ so that
$\mathbb E E_K(0)\le C_s/N$, then
$\mathbb E\|u_K(t)-\mu^N(t)\|_{H^{-s}}^2\le C_{T,K}/N$ for all $t\le T$.
\end{proof}
% By the triangle inequality,
% \begin{equation*}
% \Vert u(t)-\mu^N(t)\Vert_{H^{-s}}
% \le \Vert u(t)-u_K(t)\Vert_{H^{-s}}
% +\Vert u_K(t)-\mu^N(t)\Vert_{H^{-s}}.
% \end{equation*}
% The first term is controlled by Proposition~\ref{prop:u_u_K}, which is
% stated in $L^2$; since $\Vert v\Vert_{H^{-s}}\le\Vert v\Vert_{L^2}$ for $s\ge0$
% on $\mathbb{T}^d$, the same bound controls the $H^{-s}$ norm, contributing the
% $K^{-2m}\Vert u\Vert_{L^2(0,t;H^m)}^2$ term. The particle term is bounded
% directly in $H^{-s}$ by Proposition~\ref{prop:u_K_mu}; choosing $u_K=u$ at
% $t=0$ and sampling $\mu^N$ unbiasedly from $u_0$, so that
% $\mathbb{E}E_K(0)\le C_s/N$. Squaring, adding the two estimates, and absorbing the resulting constants into $C_T$ yields
% \begin{equation}\label{eq:u-muN-cont}
% \mathbb{E}\Vert u(t)-\mu^N(t)\Vert_{H^{-s}}^2
% \le C_{T,K}\!\left(C_s/N+K^{-2m}\Vert u\Vert_{L^2(0,t;H^m)}^2\right).
% \end{equation}


% \begin{definition}[Weak formulation]
% Let $\Omega\subset\mathbb{R}^2$ be a periodic domain and $T>0$.  
% Assume $\hat a,\hat d,\hat r \in L^\infty(\Omega_T)$ with $\hat d(t,x)\ge d_0>0$.  
% Given initial data $\hat u_0 \in L^2(\Omega)$, we call 
% \[
% \hat u \in L^2(0,T;H^1(\Omega)) \cap C([0,T];L^2(\Omega))
% \]
% a weak solution of \eqref{equ:linearlized} if for all test functions $f \in H^1(\Omega)$ and for almost every $t\in (0,T)$,
% \begin{equation}\label{eq:weak-evans}
% \begin{aligned}
%     &\left\langle\partial_t \hat u(t), f\right\rangle+ \int_\Omega \big( \hat a(t,x)\,\hat u (x)\cdot \nabla f(x) 
% + \nabla (\hat d(t,x) \hat u (x))\cdot \nabla f(x) \big)\,dx \\
% =& \left\langle\hat r(t,\cdot)\,\hat u(t,\cdot), f\right\rangle,
% \end{aligned}
% \end{equation}
% where $\langle\cdot,\cdot\rangle$ denotes the $L^2(\Omega)$ inner product.
% The initial condition is understood in the sense
% \[
% (\hat u(0), f)_{L^2} = (\hat u_0, f)_{L^2}, \quad \forall f\in L^2(\Omega).
% \]
% \end{definition}


% \begin{lemma}[Well-posedness of the linearized PDE]\label{lemma:well_posedness}
% Assume that $\hat a(t,x), \hat d(t,x)$,  and  $ \hat r(t,x) \in L^\infty(\Omega_T)$ with 
% $\hat d(t,x) \ge d_0 > 0$ for some constant $d_0$. 
% Then the linear parabolic PDE \eqref{equ:linearlized}
% admits a unique weak solution 
% \[
% \hat u \in L^2\big(0,T; H^1(\Omega)\big) \cap C\big([0,T]; L^2(\Omega)\big),
% \]
% satisfying 
% for any initial data $\hat u(0,\cdot) \in L^2(\Omega)$.
% \end{lemma}
% The proof can be found in standard PDE textbooks such as ~\cite{evans2022partial}.
% The continuous limit of our algorithm is described by a stochastic differential equation
% \[
% \intd \bX^N_i(t) = \hat a(t,\bX^N_i(t))\intd t  + \hat d (t,\bX^N_i(t)) \intd \bW^i(t).
% \]
% Also, any particle $i\in \mathcal I(N,t)$ at position $\bX^N_i(t)$ and time $t$ will induce discontinuous changes to the population by giving birth a particle $i^*\in \mathcal I(N,t)$ at $\bX^N_i$ with intensity $\hat r(x)$ if $\hat r(x)>0$ and die with intensity $-\hat r(x)$ of $\hat r(x)<0$. It induces the following two point process for $\sigma\in \mathbb R^+$ 
% \[
% \begin{aligned}
%    \mathcal R^{i^*,birth}_N( \sigma) = Q^{i,birth}_N \left(\int_{0}^{ \sigma} \mathbf 1 _{S(N,\tau)}(i)\hat r(t,\bX^N_i(\tau))\intd \tau\right),\\
%     \mathcal R^{i^*,death}_N( \sigma) = Q^{i,death}_N \left(-\int_{0}^{\sigma} \mathbf 1 _{S(N,\tau)}(i)\hat r(t,\bX^N_i(\tau))\intd \tau\right), 
% \end{aligned}
% \]
% where $Q^{i,birth}_N$, $Q^{i,death}_N$ are independent standard Poisson processes. 
% \begin{proposition}
%     Assume that $\hat a(t,x), \hat d(t,x)$,  and  $ \hat r(t,x) \in L^\infty(\Omega_T)$, therfore we have for $f \in H^2(\Omega) \subseteq H^1(\Omega)$,
%     \[
%     \mathbb E \left[ \langle f,\mu^N(t,\cdot)-\hat u(t,\cdot) \rangle \right] = 0,
%     \] and 
%     \[
%     \mathbb E \left[ \langle f,\mu^N(t,\cdot)-\hat u(t,\cdot) \rangle^2 \right] = \frac{1}{N},
%     \]
%     where $\mu^N$ is the empirical measure valued stochastic process before.
% \end{proposition}
% \begin{proof}
%     Using It\^o's formula, we have 
% \begin{equation*}
%     \begin{aligned}
%         \langle f,\mu^N(t,\cdot)\rangle = &\left\langle f(\tau,\cdot), \mu^N(0,\cdot)\right\rangle - \int_0^t\left\langle \hat a(\tau,\cdot)\cdot \nabla f(\tau,\cdot) ,\mu^N(\tau,\cdot)\right\rangle\intd \tau \\
%         &+ \int_0^t\left\langle \hat d(\tau, \cdot) \Delta f(\tau,\cdot),\mu^N(\tau,\cdot)\right\rangle \intd \tau\\
%         &+ \int_0^t\left\langle  \frac{\partial}{\partial \tau} f(\tau,\cdot),\mu^N(\tau,\cdot)\right\rangle \intd\tau\\
%         &+ \frac{1}{N} \int_0^t \sum_{i\in \mathcal I (N,\tau)} \sqrt{2 \hat d(\tau,\bX_N^i(\tau))} \nabla f(\tau,\bX_N^i(\tau))\cdot \intd \bW_N^i(\tau)\\
%         &+  \frac{1}{N} \int_0^t \sum_{i\in \mathcal I (N,\tau), \hat r(\tau, \bX^N_i(\tau))>0} f(\tau,\bX_N^i(\tau))\mathcal R_N^{i^*,birth}(\intd \tau)\\
%         &-  \frac{1}{N} \int_0^t \sum_{i\in \mathcal I (N,\tau), r(\tau, \bX^N_i(\tau))<0} f(\tau,\bX_N^i(\tau)) \mathcal R_N^{i^*,death}(\intd \tau)
%     \end{aligned}
% \end{equation*}
% One defines 
% \begin{equation*}
%     \begin{aligned}
%         \mathcal M^1(t,f) =& \frac{1}{N} \int_0^t \sum_{i\in \mathcal I(N,\tau)} \sqrt{2 \hat d(\tau,\bX_N^i(\tau))} \nabla f(\bX_N^i(\tau))\cdot \intd \bW_N^i(\tau)\\
%         \mathcal M^2(t,f) =& \frac{1}{N} \int_0^t \sum_{i\in \mathcal I(N,\tau), r(\bX^N_i(\tau))>0} f(\bX_N^i(\tau))\left(\mathcal R_N^{i^*,birth}(\intd \tau) -r(\bX_N^i(\tau))\intd \tau\right)\\
%         \mathcal M^3(t,f) =& \frac{1}{N} \int_0^t \sum_{i\in \mathcal I(N,\tau), r(\bX^N_i(\tau))<0} f(\bX_N^i(\tau)) \left(\mathcal R_N^{i^*,death}(\intd \tau)-r(\bX_N^i(\tau))\intd \tau\right).
%     \end{aligned}
% \end{equation*}
% We notice that 
% \begin{equation}
% \begin{aligned}
%     \mathbb E\left[ \left(\mathcal M^1(t,f) \right)\right] &= 0
% \end{aligned}
% \end{equation}
% and
% \begin{equation}
% \begin{aligned}
%     \mathbb E\left[ \left(\mathcal M^1(t,f) \right)^2\right] &= \frac{1}{N^2} \int_0^t \sum_{i\in \mathcal I(N,\tau)}  \left|\hat d(t,\bX^N_i(\tau))\nabla f(\bX_N^i(\tau))\right|^2 \intd \tau \\
%     & \leq \frac{C_{f',\sigma}t}{N},
% \end{aligned}
% \end{equation}
% where $C_{f',\sigma}$ is a constant only dependent on the gradient of $f$ and $\sigma$. Similarily, we have
% \begin{equation}
% \begin{aligned}
%     \mathbb E\left[ \mathcal M^2(t,f) \right] &= 0
% \end{aligned}
% \end{equation}
% and 
% \begin{equation*}
%     \begin{aligned}
%        \mathbb E\left[ \left(\mathcal M^2(t,f) \right)^2\right] &= \frac{1}{N^2}  \mathbb E\left[ \int_0^t \sum_{i\in \mathcal I(N,\tau), r(\bX^N_i(\tau))>0}  f(\bX_N^i(\tau))\mathcal R_N^{i^*,birth}(\intd \tau) \cdot\right.\\
%     & \qquad \left.\int_0^{t'} \sum_{j\in \mathcal I(N,\tau), r(\bX^j_N(\tau'))>0}  f(\bX_N^j(\tau'))\mathcal R_N^{j^*,birth}(\intd \tau')\right]\\
%        & - \frac{2}{N^2} \mathbb E\left[\int_0^t \sum_{i\in \mathcal I(N,\tau), r(\bX^N_i(\tau))>0}  f(\bX_N^i(\tau))\mathcal R_N^{i^*,birth}(\intd \tau)\right.\\
%     & \qquad \left.\int_0^{t'} \sum_{j\in \mathcal I(N,\tau), r(\bX^j_N(\tau'))>0}  f(\bX_N^j(\tau'))r(\bX_N^j(\tau'))\intd \tau'\right]\\
%        & + \frac{1}{N^2} \left(\int_0^{t} \sum_{i\in \mathcal I(N,\tau), r(\bX^N_i(\tau))>0}  f(\bX_N^i(\tau))r(\bX_N^i(\tau))\intd \tau\right)^2\\
%        & = \frac{1}{N^2}I_1-\frac{2}{N^2}I_2+\frac{1}{N^2}I_3.
%     \end{aligned}
% \end{equation*}
% For $I_1$, we have 
% \begin{equation*}
%     \begin{aligned}
%         I_1 =& \sum_{i\in \mathcal I(N,\tau), r(\bX^N_i(\tau))>0}\int_0^{t}f^2(\bX_N^i(\tau))r(\bX_N^i(\tau)) \intd \tau\\
%         &+\left(\int_0^{t} \sum_{i\in \mathcal I(N,\tau), r(\bX^N_i(\tau))>0}  f(\bX_N^i(\tau))r(\bX_N^i(\tau))\intd \tau\right)^2.
%     \end{aligned}
% \end{equation*}
% For $I_2$, we have
% \begin{equation*}
%     \begin{aligned}
%         I_2 &= \left(\int_0^{t} \sum_{i\in \mathcal I(N,\tau), r(\bX^N_i(\tau))>0}  f(\bX_N^i(\tau))r(\bX_N^i(\tau))\intd \tau\right)^2.
%     \end{aligned}
% \end{equation*}
% Therefore, we have
% $$
% \begin{aligned}
% \mathbb E\left[ \left(\mathcal M^2(t,f) \right)^2\right] &= \frac{1}{N^2}\sum_{i\in \mathcal I(N,\tau), r(\bX^N_i(\tau))>0}\int_0^{t}f^2(\bX_N^i(\tau))r(\bX_N^i(\tau)) \intd \tau\\
%     & = \frac{t C_{f,r}}{N},
% \end{aligned}
% $$
% where $C_{f,r}$ is a constant dependent on $f$ and $r$. Similarly, we can show that 
% $$
% \begin{aligned}
% \mathbb E\left[ \mathcal M^3(t,f) \right] 
%     & = 0,
% \end{aligned}
% $$
% and 
% $$
% \begin{aligned}
% \mathbb E\left[ \left(\mathcal M^3(t,f) \right)^2\right] 
%     & = \frac{t C_{f,r}}{N}.
% \end{aligned}
% $$
% % As a result, we have
% % \begin{equation*}
% %     \begin{aligned}
% %          \langle f(t,\cdot),\mu^N(t,\cdot)\rangle = &\langle f(0,\cdot), \mu^N(0,\cdot)\rangle \\
% %          & + \int_0^t\langle -\hat a(\cdot)\cdot \nabla f (\tau,\cdot) + \hat d(\tau, \cdot) \Delta f(\tau,\cdot) + r(\cdot) ,\mu^N(\tau,\cdot)\rangle \intd \tau\\
% %          & + \int_0^t\langle \frac{\partial}{\partial \tau} f(\tau,\cdot),\mu^N(\tau,\cdot)\rangle \intd \tau \\
% %          & + \mathcal M^1+\mathcal M^2 + \mathcal M^3.
% %     \end{aligned}
% % \end{equation*}
% Notice that 
% \begin{equation*}
%     \begin{aligned}
%          \langle f,\mu(t,\cdot)-\hat u(t,\cdot) \rangle = &\langle f, \mu^N(0,\cdot)-\hat u(0,\cdot)\rangle \\
%          &  +\int_0^t\langle -\hat a(\tau ,\cdot )\cdot \nabla f  + \hat d(\tau, \cdot) \Delta f+r(\tau,\cdot),\mu(\tau,\cdot)-\hat u(\tau,\cdot)\rangle\intd \tau\\
%          & + \int_0^t\langle f ,\frac{\partial}{\partial \tau} \mu(\tau,\cdot)-\frac{\partial}{\partial \tau} \hat u(\tau,\cdot)\rangle\intd \tau\\
%          & + \mathcal M^1+\mathcal M^2 + \mathcal M^3.
%     \end{aligned}
% \end{equation*}
% Therefore, by taking the expectation, we have 
% \begin{equation}
% \begin{aligned}
%         \mathbb E \left[\left\langle f,\mu(t,\cdot) \right\rangle\right] = &
%           \mathbb E \left[\langle f, \mu^N(0,\cdot) \rangle\right] \\
%          &  +\int_0^t\mathbb E \left[ \langle -\hat a(\tau,\cdot)\cdot \nabla f  + \hat d(\tau,\cdot) \Delta f +\hat r(\tau,\cdot),\mu(\tau,\cdot)\rangle\right]\intd \tau\\
%          & + \int_0^t\mathbb E \left[ \langle f ,\frac{\partial}{\partial \tau} \mu(\tau,\cdot)\rangle\right]\intd \tau\\
% \end{aligned}
% \end{equation}
% Based on the uniqueness of the weak solution of the \eqref{equ:linearlized} in Lemma \ref{lemma:well_posedness}, we  have 
% \begin{equation}
%     \mathbb E \left[\left\langle f,\mu(t,\cdot)-\hat u(t,\cdot) \right\rangle\right] = 0.
% \end{equation}
% For the second moment, the quadratic variation of each martingale term gives
% \[
% \mathbb{E}\!\left[ M_1(t,f)^2 \right] \;\le\; \frac{C_{f,\sigma}\,t}{N}, 
% \quad 
% \mathbb{E}\!\left[ M_2(t,f)^2 \right] \;\le\; \frac{C_{f,r}\,t}{N}, 
% \quad 
% \mathbb{E}\!\left[ M_3(t,f)^2 \right] \;\le\; \frac{C_{f,r}\,t}{N},
% \]
% where $C_{f,\sigma}$ and $C_{f,r}$ are constants depending on the test function 
% $f$, the bounded diffusion coefficient, and the reaction coefficient. 
% Cross-variations vanish by independence. Therefore,
% \[
% \mathbb{E}\!\left[\langle f(t), \mu^N(t) - \hat u(t) \rangle^2\right]
% = \mathbb{E}\!\left[ M_1(t,f)^2 + M_2(t,f)^2 + M_3(t,f)^2 \right] 
% \;\le\; \frac{C\,t}{N},
% \]
% which establishes the claimed $\mathcal{O}(1/N)$ rate.
% \end{proof}
% \newpage

% \begin{proof}
%     Using It\^o's formula, we have 
% \begin{equation*}
%     \begin{aligned}
%         \langle f(t,\cdot),\hat{\mu}^N(t,\cdot)\rangle = &\left\langle f(0,\cdot), \hat{\mu}^N(0,\cdot)\right\rangle - \int_0^t\left\langle a\left(\left[\mathcal P_K \hat \mu^N\right](\cdot) \right)\cdot \nabla f (\tau,\cdot),\hat{\mu}^N(\tau,\cdot)\right\rangle\intd \tau \\
%         &+ \int_0^t\left\langle d\left(\left[\mathcal P_K \hat \mu^N\right](\cdot) \right) \Delta f(\tau,\cdot),\hat \mu^N(\tau,\cdot)\right\rangle \intd \tau\\
%         &+ \int_0^t\left\langle  \frac{\partial}{\partial \tau} f(\tau,\cdot),\hat \mu^N(\tau,\cdot)\right\rangle \intd\tau\\
%         &+ \frac{1}{N} \int_0^t \sum_{i\in \mathcal I (N,\tau)} \sqrt{2 d\left(\left[\mathcal P_K \hat \mu^N\right]\left(\bX_N^i(\tau)\right) \right)} \nabla f(\tau,\bX_N^i(\tau))\cdot \intd \bW_N^i(\tau)\\
%         &+  \frac{1}{N} \int_0^t \sum_{i\in \mathcal I (N,\tau), r\left(\left[\mathcal P_K \hat \mu^N\right]\left(\bX_N^i(\tau)\right) \right)>0} f(\tau,\bX_N^i(\tau))\mathcal R_N^{i^*,birth}(\intd \tau)\\
%         &-  \frac{1}{N} \int_0^t \sum_{i\in \mathcal I (N,\tau), r\left(\left[\mathcal P_K \hat \mu^N\right]\left(\bX_N^i(\tau)\right) \right)<0} f(\tau,\bX_N^i(\tau)) \mathcal R_N^{i^*,death}(\intd \tau)
%     \end{aligned}
% \end{equation*}
% Subtract the exact solution, gives
% \begin{equation*}
%     \begin{aligned}
%         &\langle f(t,\cdot),\hat{\mu}^N(t,\cdot)-u(t,\cdot)\rangle \\= &\left\langle f(0,\cdot), \hat{\mu}^N(0,\cdot)-u(0,\cdot)\right\rangle \\
%         & - \int_0^t\left\langle a\left(\left[\mathcal P_K \hat \mu^N\right](\cdot) \right)\cdot \nabla f (\tau,\cdot),\hat{\mu}^N(\tau,\cdot)\right\rangle - \left\langle a(u(\tau,\cdot))\cdot \nabla f(\tau,\cdot) ,u(\tau,\cdot) \right\rangle \intd \tau \\
%         &+ \int_0^t\left\langle d\left(\left[\mathcal P_K \hat \mu^N\right](\cdot) \right) \Delta f(\tau,\cdot),\hat \mu^N(\tau,\cdot)\right\rangle - 
%         \left\langle d\left(u(\tau,\cdot)  \right) \Delta f(\tau,\cdot),u(\tau,\cdot)\right\rangle
%         \intd \tau\\
%         &+ \int_0^t\left\langle  \frac{\partial}{\partial \tau} f(\tau,\cdot),\hat \mu^N(\tau,\cdot) -  u(\tau,\cdot)\right\rangle \intd\tau\\
%         &+ \frac{1}{N} \int_0^t \sum_{i\in \mathcal I (N,\tau)} \sqrt{2 d\left(\left[\mathcal P_K \hat \mu^N\right]\left(\bX_N^i(\tau)\right) \right)} \nabla f(\tau,\bX_N^i(\tau))\cdot \intd \bW_N^i(\tau)\\
%         &+  \frac{1}{N} \int_0^t \sum_{i\in \mathcal I (N,\tau), r\left(\left[\mathcal P_K \hat \mu^N\right]\left(\bX_N^i(\tau)\right) \right)>0} f(\tau,\bX_N^i(\tau))\mathcal R_N^{i^*,birth}(\intd \tau) \\
%         & - \int_0^t  \int_{r(u(\tau,\bx))>0} f(\tau,\cdot), r(u(\tau,\cdot))u(\tau,\cdot)\intd \bx \intd \tau \\
%         &-  \frac{1}{N} \int_0^t \sum_{i\in \mathcal I (N,\tau), r\left(\left[\mathcal P_K \hat \mu^N\right]\left(\bX_N^i(\tau)\right) \right)<0} f(\tau,\bX_N^i(\tau)) \mathcal R_N^{i^*,death}(\intd \tau)\\
%         & - \int_0^t  \int_{r(u(\tau,\bx))<0} f(\tau,\cdot), r(u(\tau,\cdot))u(\tau,\cdot)\intd \bx \intd \tau \\
%         & = \mathcal  E_0 + \mathcal E_a + \mathcal E_d + \mathcal E_f + \mathcal N_w +\mathcal N_p +\mathcal N_n.
%     \end{aligned}
% \end{equation*}
% For the $\mathcal E_a$, we have that
% \[
% \begin{aligned}
% \mathcal E_a = &  \int_{0}^t \left\langle a\left(\left[\mathcal P_K \hat \mu^N\right](\cdot) \right)\cdot \nabla f (\tau,\cdot),\hat{\mu}^N(\tau,\cdot)\right\rangle - \left\langle a(u(\tau,\cdot))\cdot \nabla f(\tau,\cdot) ,u(\tau,\cdot) \right\rangle \intd \tau\\
%      = & \int_{0}^t \left\langle a\left(\left[\mathcal P_K \hat \mu^N\right](\cdot) \right)\cdot \nabla f (\tau,\cdot),\hat{\mu}^N(\tau,\cdot)  - u(\tau,\cdot) \right\rangle  \intd \tau
%     \\ & + \int_{0}^t  \left\langle \left(a\left(\left[\mathcal P_K \hat \mu^N\right](\cdot) \right) -a(u(\tau,\cdot))\right)\cdot \nabla f(\tau,\cdot) , u(\tau,\cdot) \right\rangle \intd \tau
% \end{aligned}
% \]

% \end{proof}

% \newpage

%


The weak formulation \eqref{eq:weak-form} defines the continuous-time particle
process $\mu^N(t)$: its transport is driven by the true Brownian motions
$\mathbf W^N_i$, and its reactions by the compensated Poisson martingales
$\widetilde{\mathcal R}^{+,N}_i,\widetilde{\mathcal R}^{-,N}_i$ of
\eqref{eq:weak-form}. This process evolves in continuous time; the step size
$\tau$ does not enter its definition.

The implemented algorithm, by contrast, advances in discrete steps and produces
the measure $\mu^N_n$ at $t_n=n\tau$ by composing, at each step, one
Euler--Maruyama increment for transport with one frozen-rate Poisson branching
for reaction. The time-discretization error is the gap between these two processes, $\mu^N(t_n)-\mu^N_n$, so the two must be compared
pathwise.

\emph{Coupling convention.} To enable such a pathwise comparison, we realize the
continuous process $\mu^N$, the discrete algorithm $\mu^N_n$, and all the
auxiliary flows introduced below on a single probability space, the common noise
space of the graphical (Harris) construction~\cite{harris1972nearest,harris1978additive}: the same Brownian streams $\mathbf W^N_i$ drive
transport in every flow, and the same Poisson random measures drive every
branching mechanism. This is a coupling convention, not a structural hypothesis
on the dynamics --- it fixes how the processes share randomness so that their
difference is meaningful pathwise, while leaving each process's law unchanged.
Concretely, under this convention,
\begin{itemize}
\item[(C1)] the Euler--Maruyama increment over $[t_n,t_{n+1}]$ uses the true
Brownian increment, $\sqrt{\tau}\,\boldsymbol\xi^i_{N,n}
:=\big(\mathbf W^N_i(t_{n+1})-\mathbf W^N_i(t_n)\big)$;
\item[(C2)] the one-step branching of particle $i$ is driven by the same Poisson random
measures $\mathcal N^\pm_i$ as the continuous process in \eqref{eq:weak-form};
the only change is that the time-varying continuous intensity
$r(\mathcal P_K\mu^N(\theta))$ is \emph{frozen} at its post-transport
left-endpoint value $c_i= r\!\big(\mathcal P_K\mu^{N,*}_n(\mathbf X^i)\big)$
over $[t_n,t_{n+1}]$. Thus the continuous time-integrated directing intensity
\[
\Lambda^{\mathrm{cont}}_i=\int_{t_n}^{t_{n+1}}
\mathbf 1_{S(N,\tau)}(i)\,r\!\big(\mathcal P_K\mu^N(\theta)\big)\,\mathrm d\theta
\]
is replaced by the frozen value $c_i\tau$, exactly as the EM increment in (C1)
replaces the exact transport drift by its left-endpoint value while reading the
same Brownian increment. The clock is sign-split so that birth and death are
driven by nonnegative intensities: if $c_i\ge0$, the $+$-clock fires with
cumulative intensity $\Lambda^{\mathrm{disc}}_{i,+}=e^{c_i\tau}-1$ and particle
$i$ outputs $1+Q^+_i(\Lambda^{\mathrm{disc}}_{i,+})$ copies; if $c_i<0$, the
$-$-clock applies the constant killing rate $-c_i$, so the particle survives iff
$Q^-_i\big((-c_i)\tau\big)=0$. In both branches the discrete update has the same
conditional first moment $e^{c_i\tau}$ as one step of the exact reaction subflow.
\end{itemize}
Under this convention the two processes differ \emph{only} through the
discretization of their generators, not through independent randomness. We have the following proposition

\begin{assumption}\label{ass:time}
Fix a finite horizon $[0,T]$, a spectral cutoff $K\ge1$, and let
$\mathbb T^d=(\mathbb R/2\pi\mathbb Z)^d$. Set $m:=\lceil s\rceil+1$. We assume:
\begin{itemize}
\item[(B1)] $a\in C^m_b(\mathbb R;\mathbb R^d)$ and $r\in C^m_b(\mathbb R)$;
in particular both are bounded and Lipschitz with a common constant
$C_{b,\mathrm{lip}}$.
\item[(B2)] along the dynamics the projected empirical field is uniformly bounded,
$\sup_{0\le t\le T}\|\mathcal P_K\mu^N(t)\|_{L^\infty}\le C_R$ (note $\mu^N$ is a
measure and need not possess a density; only its smooth projection
$\mathcal P_K\mu^N$ enters the coefficients).
\item[(B3)] the population is uniformly bounded, $|\mathcal I(N,t)|\le C_N N$ for all
$t\in[0,T]$.
\item[(B4)] $s>d/2+4$, so that $H^s\hookrightarrow C^4(\mathbb T^d)$ with
$\|f\|_{C^4}\le C_s\|f\|_{H^s}$, and $H^{s-1}$ is a multiplication algebra.
\end{itemize}
\end{assumption}
\begin{remark}[Admissible states and stopping]\label{rem:admissible}
Conditions (B2)--(B3) are most naturally read through the \emph{admissible class}
\[
\mathcal A_{K,R,N}:=\Bigl\{\nu=\tfrac1N\textstyle\sum_{i\in I}\delta_{y_i}:
I\subset\mathbb N,\ |I|\le C_NN,\ y_i\in\mathbb T^d,\ \|\mathcal P_K\nu\|_{L^\infty}\le C_R\Bigr\},
\]
and the two exit times
\[
\theta_{K,R,N}:=\inf\{t\ge0:\mu^N(t)\notin\mathcal A_{K,R,N}\},\qquad
\Theta_{K,R,N}:=\inf\{t\ge0:\text{some coupled flow used below exits }\mathcal A_{K,R,N}\}.
\]
Since the true process $\mu^N$ is one of the coupled flows,
$\Theta_{K,R,N}\le\theta_{K,R,N}$ pathwise. All local estimates of
Lemmas~\ref{lem:stability}--\ref{lem:branch} are first proved in \emph{stopped}
form on $\{t_n\le T\wedge\Theta_{K,R,N}\}$, where every coupled flow---the exact
Lie split $\widetilde\Phi^{N,\omega}_\tau$, the Euler--Maruyama transport
$\mathsf A^{N,\omega}_{\tau,h}$, the exact reaction subflow
$\mathsf B^{N,\omega}_\tau$, the frozen-rate branching
$\mathsf B^{N,\omega}_{\tau,h}$, and the intermediate flows in their proofs---
remains in $\mathcal A_{K,R,N}$, so that the coefficient bounds of
Lemma~\ref{lem:coeff} apply. Accumulating them yields the stopped bound
\eqref{eq:P1-stopped}, which we state explicitly in
Proposition~\ref{prop:time_discrete_particle} and suppress in the notation
thereafter; the population and field bounds thus serve only to keep the dynamics
in $\mathcal A_{K,R,N}$.

Imposed on the true dynamics alone, (B2)--(B3) give $\theta_{K,R,N}>T$ a.s., but
this does not imply the stronger auxiliary-admissibility condition
$\Theta_{K,R,N}>T$. The unstopped bound \eqref{eq:P1} follows precisely when
$\Theta_{K,R,N}>T$ a.s., which is secured by combining globally truncated
coefficients---smooth bounded extensions of $-a(\mathcal P_K\nu),
r(\mathcal P_K\nu)$ agreeing with them on $\mathcal A_{K,R,N}$, removing the
field constraint---with a population cutoff enforcing $|I|\le C_NN$ for all coupled
flows. The truncation alone does not control the population, since a bounded birth
rate can still exceed $C_NN$ particles in finite time; the cutoff closes this last
gap in applying Lemma~\ref{lem:coeff} to the auxiliary flows.
\end{remark}

\begin{proposition}
\label{prop:time_discrete_particle}
Let Assumption~\ref{ass:time} hold, let $\mu^N(t)$ be the weak solution of
\eqref{eq:weak-form} and $\mu^N_n$ the fully discrete scheme at $t_n=n\tau$, coupled on
the common noise space under the coupling convention above, with $\mu^N(0)=\mu^N_0$.
Write $\Theta:=\Theta_{K,R,N}$ (the auxiliary-flow exit time of
Remark~\ref{rem:admissible}, $\Theta_{K,R,N}\le\theta_{K,R,N}$ pathwise) and
$n_\Theta:=\lfloor\Theta/\tau\rfloor$ for the index
of the last grid point at or before the exit time. Then there exist $\tau_0>0$ and a
constant
\[
C_{T,K}=C_{T,K}(D,C_{b,\mathrm{lip}},C_N,C_R,s,d,T,K),
\]
independent of $N$ and $\tau$, such that for every $t_n\le T$ and $0<\tau\le\tau_0$ the
\emph{stopped} error satisfies, with no further assumption on the dynamics,
\begin{equation}\label{eq:P1-stopped}
\mathbb E\big\|\mu^N(t_{n\wedge n_\Theta})-\mu^N_{n\wedge n_\Theta}\big\|_{H^{-s}}^2
\le C_{T,K}\Big(\tau^2+\tfrac{\tau^{-1/2}}N\Big).
\end{equation}
If, in addition, $\Theta_{K,R,N}>T$ a.s.\ (the auxiliary-admissibility condition
of Remark~\ref{rem:admissible}), the stopping is vacuous on $[0,T]$, so
$n\wedge n_\Theta=n$ for all $t_n\le T$ and \eqref{eq:P1-stopped} holds for the
unstopped error:
\begin{equation}\label{eq:P1}
\mathbb E\big\|\mu^N(t_n)-\mu^N_n\big\|_{H^{-s}}^2
\le C_{T,K}\Big(\tau^2+\tfrac{\tau^{-1/2}}N\Big).
\end{equation}
Equivalently, the \emph{stopped} square-root error satisfies
\begin{equation}\label{eq:P2}
\Big(\mathbb E\big\|\mu^N(t_{n\wedge n_\Theta})-\mu^N_{n\wedge n_\Theta}\big\|_{H^{-s}}^2\Big)^{1/2}
\le C_{T,K}\Big(\tau+\tfrac1{\sqrt N\,\tau^{1/4}}\Big),
\end{equation}
and if $N\tau^{5/2}\ge c_*>0$ then
$\big(\mathbb E\|\mu^N(t_{n\wedge n_\Theta})-\mu^N_{n\wedge n_\Theta}\|_{H^{-s}}^2\big)^{1/2}\le C_{T,K,c_*}\,\tau$.
Under the auxiliary-admissibility condition $\Theta_{K,R,N}>T$ a.s.\ these bounds again hold unstopped.
\end{proposition}

The proof is organized into the lemmas below, combined at the end. Throughout,
$C$ denotes a constant independent of $N,\tau,K$, and $C_K$ a constant of the
form $C\,K^{p}$ for a fixed power $p$, whose precise value may change from line
to line and which we do not track. For a finite measure $\nu$ we write
\[
b_\nu:=-a(\mathcal P_K\nu),\qquad
c_\nu:=r(\mathcal P_K\nu),\qquad
R:=(I-\Delta)^{-s},
\]
so that $b_\nu$ is the reconstructed drift and $c_\nu$ the reconstructed reaction
rate associated with $\nu$.

\begin{lemma}[Coefficient bounds]\label{lem:coeff}
Under Assumption~\ref{ass:time}, with $b_\nu=-a(\mathcal P_K\nu)$ and
$c_\nu=r(\mathcal P_K\nu)$ as above, the following hold for every finite
measure $\nu$ arising along the dynamics and every multi-index $\beta$ with
$|\beta|=q\le m$:
\begin{equation}\label{equ:Bernstein_inequality}
\|\partial^\beta\mathcal P_K\nu\|_{L^\infty}
\le C_{q,d}\,K^{q}\,\|\mathcal P_K\nu\|_{L^\infty}
\le C_{q,d}\,C_R\,K^{q},
\end{equation}
with $C_{q,d}$ independent of $K$. Consequently
\begin{equation}\label{eq:Bb}
\|b_\nu\|_{C^m}\le C\,K^{m},\qquad \|c_\nu\|_{C^m}\le C\,K^{m}.
\end{equation}
Moreover, for any two such measures $\nu,\eta$ and $|\beta|=q\le m$,
\begin{equation}\label{eq:proj-diff}
\|\partial^\beta\mathcal P_K(\nu-\eta)\|_{L^\infty}
\le C\,K^{q+s+d/2}\,\|\nu-\eta\|_{H^{-s}},
\end{equation}
and the coefficient differences satisfy, for $0\le q\le m-1$,
\begin{equation}\label{eq:Blip}
\|b_\nu-b_\eta\|_{C^{q}}+\|c_\nu-c_\eta\|_{C^{q}}
\le C\,K^{\,q+s+d/2}\,\|\nu-\eta\|_{H^{-s}}.
\end{equation}
Finally, the birth/death rates are the positive and negative parts
$c^\pm_\nu:=(c_\nu)^\pm$, with $x\mapsto x^\pm=\max(\pm x,0)$. Wherever $c^\pm$ appear
below we use \emph{only} that they are bounded ($\|c^\pm_\nu\|_{L^\infty}\le\|c_\nu\|_{L^\infty}\le CK^{m}$)
and that $x\mapsto x^\pm$ is $1$-Lipschitz, so that
$|c^\sigma_\nu(x)-c^\sigma_\eta(y)|\le|c_\nu(x)-c_\eta(y)|$ for $\sigma\in\{+,-\}$;
we never claim $C^m$-regularity of $c^\pm$, which fails in general even when
$c_\nu\in C^m$ (the parts $x^\pm$ are merely Lipschitz, not $C^2$).
\end{lemma}

\begin{proof}
Since $\mathcal P_K\nu\in\mathcal T_K$,
it is a trigonometric polynomial of degree $\le K$ in each coordinate. By the
one-dimensional Bernstein--Zygmund inequality~\cite[Theorem~3.13,
Chapter~X]{zygmund2002trigonometric}, every univariate trigonometric polynomial
$T_n$ of degree $\le n$ satisfies $\|T_n'\|_{L^\infty}\le n\,\|T_n\|_{L^\infty}$,
and hence $\|T_n^{(\ell)}\|_{L^\infty}\le n^{\ell}\,\|T_n\|_{L^\infty}$ for
$\ell\in\mathbb N$. Applying this coordinatewise, for any multi-index $\beta$ with
$|\beta|=q$ and any $f\in\mathcal T_K$,
\[
\|\partial^\beta f\|_{L^\infty}
\le K^{\beta_1}\cdots K^{\beta_d}\,\|f\|_{L^\infty}
=K^{q}\,\|f\|_{L^\infty}.
\]
Taking $f=\mathcal P_K\nu$ gives the first inequality in
\eqref{equ:Bernstein_inequality}; the second then follows from the uniform
density bound~(B2).

Differentiating the composition $b_\nu=-a(\mathcal P_K\nu)$ up to order $m$ (chain
and product rules), each $\partial^\beta b_\nu$ with $|\beta|=q\le m$ is a finite
sum of terms $a^{(j)}(\mathcal P_K\nu)\prod_\ell\partial^{\beta_\ell}\mathcal P_K\nu$
with $\sum_\ell|\beta_\ell|=q$ and $j\le q\le m$. By (B1) the prefactors satisfy
$\|a^{(j)}\|_{L^\infty}\le C$, and by \eqref{equ:Bernstein_inequality} each factor
obeys $\|\partial^{\beta_\ell}\mathcal P_K\nu\|_{L^\infty}\le CK^{|\beta_\ell|}$;
since the orders sum to $q\le m$, $\|b_\nu\|_{C^m}\le CK^m$. The same argument
gives $\|c_\nu\|_{C^m}\le CK^m$, proving \eqref{eq:Bb}.

For \eqref{eq:proj-diff}, write $g:=\nu-\eta$. By
\eqref{equ:Bernstein_inequality} applied to the signed measure $g$,
$\|\partial^\beta\mathcal P_K g\|_{L^\infty}\le C_{q,d}K^{q}\|\mathcal P_K g\|_{L^\infty}$.
Expanding in Fourier modes and using Cauchy--Schwarz,
\[
\begin{aligned}
\|\mathcal P_K g\|_{L^\infty}
&\le\sum_{|\mathbf k|_\infty\le K}|\hat g(\mathbf k)|
\le\Bigl(\sum_{|\mathbf k|_\infty\le K}(1+|\mathbf k|^2)^{s}\Bigr)^{1/2}
   \Bigl(\sum_{\mathbf k}(1+|\mathbf k|^2)^{-s}|\hat g(\mathbf k)|^2\Bigr)^{1/2}\\
&\le\bigl(C\,K^d\!\cdot\!(CK)^{2s}\bigr)^{1/2}\|g\|_{H^{-s}}
=C\,K^{s+d/2}\,\|g\|_{H^{-s}},
\end{aligned}
\]
using $(1+|\mathbf k|^2)^{1/2}\le CK$ and $\#\{|\mathbf k|_\infty\le K\}\le CK^d$
in the first factor and the Fourier definition
$\|g\|_{H^{-s}}^2=\sum_{\mathbf k}(1+|\mathbf k|^2)^{-s}|\hat g(\mathbf k)|^2$ in
the second. Substituting this into the Bernstein bound
$\|\partial^\beta\mathcal P_K g\|_{L^\infty}\le K^{q}\|\mathcal P_K g\|_{L^\infty}$
obtained above yields \eqref{eq:proj-diff}.

Finally we prove \eqref{eq:Blip} for $b$. We take $q=1$ and
bound the two parts of the $C^1$ norm,
$\|b_\nu-b_\eta\|_{C^1}=\|b_\nu-b_\eta\|_{L^\infty}+\|\nabla b_\nu-\nabla b_\eta\|_{L^\infty}$,
in turn. For the zeroth-order part,
\[
\|b_\nu-b_\eta\|_{L^\infty}
=\|a(\mathcal P_K\nu)-a(\mathcal P_K\eta)\|_{L^\infty}
\le\|a'\|_{L^\infty}\,\|\mathcal P_K(\nu-\eta)\|_{L^\infty}
\le C\,K^{s+d/2}\,\|\nu-\eta\|_{H^{-s}},
\]
by the Lipschitz continuity of $a$ and \eqref{eq:proj-diff} with $q=0$.
At first order, the chain rule gives
$\nabla b_\mu=-a'(\mathcal P_K\mu)\nabla\mathcal P_K\mu$, so
\[
\nabla b_\nu-\nabla b_\eta
=-\bigl(a'(\mathcal P_K\nu)-a'(\mathcal P_K\eta)\bigr)\nabla\mathcal P_K\nu
 -a'(\mathcal P_K\eta)\,\nabla\mathcal P_K(\nu-\eta),
\]
which we bound in $L^\infty$ term by term. For the first term,
\[
\bigl\|(a'(\mathcal P_K\nu)-a'(\mathcal P_K\eta))\nabla\mathcal P_K\nu\bigr\|_{L^\infty}
\le\|a''\|_{L^\infty}\,\|\mathcal P_K(\nu-\eta)\|_{L^\infty}\,
   \|\nabla\mathcal P_K\nu\|_{L^\infty}
\le C\,K^{1+s+d/2}\,\|\nu-\eta\|_{H^{-s}},
\]
by the Lipschitz continuity of $a'$, \eqref{eq:proj-diff} with $q=0$, and
\eqref{equ:Bernstein_inequality} with $q=1$. For the second term,
\[
\bigl\|a'(\mathcal P_K\eta)\,\nabla\mathcal P_K(\nu-\eta)\bigr\|_{L^\infty}
\le\|a'\|_{L^\infty}\,\|\nabla\mathcal P_K(\nu-\eta)\|_{L^\infty}
\le C\,K^{1+s+d/2}\,\|\nu-\eta\|_{H^{-s}},
\]
by the boundedness of $a'$ and \eqref{eq:proj-diff} with $q=1$. Adding
the zeroth- and first-order bounds gives \eqref{eq:Blip} for $q=1$.

For general $q\le m-1$ the same mechanism applies to the higher-order expansion
of $\partial^\beta b_\nu-\partial^\beta b_\eta$: after telescoping, each term
contains exactly one difference factor---either
$a^{(\ell)}(\mathcal P_K\nu)-a^{(\ell)}(\mathcal P_K\eta)$ (Lipschitz,
$\ell\le m-1$, contributing $K^{s+d/2}$ via \eqref{eq:proj-diff} with $q=0$) or
$\partial^{\beta_\ell}\mathcal P_K(\nu-\eta)$ (contributing $K^{|\beta_\ell|+s+d/2}$)---while
the remaining factors are bounded by $CK^{|\cdot|}$ via
\eqref{equ:Bernstein_inequality}. Since the derivative orders sum to $q$, every
term is $\le CK^{q+s+d/2}\|\nu-\eta\|_{H^{-s}}$, giving \eqref{eq:Blip} for $b$.
Replacing $a$ by $r$ throughout, the same argument yields the
corresponding bound for $c$, completing the proof of \eqref{eq:Blip}.
\end{proof}

We denote by $\Phi^{N,\omega}_\tau$ the exact full stochastic particle flow of the
weak formulation \eqref{eq:weak-form} over one step: a random map, driven by the
Brownian motions and Poisson clocks (C1)--(C2), that evolves a given particle
measure by time $\tau$. The next lemma is a stability estimate for this flow with
respect to its initial data.

The proof uses the same $H^{-s}$ energy method as
Proposition~\ref{prop:u_K_mu}, but compares two different objects, so we state and
prove it separately. Proposition~\ref{prop:u_K_mu} compares the \emph{stochastic}
particle process with the \emph{deterministic} projected solution $u_K$, so its
only martingale comes from the single process $\mu^N$. Here, by contrast, both
$\Phi^{N,\omega}_\tau(\mu)$ and $\Phi^{N,\omega}_\tau(\mu')$ are stochastic, and
the energy balance involves the difference of two such martingales.

\begin{lemma}[One-step stability of the exact flow]\label{lem:stability}
Under Assumption~\ref{ass:time}, let $\mu,\mu'$ be two particle measures driven by
the common graphical noise (C1)--(C2), without labelwise correspondence assumed
between them. Then there is $C_K$ such that, for $\tau\le\tau_0$ with $C_K\tau_0\le1$,
\begin{equation}\label{eq:stab}
\mathbb E\Bigl[\|\Phi^{N,\omega}_\tau(\mu)-\Phi^{N,\omega}_\tau(\mu')\|_{H^{-s}}^2
\,\Big|\,\mathcal F\Bigr]
\le(1+C_K\tau)\,\|\mu-\mu'\|_{H^{-s}}^2+C_K\,\frac{\tau}{N},
\end{equation}
where $\mathcal F$ is the $\sigma$-algebra of the data at the start of the step.
The additive term $C_K\tau/N$ vanishes when $\mu,\mu'$ coincide labelwise.
\end{lemma}

\begin{proof}
For $0\le\theta\le\tau$ let $\mu_\theta,\mu'_\theta$ be the exact evolutions of
\eqref{eq:weak-form} from $\mu,\mu'$ under the same drivers, and
$\Delta_\theta:=\mu_\theta-\mu'_\theta$, so $\Delta_0=\mu-\mu'$ and
$\Delta_\tau=\Phi^{N,\omega}_\tau(\mu)-\Phi^{N,\omega}_\tau(\mu')$. Write
$\mathbb E[\cdot\mid\mathcal F]=\mathbb E_{\mathcal F}[\cdot]$. For $f\in H^s$,
subtracting the weak formulations,
\[
\mathrm d\langle f,\Delta_\theta\rangle
=\bigl[\langle\nabla f,b_{\mu_\theta}\mu_\theta-b_{\mu'_\theta}\mu'_\theta\rangle
+\langle\Delta f,D\Delta_\theta\rangle
+\langle f,c_{\mu_\theta}\mu_\theta-c_{\mu'_\theta}\mu'_\theta\rangle\bigr]\mathrm d\theta
+\mathrm dM_f(\theta),
\]
$M_f=M^\mu_f-M^{\mu'}_f$ being the difference of the two flows' martingale parts. Let
$\phi_\theta:=R\Delta_\theta$, so $\|\phi_\theta\|_{H^s}^2
=\langle\Delta_\theta,\phi_\theta\rangle=\|\Delta_\theta\|_{H^{-s}}^2$. It\^o's
formula applied to $\|\Delta_\theta\|_{H^{-s}}^2$ and the conditional expectation
$\mathbb E_{\mathcal F}$ remove the martingale, leaving
\[
\tfrac{\mathrm d}{\mathrm d\theta}\mathbb E_{\mathcal F}\|\Delta_\theta\|_{H^{-s}}^2
=\mathbb E_{\mathcal F}[J_1+J_2+J_3+J_4+J_5],
\]
with $J_1=\langle\nabla\phi_\theta,b_{\mu_\theta}\mu_\theta-b_{\mu'_\theta}\mu'_\theta\rangle$,
$J_2=D\langle\Delta\phi_\theta,\Delta_\theta\rangle$,
$J_3=\langle\phi_\theta,c_{\mu_\theta}\mu_\theta-c_{\mu'_\theta}\mu'_\theta\rangle$,
and $J_4,J_5$ the diffusion and branching quadratic variations.

\emph{Advection $J_1$.} Splitting
$b_{\mu_\theta}\mu_\theta-b_{\mu'_\theta}\mu'_\theta
=b_{\mu'_\theta}\Delta_\theta+(b_{\mu_\theta}-b_{\mu'_\theta})\mu_\theta$ gives
$J_1=J_1^{\mathrm m}+J_1^{\mathrm c}$. For the measure part,
\[
|J_1^{\mathrm m}|
=|\langle b_{\mu'_\theta}\nabla\phi_\theta,\Delta_\theta\rangle|
\le\|b_{\mu'_\theta}\nabla\phi_\theta\|_{H^{s-1}}\|\Delta_\theta\|_{H^{-s+1}}
\le C_K\|\phi_\theta\|_{H^s}\|\Delta_\theta\|_{H^{-s+1}},
\]
by moving $b_{\mu'_\theta}$ onto the first factor, $\langle H^{s-1},H^{-s+1}\rangle$
duality, and \eqref{eq:Bb} together with
$\|\nabla\phi_\theta\|_{H^{s-1}}\le\|\phi_\theta\|_{H^s}$. For the coefficient part,
\[
|J_1^{\mathrm c}|
\le\|(b_{\mu_\theta}-b_{\mu'_\theta})\nabla\phi_\theta\|_{H^{s-1}}\|\mu_\theta\|_{H^{-s+1}}
\le C_K\|\phi_\theta\|_{H^s}\|\Delta_\theta\|_{H^{-s}},
\]
by $\|\mu_\theta\|_{H^{-s+1}}\le C_N$ (mass $\le C_N$, $s-1>d/2$) and
\eqref{eq:Blip} with $q=m-1$.

\emph{Reaction $J_3$.} Likewise $J_3=J_3^{\mathrm m}+J_3^{\mathrm c}$; no
derivative falls on $\phi_\theta$, so both close in $H^{-s}$. For the measure part,
\[
|J_3^{\mathrm m}|
=|\langle c_{\mu'_\theta}\phi_\theta,\Delta_\theta\rangle|
\le\|c_{\mu'_\theta}\|_{C^{\lceil s\rceil}}\|\phi_\theta\|_{H^s}\|\Delta_\theta\|_{H^{-s}}
\le C_K\|\Delta_\theta\|_{H^{-s}}^2,
\]
by \eqref{eq:Bb} and $\|\phi_\theta\|_{H^s}=\|\Delta_\theta\|_{H^{-s}}$. For the
coefficient part,
\[
|J_3^{\mathrm c}|
\le\|(c_{\mu_\theta}-c_{\mu'_\theta})\phi_\theta\|_{H^s}\|\mu_\theta\|_{H^{-s}}
\le C_K\|\Delta_\theta\|_{H^{-s}}^2,
\]
by $\|\mu_\theta\|_{H^{-s}}\le C_N$ and \eqref{eq:Blip} with
$q=\lceil s\rceil\le m-1$.

\emph{Diffusion $J_2$.} Writing $\Delta=I-(I-\Delta)$,
\[
J_2=2D\|\Delta_\theta\|_{H^{-s}}^2-D\|\Delta_\theta\|_{H^{-s+1}}^2.
\]
Only $J_1^{\mathrm m}$ carries $\|\Delta_\theta\|_{H^{-s+1}}$, and Young's
inequality with weight $\epsilon=D/C_K$ gives
\[
C_K\|\phi_\theta\|_{H^s}\|\Delta_\theta\|_{H^{-s+1}}
\le\tfrac{C_K^2}{2D}\|\phi_\theta\|_{H^s}^2+\tfrac{D}{2}\|\Delta_\theta\|_{H^{-s+1}}^2
=\tfrac{C_K^2}{2D}\|\Delta_\theta\|_{H^{-s}}^2+\tfrac{D}{2}\|\Delta_\theta\|_{H^{-s+1}}^2,
\]
whose second term is absorbed by the negative diffusion contribution in $J_2$.
Hence $\mathbb E_{\mathcal F}[J_1+J_2+J_3]\le C_K\mathbb E_{\mathcal F}\|\Delta_\theta\|_{H^{-s}}^2$.

\emph{Quadratic variations $J_4,J_5$.} The two processes have distinct particle
populations $\mathcal I^\mu_\theta,\mathcal I^{\mu'}_\theta$ with no assumed
correspondence, so their martingales are driven by independent noise; we
bound the quadratic variation of the difference by the triangle inequality
$\langle M^\mu_f-M^{\mu'}_f\rangle_\theta\le2\langle M^\mu_f\rangle_\theta+2\langle M^{\mu'}_f\rangle_\theta$
and compute the energy quadratic variation of each single process. As in
Proposition~\ref{prop:u_K_mu}, the It\^o correction of $\|\cdot\|_{H^{-s}}^2$
contracts each diffusion increment / branching $\tfrac1N$-jump with the Hessian
$R=(I-\Delta)^{-s}$, producing the same \emph{$x$-independent} per-particle
constants $C_s'=\|\nabla(I-\Delta)^{-s/2}\delta_x\|_{L^2}^2$ ($s>d/2+1$) and
$C_s=\|\delta_x\|_{H^{-s}}^2$ ($s>d/2$) as in Proposition~\ref{prop:u_K_mu},
\emph{not} the random potential $\nabla\phi_\theta(x)$ or $\phi_\theta(x)$ evaluated
at the particle. Hence, by $|\mathcal I^\mu_\theta|,|\mathcal I^{\mu'}_\theta|\le C_NN$
and $|c_{\mu_\theta}|,|c_{\mu'_\theta}|\le C_K$,
\[
J_4\le\frac{4D}{N^2}\bigl(|\mathcal I^\mu_\theta|+|\mathcal I^{\mu'}_\theta|\bigr)C_s'
\le\frac{8DC_NC_s'}{N}=:\frac{C_4}{N},
\qquad
J_5\le\frac{2}{N^2}\bigl(|\mathcal I^\mu_\theta|+|\mathcal I^{\mu'}_\theta|\bigr)C_KC_s
\le\frac{4C_KC_NC_s}{N}=:\frac{C_5}{N}.
\]
Both are \emph{additive}, of order $1/N$ and \emph{independent of}
$\|\Delta_\theta\|$: with no label correspondence the two empirical processes read
independent noise, so even at $\Delta_0=\mu-\mu'=0$ the difference fluctuates at
order $\tau/N$, and no cancellation to $O(\tau/N^2)$ is available. (The
first-variation contribution $2\langle\frac1N\delta,\phi_\theta\rangle$ of each
branching event is already counted in $J_3$.) The diffusion constant
$C_4=8DC_NC_s'$ carries no power of $K$. Thus
$\mathbb E_{\mathcal F}[J_4+J_5]\le C_K/N$.

Combining the above,
$\tfrac{\mathrm d}{\mathrm d\theta}\mathbb E_{\mathcal F}\|\Delta_\theta\|_{H^{-s}}^2
\le C_K\mathbb E_{\mathcal F}\|\Delta_\theta\|_{H^{-s}}^2+C_K/N$. Gr\"onwall over
$[0,\tau]$ with $\Delta_0=\mu-\mu'$, $e^{C_K\tau}\le1+C_K\tau$ and
$(e^{C_K\tau}-1)/C_K\le\tau$ (since $C_K\tau_0\le1$), gives \eqref{eq:stab}, with
additive forcing $C_K\tau/N$.
\end{proof}

The error analysis below repeatedly compares two flows of the \emph{same}
particle system --- the exact flow against its Lie splitting, and against its
Euler--Maruyama and frozen-rate discretizations. For such a difference to be a
pathwise quantity, and not merely a comparison of laws, the two flows must be
driven by the same randomness.

\emph{Common noise space and graphical (Harris) construction.} All the continuous and discrete particle flows of this section
are realized on one probability space by a graphical construction in the sense of
Harris~\cite{harris1972nearest,harris1978additive}, in the measure-valued SDE form
of~\cite{fournier2004microscopic,bansaye2015stochastic}. Let
\[
\mathcal U_N:=\{1,\ldots,N\}\times\mathbb N^{<\infty}
\]
be the Ulam--Harris label set: each of the $N$ initial particles is a root, and a
label $\alpha$ spawns a child by appending a symbol. To each $\alpha\in\mathcal U_N$
we attach, \emph{independently across labels}, a $d$-dimensional Brownian stream
$\mathbf W^N_\alpha$ and a pair of unit-intensity Poisson random measures
$\mathcal N^{\pm}_\alpha(\mathrm dt,\mathrm du)$ on $[0,T]\times[0,\infty)$ (here $t$ is
time and $u\ge0$ is an auxiliary height used for thinning, not a spatial variable)
governing births ($+$) and deaths ($-$); $(\mathcal H_t)_{t\ge0}$ is the filtration generated by
the initial configuration and all these noises. A particle system is a labeled
configuration $\Xi=\{(\alpha,x_\alpha):\alpha\in I\}$ with finite alive label set
$I\subset\mathcal U_N$ and $x_\alpha\in\mathbb T^d$, projected to the empirical measure
$\Pi(\Xi):=\frac1N\sum_{\alpha\in I}\delta_{x_\alpha}$; we write
$\Phi^{N,\omega}_\tau(\mu)$ for $\Pi(\Phi^{N,\omega}_\tau(\Xi))$ with $\Pi(\Xi)=\mu$
(and likewise for the substep flows $\mathsf A^{N,\omega}_\tau,\mathsf B^{N,\omega}_\tau,
\mathsf A^{N,\omega}_{\tau,h},\mathsf B^{N,\omega}_{\tau,h}$ defined below).

Each flow runs with its own predictable rate field $c$, the reconstructed reaction rate
$c_\nu=r(\mathcal P_K\nu)$ along the flow's running measure $\nu$; write
$c^\pm_t(x):=\max(\pm c_t(x),0)$ for its positive and negative parts, both bounded
pointwise by $C_K$ (by~\eqref{eq:Bb}). Given $c$, the births and deaths of a label
$\alpha$ are the points of $\mathcal N^\pm_\alpha$ lying below the rate: $\alpha$ gives
birth at each point $(t,u)$ of $\mathcal N^{+}_\alpha$ with $u\le c^+_t(x_\alpha(t^-))$,
the new child taking the first unused child label of $\alpha$ and starting at
$x_\alpha(t^-)$, and dies at each point $(t,u)$ of $\mathcal N^{-}_\alpha$ with
$u\le c^-_t(x_\alpha(t^-))$; between jumps each alive particle moves by the transport SDE
driven by its own stream $\mathbf W^N_\alpha$. Equivalently, the number of accepted
births (resp.\ deaths) of $\alpha$ on $[0,t]$ is the time-changed unit-rate process
$Q^{\pm,N}_\alpha\!\bigl(\Lambda^{\pm,N}_\alpha(t)\bigr)$ of the weak
formulation~\eqref{eq:weak-form}, with
$\Lambda^{\pm,N}_\alpha(t)=\int_0^t c^\pm_s(x_\alpha(s^-))\,\mathrm ds$; in particular the
root label of particle $i\in\mathcal I_n$ carries the stream $\mathbf W^N_i$ and the
clocks $Q^{\pm,N}_i$ of~\eqref{eq:weak-form} and~(C1)--(C2).

Two flows of the same particle system are coupled by reading the \emph{same} noise
$\{\mathbf W^N_\alpha,\mathcal N^\pm_\alpha\}$ but with their own rate fields $c=c_\nu$ and
$\bar c=c_{\bar\nu}$, reconstructed along the two running measures $\nu,\bar\nu$. For each
sign $\sigma\in\{+,-\}$, write $c^\sigma,\bar c^\sigma$ for the $\sigma$-rates of a label
$\alpha$ in the two flows, each evaluated at that flow's position. A point $(t,u)$ of the
shared clock $\mathcal N^\sigma_\alpha$ is \emph{matched} if both flows accept it,
$u\le\min(c^\sigma,\bar c^\sigma)$, and \emph{surplus} if exactly one does,
$\min(c^\sigma,\bar c^\sigma)<u\le\max(c^\sigma,\bar c^\sigma)$.

We now apply the flow to the coupled processes. The exact full stochastic
particle flow $\Phi^{N,\omega}_\tau$ of the weak formulation
\eqref{eq:weak-form} advances the exact process by one step,
$\mu^N(t_{n+1})=\Phi^{N,\omega}_\tau(\mu^N(t_n))$; in $\Phi^{N,\omega}_\tau$
transport and branching occur \emph{simultaneously}. The implemented algorithm,
by contrast, performs them \emph{successively}: first an Euler--Maruyama
transport step, then a frozen-rate Poisson branching step. To compare them
without conflating the two effects, we introduce two exact substep flows
mirroring the algorithm's split structure:
\begin{itemize}
\item $\mathsf A^{N,\omega}_\tau$ --- the \emph{exact nonlinear transport--diffusion
subflow}, with no branching. Starting from $\nu_0=\mu^N_n$,
\[
\mathrm dY_i(\theta)=b_{\nu_\theta}(Y_i(\theta))\,\mathrm d\theta
+\sqrt{2D}\,\mathrm dW^N_i(t_n+\theta),
\qquad Y_i(0)=\mathbf X^i_n,
\qquad
\nu_\theta=\frac1N\sum_{i\in\mathcal I_n}\delta_{Y_i(\theta)},
\]
$0\le\theta\le\tau$; the drift is the reconstructed $b_{\nu_\theta}$ associated
with the running measure $\nu_\theta$ (not frozen).
\item $\mathsf B^{N,\omega}_\tau$ --- the \emph{exact pure reaction (branching)
subflow}, with no spatial motion: particles stay fixed at their starting
positions and undergo birth--death with rate $c_{\eta_\theta}
=r(\mathcal P_K\eta_\theta)$ associated with the running measure
$\eta_\theta$ of the pure reaction flow.
\end{itemize}
Their composition
\[
\widetilde\Phi^{N,\omega}_\tau:=\mathsf B^{N,\omega}_\tau\circ\mathsf A^{N,\omega}_\tau
\]
is the \emph{exact nonlinear Lie-split flow}: do exact transport--diffusion for
time $\tau$ (without branching), then exact pure branching for time $\tau$
(without motion). Both substep flows use the noise family (C1)--(C2). The
algorithm $\Psi^{N,\omega}_\tau$ defined in (i)--(ii) above is, in this notation,
the discrete counterpart
$\Psi^{N,\omega}_\tau=\mathsf B^{N,\omega}_{\tau,h}\circ\mathsf A^{N,\omega}_{\tau,h}$,
where:
\begin{itemize}
\item[(i)] $\mathsf A^{N,\omega}_{\tau,h}$ is the \emph{Euler--Maruyama transport
substep} with drift frozen at the left endpoint,
\[
\mathbf X^*_i
:=\mathbf X^i_n+b_{\mu^N_n}(\mathbf X^i_n)\,\tau
+\sqrt{2D}\,\bigl(\mathbf W^N_i(t_{n+1})-\mathbf W^N_i(t_n)\bigr),
\quad i\in\mathcal I_n,
\quad
\mu^{N,*}_n:=\frac1N\sum_{i\in\mathcal I_n}\delta_{\mathbf X^*_i};
\]
\item[(ii)] $\mathsf B^{N,\omega}_{\tau,h}$ is the \emph{frozen-rate Poisson
branching substep}: for each $i\in\mathcal I_n$, with
$c_i:=c_{\mu^{N,*}_n}(\mathbf X^*_i)$ and driven by $Q^{\pm,N}_i$, particle
$i$ undergoes a birth of cumulative intensity $e^{c_i\tau}-1$ (if $c_i\ge0$) or a
killing of cumulative intensity $(-c_i)\tau$ with survival probability
$e^{c_i\tau}$ (if $c_i<0$); the surviving and newly born particles stay at
$\mathbf X^*_i$, and
$\mu^N_{n+1}=\frac1N\sum_{i\in\mathcal I_{n+1}}\delta_{\mathbf X^*_i}$.
\end{itemize}
Under (C1)--(C2) all four maps act on the same input $\mu^N_n$ and share the
graphical noise fixed in the construction above.

We record the two filtrations used throughout the local one-step error analysis. Recall that
$\mathcal F_n:=\mathcal H_{t_n}$ is the $\sigma$-algebra of all information
available up to time $t_n$, i.e.\ \emph{before} the step-$n$ update;
write $\mathbb E_n[\cdot]:=\mathbb E[\cdot\mid\mathcal F_n]$ for the corresponding
conditional expectation. The larger $\sigma$-algebra
$\mathcal G_n$ additionally fixes the step-$n$ transport Brownian paths,
\begin{equation}\label{eq:filtrations}
\mathcal G_n:=\mathcal F_n\vee\sigma\bigl(\mathbf W^N_\alpha(t_n+\theta)-\mathbf W^N_\alpha(t_n):
0\le\theta\le\tau,\ \alpha\in\mathcal U_N\bigr),
\end{equation}
where $(\mathcal H_t)_{t\ge0}$ and the label set $\mathcal U_N$ are those of the
graphical construction above. Conditioning on $\mathcal G_n$ fixes the
post-transport positions $\mathbf X^*_i$, the post-EM measure $\mu^{N,*}_n$ and
the frozen rates $c_i$, leaving only the branching randomness; this is the
structure used in Lemmas~\ref{lem:Bstab},~\ref{lem:transport}
and~\ref{lem:branch}.

The local one-step defect
$\ell_{n+1}:=\Phi^{N,\omega}_\tau(\mu^N_n)-\Psi^{N,\omega}_\tau(\mu^N_n)$
admits the three-term decomposition obtained by inserting $\widetilde\Phi$ and
$\mathsf B^{N,\omega}_\tau\circ\mathsf A^{N,\omega}_{\tau,h}$:
\[
\ell_{n+1}=\ell^{\mathrm{split}}_{n+1}+\ell^{\mathrm{tr}}_{n+1}+\ell^{\mathrm{br}}_{n+1},
\]
where
\begin{align*}
\ell^{\mathrm{split}}_{n+1}
&:=\Phi^{N,\omega}_\tau(\mu^N_n)-\widetilde\Phi^{N,\omega}_\tau(\mu^N_n)
&&\text{(full simultaneous vs exact Lie split),}\\
\ell^{\mathrm{tr}}_{n+1}
&:=\mathsf B^{N,\omega}_\tau(\mathsf A^{N,\omega}_\tau(\mu^N_n))
   -\mathsf B^{N,\omega}_\tau(\mathsf A^{N,\omega}_{\tau,h}(\mu^N_n))
&&\text{(exact transport vs EM, then exact branching),}\\
\ell^{\mathrm{br}}_{n+1}
&:=\mathsf B^{N,\omega}_\tau(\mathsf A^{N,\omega}_{\tau,h}(\mu^N_n))
   -\mathsf B^{N,\omega}_{\tau,h}(\mathsf A^{N,\omega}_{\tau,h}(\mu^N_n))
&&\text{(exact pure reaction vs frozen-rate branching),}
\end{align*}
bounded in the three lemmas below.

\begin{lemma}[Finite-dimensional generator bound for the reaction-to-transport feedback]\label{lem:drift-feedback-generator}
On the stopped admissible set $\{t_n\le T\wedge\Theta_{K,R,N}\}$, set
\[
\nu^\Phi_\theta:=\Phi^{N,\omega}_\theta(\mu^N_n),\qquad
\nu^{\mathsf A}_\theta:=\mathsf A^{N,\omega}_\theta(\mu^N_n),\qquad 0\le\theta\le\tau,
\]
the empirical measures of the full simultaneous flow $\Phi$ and of the branching-free
transport subflow $\mathsf A$. For each root label $i\in\mathcal I_n$ let
\[
\begin{aligned}
\mathrm d\widehat Y_i(\theta)
&= b_{\nu^\Phi_\theta}(\widehat Y_i(\theta))\,\mathrm d\theta
   +\sqrt{2D}\,\mathrm dW_i(t_n+\theta),\\
\mathrm dY_i(\theta)
&= b_{\nu^{\mathsf A}_\theta}(Y_i(\theta))\,\mathrm d\theta
   +\sqrt{2D}\,\mathrm dW_i(t_n+\theta),
\end{aligned}
\qquad
\widehat Y_i(0)=Y_i(0)=\mathbf X^i_n,
\]
driven by the common stream $W_i$, with $\widehat Y_i$ a ghost root carried by the
full-flow drift $b_{\nu^\Phi_\theta}$ regardless of its genealogical survival. Set
\[
\ell^{\mathrm{drift}}_{n+1}
:=\frac1N\sum_{i\in\mathcal I_n}\bigl(\delta_{\widehat Y_i(\tau)}-\delta_{Y_i(\tau)}\bigr).
\]
Then there is $C_K$, independent of $N$ and $\tau$, such that for $0<\tau\le\tau_0(K)$
\[
\bigl\|\mathbb E_n\ell^{\mathrm{drift}}_{n+1}\bigr\|_{H^{-s}}\le C_K\tau^2,
\]
\[
\mathbb E_n\bigl\|\ell^{\mathrm{drift}}_{n+1}-\mathbb E_n\ell^{\mathrm{drift}}_{n+1}\bigr\|_{H^{-s}}^2
\le C_K\Bigl(\frac{\tau^3}{N}+\tau^4\Bigr)
\le C_K\Bigl(\frac{\tau^2}{N}+\tau^4\Bigr).
\]
\end{lemma}

\begin{proof}
Throughout,
$\mathcal M_K=\{\mathbf k\in\mathbb Z^d:|\mathbf k|_\infty\le K\}$ is the truncation index
set and $\{\psi_{\mathbf k}\}_{\mathbf k\in\mathcal M_K}$ the reconstruction basis, so that
$b_\nu=-a(\mathcal P_K\nu)$ depends on $\nu$ only through the finitely many modal
coordinates $\{\langle\psi_{\mathbf k},\nu\rangle\}_{\mathbf k\in\mathcal M_K}$.

Write
\[
Z_{\mathbf k}^\Phi(\theta):=\langle\psi_{\mathbf k},\nu^\Phi_\theta\rangle,\qquad
Z_{\mathbf k}^{\mathsf A}(\theta):=\langle\psi_{\mathbf k},\nu^{\mathsf A}_\theta\rangle,\qquad
\delta Z_{\mathbf k}(\theta):=Z_{\mathbf k}^\Phi(\theta)-Z_{\mathbf k}^{\mathsf A}(\theta),
\]
and set
$|\delta Z_\theta|_K^2:=\sum_{\mathbf k\in\mathcal M_K}|\delta Z_{\mathbf k}(\theta)|^2$.
Since $b_\nu=-a(\mathcal P_K\nu)$ depends on $\nu$ only through the coordinates
$\{Z_{\mathbf k}\}_{\mathbf k\in\mathcal M_K}$, Lemma~\ref{lem:coeff} and the equivalence of
norms on the finite-dimensional space $\mathcal T_K$ give, for $q\le2$,
\begin{equation}\label{eq:L4-coeff}
\|b_{\nu^\Phi_\theta}-b_{\nu^{\mathsf A}_\theta}\|_{C^q}\le C_K|\delta Z_\theta|_K .
\end{equation}

\medskip
\noindent\textbf{Step 1: finite-mode semimartingale bounds.}
For each $\mathbf k\in\mathcal M_K$, It\^o's formula (equivalently the weak
formulation~\eqref{eq:weak-form}) applied to the fixed smooth mode $\psi_{\mathbf k}$ gives,
for the full simultaneous flow,
\[
\mathrm dZ_{\mathbf k}^\Phi
=\bigl\langle D\Delta\psi_{\mathbf k}+\nabla\psi_{\mathbf k}\!\cdot\! b_{\nu^\Phi_\theta}
   +\psi_{\mathbf k}\,c_{\nu^\Phi_\theta},\,\nu^\Phi_\theta\bigr\rangle\,\mathrm d\theta
 +\mathrm dM_{\mathbf k}^\Phi(\theta),
\]
and, for the branching-free transport subflow (no reaction term),
\[
\mathrm dZ_{\mathbf k}^{\mathsf A}
=\bigl\langle D\Delta\psi_{\mathbf k}+\nabla\psi_{\mathbf k}\!\cdot\! b_{\nu^{\mathsf A}_\theta},\,
   \nu^{\mathsf A}_\theta\bigr\rangle\,\mathrm d\theta+\mathrm dM_{\mathbf k}^{\mathsf A}(\theta).
\]
Subtracting, $\delta Z_{\mathbf k}$ is the one-dimensional semimartingale
$\mathrm d\,\delta Z_{\mathbf k}=B_{\mathbf k}\,\mathrm d\theta+\mathrm dM_{\mathbf k}$, with
$M_{\mathbf k}=M_{\mathbf k}^\Phi-M_{\mathbf k}^{\mathsf A}$ and
\[
B_{\mathbf k}(\theta)
=\bigl\langle D\Delta\psi_{\mathbf k}+\nabla\psi_{\mathbf k}\!\cdot\! b_{\nu^\Phi_\theta},\,\nu^\Phi_\theta\bigr\rangle
-\bigl\langle D\Delta\psi_{\mathbf k}+\nabla\psi_{\mathbf k}\!\cdot\! b_{\nu^{\mathsf A}_\theta},\,\nu^{\mathsf A}_\theta\bigr\rangle
+\bigl\langle\psi_{\mathbf k}\,c_{\nu^\Phi_\theta},\,\nu^\Phi_\theta\bigr\rangle .
\]
Here $M_{\mathbf k}^\Phi$ collects the Brownian-transport integral
$\tfrac{\sqrt{2D}}N\sum_\alpha\nabla\psi_{\mathbf k}(x_\alpha)\!\cdot\!\mathrm dW_\alpha$ and the
compensated birth--death integrals against $\mathcal N^\pm_\alpha$, while
$M_{\mathbf k}^{\mathsf A}$ collects only the former. Each bracket of $B_{\mathbf k}$ is a finite
sum $\langle g,\nu\rangle=\frac1N\sum_\alpha g(x_\alpha)$ with $\|g\|_\infty\le C_K$ (from
$\|\psi_{\mathbf k}\|_{C^2},\|b_\nu\|_\infty,\|c_\nu\|_\infty\le C_K$ by~\eqref{eq:Bb})
against a measure of mass $\|\nu_\theta\|_{\mathrm{TV}}\le C_N$, so
\[
|B_{\mathbf k}(\theta)|
\le C_K\bigl(\|\nu^\Phi_\theta\|_{\mathrm{TV}}+\|\nu^{\mathsf A}_\theta\|_{\mathrm{TV}}\bigr)
\le C_KC_N \leq C_K.
\]
For the martingale, $\langle M_{\mathbf k}\rangle_\theta\le2\langle M_{\mathbf k}^\Phi\rangle_\theta+2\langle M_{\mathbf k}^{\mathsf A}\rangle_\theta$, with
\[
\frac{\mathrm d}{\mathrm d\theta}\langle M_{\mathbf k}^\Phi\rangle_\theta
=\frac{2D}{N^2}\sum_\alpha|\nabla\psi_{\mathbf k}(x_\alpha)|^2
+\frac1{N^2}\sum_\alpha\sum_{\sigma=\pm}c^\sigma_{\nu^\Phi_\theta}(x_\alpha)\,\psi_{\mathbf k}(x_\alpha)^2
\le\frac{C_K}{N^2}\,|\mathcal I(N,t_n+\theta)|\le\frac{C_K}N ,
\]
and $\tfrac{\mathrm d}{\mathrm d\theta}\langle M_{\mathbf k}^{\mathsf A}\rangle_\theta$ carrying only the
Brownian term; hence
\[
\frac{\mathrm d}{\mathrm d\theta}\langle M_{\mathbf k}\rangle_\theta\le\frac{C_K}{N}.
\]
Since $\delta Z_{\mathbf k}(0)=0$, we have
$\delta Z_{\mathbf k}(\theta)=\int_0^\theta B_{\mathbf k}(r)\,\mathrm dr+M_{\mathbf k}(\theta)$,
so by $(a+b)^2\le2a^2+2b^2$ and Cauchy--Schwarz in $r$, together with $|B_{\mathbf k}|\le C_K$,
\[
|\delta Z_{\mathbf k}(\theta)|^2
\le 2\theta\int_0^\theta B_{\mathbf k}(r)^2\,\mathrm dr+2M_{\mathbf k}(\theta)^2
\le 2C_K^2\theta^2+2M_{\mathbf k}(\theta)^2 .
\]
Taking $\mathbb E_n$ and using the isometry together with the rate bound just obtained,
$\mathbb E_nM_{\mathbf k}(\theta)^2=\mathbb E_n\langle M_{\mathbf k}\rangle_\theta\le C_K\theta/N$,
then summing over the $|\mathcal M_K|\le C_K$ modes,
\begin{equation}\label{eq:L4-mode2}
\mathbb E_n|\delta Z_\theta|_K^2\le C_K\Bigl(\theta^2+\frac{\theta}{N}\Bigr),\qquad 0\le\theta\le\tau .
\end{equation}

\medskip
\noindent\textbf{Step 2: labelwise displacement bound.}
Let
\[
\Delta_i(\theta):=\widehat Y_i(\theta)-Y_i(\theta),\qquad
\Gamma_\theta^2:=\frac1N\sum_{i\in\mathcal I_n}|\Delta_i(\theta)|^2 .
\]
Since the two root trajectories are driven by the same Brownian stream, their difference
satisfies
\[
\frac{\mathrm d}{\mathrm d\theta}\Delta_i(\theta)
=b_{\nu^\Phi_\theta}(\widehat Y_i(\theta))-b_{\nu^{\mathsf A}_\theta}(Y_i(\theta)).
\]
Using \eqref{eq:L4-coeff} and the $C^1$-bound on $b_{\nu^{\mathsf A}_\theta}$,
\[
|\dot\Delta_i(\theta)|\le C_K\bigl(|\delta Z_\theta|_K+|\Delta_i(\theta)|\bigr).
\]
Since $\Delta_i(0)=0$, integrating and applying Cauchy--Schwarz in $r$ gives
\[
|\Delta_i(\theta)|^2
=\Bigl|\int_0^\theta\dot\Delta_i(r)\,\mathrm dr\Bigr|^2
\le\theta\int_0^\theta|\dot\Delta_i(r)|^2\,\mathrm dr
\le 2C_K^2\,\theta\int_0^\theta\bigl(|\delta Z_r|_K^2+|\Delta_i(r)|^2\bigr)\,\mathrm dr,
\]
using $(a+b)^2\le2a^2+2b^2$ in the last step. Averaging over $i\in\mathcal I_n$ and
using that $|\delta Z_r|_K$ is independent of $i$, so
$\frac1N\sum_{i\in\mathcal I_n}|\delta Z_r|_K^2=\frac{|\mathcal I_n|}{N}|\delta Z_r|_K^2\le C_N|\delta Z_r|_K^2$
on the admissible set, with $C_N$ absorbed into $C_K$,
\[
\Gamma_\theta^2\le C_K\theta\int_0^\theta\bigl(|\delta Z_r|_K^2+\Gamma_r^2\bigr)\,\mathrm dr .
\]
Taking conditional expectation and using \eqref{eq:L4-mode2}, followed by Gr\"onwall, gives
\begin{equation}\label{eq:L4-disp2}
\mathbb E_n\Gamma_\theta^2\le C_K\Bigl(\theta^4+\frac{\theta^3}{N}\Bigr).
\end{equation}

\medskip
\noindent\textbf{Step 3: weak generator representation of the feedback.}
For $f\in C^4(\mathbb T^d)$, define the cylinder functional
\[
H_f(\theta):=\bigl\langle f,\ell^{\mathrm{drift}}_\theta\bigr\rangle
=\frac1N\sum_{i\in\mathcal I_n}\bigl[f(\widehat Y_i(\theta))-f(Y_i(\theta))\bigr],
\]
where
\[
\ell^{\mathrm{drift}}_\theta:=\frac1N\sum_{i\in\mathcal I_n}\bigl(\delta_{\widehat Y_i(\theta)}-\delta_{Y_i(\theta)}\bigr).
\]
Applying It\^o's formula to the paired trajectories gives
\begin{equation}\label{eq:L4-ito}
H_f(\theta)=\int_0^\theta A_f(r)\,\mathrm dr+M_f(\theta),
\end{equation}
where
\begin{equation}\label{eq:L4-Af}
\begin{aligned}
A_f(r):=\frac1N\sum_{i\in\mathcal I_n}\Bigl[
&\,b_{\nu^\Phi_r}(\widehat Y_i(r))\cdot\nabla f(\widehat Y_i(r))
-b_{\nu^{\mathsf A}_r}(Y_i(r))\cdot\nabla f(Y_i(r))\\
&\quad+D\Delta f(\widehat Y_i(r))-D\Delta f(Y_i(r))\Bigr],
\end{aligned}
\end{equation}
and
\[
M_f(\theta)=\frac{\sqrt{2D}}{N}\sum_{i\in\mathcal I_n}\int_0^\theta
\bigl[\nabla f(\widehat Y_i(r))-\nabla f(Y_i(r))\bigr]\cdot\mathrm dW_i(t_n+r).
\]
There is no direct jump martingale in $H_f$, because the auxiliary root positions do not
jump when the empirical field $\nu^\Phi$ branches: the branching randomness enters $H_f$
only through the finite-dimensional coordinates $Z^\Phi$.

We bound the $i$-th summand of \eqref{eq:L4-Af} termwise. For the drift part, insert
the intermediate value $b_{\nu^{\mathsf A}_r}(\widehat Y_i)\cdot\nabla f(\widehat Y_i)$:
\[
b_{\nu^\Phi_r}(\widehat Y_i)\!\cdot\!\nabla f(\widehat Y_i)
-b_{\nu^{\mathsf A}_r}(Y_i)\!\cdot\!\nabla f(Y_i)
=\underbrace{\bigl[b_{\nu^\Phi_r}-b_{\nu^{\mathsf A}_r}\bigr](\widehat Y_i)\!\cdot\!\nabla f(\widehat Y_i)}_{(\mathrm I)}
+\underbrace{\bigl[(b_{\nu^{\mathsf A}_r}\!\cdot\!\nabla f)(\widehat Y_i)-(b_{\nu^{\mathsf A}_r}\!\cdot\!\nabla f)(Y_i)\bigr]}_{(\mathrm{II})}.
\]
Term $(\mathrm I)$ is a coefficient difference at the fixed point $\widehat Y_i$, bounded
by~\eqref{eq:L4-coeff} (with $q=0$),
\[
|(\mathrm I)|\le\|b_{\nu^\Phi_r}-b_{\nu^{\mathsf A}_r}\|_{C^0}\,\|\nabla f\|_{C^0}
\le C_K\|f\|_{C^1}|\delta Z_r|_K .
\]
Term $(\mathrm{II})$ is the same function $b_{\nu^{\mathsf A}_r}\!\cdot\!\nabla f$ evaluated at
two points, so by the mean value theorem and the $C^1$-bound on $b_{\nu^{\mathsf A}_r}$,
\[
|(\mathrm{II})|\le\|b_{\nu^{\mathsf A}_r}\!\cdot\!\nabla f\|_{C^1}|\Delta_i|
\le\|b_{\nu^{\mathsf A}_r}\|_{C^1}\|\nabla f\|_{C^1}|\Delta_i|
\le C_K\|f\|_{C^2}|\Delta_i| .
\]
The diffusion part is likewise a single function evaluated at two points,
\[
\bigl|D\Delta f(\widehat Y_i)-D\Delta f(Y_i)\bigr|
\le D\|\nabla\Delta f\|_{C^0}|\Delta_i|\le C_K\|f\|_{C^3}|\Delta_i| ,
\]
with $D$ absorbed into $C_K$. Adding the three bounds, the $i$-th bracket of \eqref{eq:L4-Af} satisfies
\[
\begin{aligned}
&\bigl|b_{\nu^\Phi_r}(\widehat Y_i)\!\cdot\!\nabla f(\widehat Y_i)
-b_{\nu^{\mathsf A}_r}(Y_i)\!\cdot\!\nabla f(Y_i)\\
&\qquad\qquad+D\Delta f(\widehat Y_i)-D\Delta f(Y_i)\bigr|\\
&\quad\le|(\mathrm I)|+|(\mathrm{II})|+\bigl|D\Delta f(\widehat Y_i)-D\Delta f(Y_i)\bigr|\\
&\quad\le C_K\|f\|_{C^1}|\delta Z_r|_K+C_K\|f\|_{C^2}|\Delta_i|+C_K\|f\|_{C^3}|\Delta_i|\\
&\quad\le C_K\|f\|_{C^3}\bigl(|\delta Z_r|_K+|\Delta_i|\bigr),
\end{aligned}
\]
the last step using $\|f\|_{C^1}\le\|f\|_{C^2}\le\|f\|_{C^3}$. Summing over $i$ and applying the
triangle inequality to \eqref{eq:L4-Af},
\[
|A_f(r)|
\le\frac1N\sum_{i\in\mathcal I_n}C_K\|f\|_{C^3}\bigl(|\delta Z_r|_K+|\Delta_i|\bigr)
=C_K\|f\|_{C^3}\Bigl(\frac1N\sum_{i\in\mathcal I_n}|\delta Z_r|_K
+\frac1N\sum_{i\in\mathcal I_n}|\Delta_i|\Bigr).
\]
On the admissible set $\{t_n\le T\wedge\Theta_{K,R,N}\}$ one has $|\mathcal I_n|\le C_NN$, so the
$i$-independent mode sum is
\[
\frac1N\sum_{i\in\mathcal I_n}|\delta Z_r|_K
=\frac{|\mathcal I_n|}N\,|\delta Z_r|_K
\le C_N|\delta Z_r|_K,
\]
while Cauchy--Schwarz together with $\Gamma_r^2=\frac1N\sum_{i\in\mathcal I_n}|\Delta_i|^2$ gives
\[
\frac1N\sum_{i\in\mathcal I_n}|\Delta_i|
\le\frac1N\Bigl(\sum_{i\in\mathcal I_n}1\Bigr)^{1/2}\Bigl(\sum_{i\in\mathcal I_n}|\Delta_i|^2\Bigr)^{1/2}
=\Bigl(\frac{|\mathcal I_n|}N\Bigr)^{1/2}\Bigl(\frac1N\sum_{i\in\mathcal I_n}|\Delta_i|^2\Bigr)^{1/2}
\le C_N^{1/2}\,\Gamma_r.
\]
Absorbing $C_N,C_N^{1/2}$ into $C_K$ and using $\|f\|_{C^3}\le\|f\|_{C^4}$,
\[
|A_f(r)|\le C_K\|f\|_{C^4}\bigl(|\delta Z_r|_K+\Gamma_r\bigr).
\]
Squaring and using $(a+b)^2\le2a^2+2b^2$,
\[
|A_f(r)|^2
\le C_K\|f\|_{C^4}^2\bigl(|\delta Z_r|_K+\Gamma_r\bigr)^2
\le 2C_K\|f\|_{C^4}^2\bigl(|\delta Z_r|_K^2+\Gamma_r^2\bigr),
\]
so, taking $\mathbb E_n$ and inserting \eqref{eq:L4-mode2} and \eqref{eq:L4-disp2},
\[
\mathbb E_n|A_f(r)|^2
\le 2C_K\|f\|_{C^4}^2\bigl(\mathbb E_n|\delta Z_r|_K^2+\mathbb E_n\Gamma_r^2\bigr)
\le C_K\|f\|_{C^4}^2\Bigl(r^2+\frac rN+r^4+\frac{r^3}N\Bigr).
\]
For $0\le r\le\tau\le1$ we have $r^4\le r^2$ and $r^3/N\le r/N$, hence
\begin{equation}\label{eq:L4-Af2}
\mathbb E_n|A_f(r)|^2\le C_K\|f\|_{C^4}^2\Bigl(r^2+\frac rN\Bigr),
\qquad 0\le r\le\tau\le1.
\end{equation}

\medskip
\noindent\textbf{Step 4: generator bound and mean estimate.}
Fix $N$. On $\{t_n\le T\wedge\Theta_{K,R,N}\}$ the active population is $\le C_NN$, so
\[
\Xi_\theta=\Bigl(\{x_\alpha(\theta)\}_{\alpha\in\mathcal I(N,t_n+\theta)},\,
\{\widehat Y_i(\theta),Y_i(\theta)\}_{i\in\mathcal I_n},\,
Z^\Phi(\theta),Z^{\mathsf A}(\theta)\Bigr)
\]
is a finite-dimensional jump--diffusion, and $A_f$ from \eqref{eq:L4-Af} is a $C^2$ cylinder
functional of it with $A_f(0)=0$ (since $\widehat Y_i(0)=Y_i(0)$ and $Z^\Phi(0)=Z^{\mathsf A}(0)$).
Let $\mathcal G^{\mathrm{aux}}$ be the infinitesimal generator of $\Xi_\theta$. We apply it only to
cylinder functionals of the root and mode coordinates --- as $A_f$ is --- that is, to
\[
g\in C^2\bigl((\mathbb T^d\times\mathbb T^d)^{|\mathcal I_n|}\times\mathbb R^{|\mathcal M_K|}\times\mathbb R^{|\mathcal M_K|};\,\mathbb R\bigr),
\]
a function of the root pairs $\{(\widehat Y_i,Y_i)\}_{i\in\mathcal I_n}$ and the two real mode families
$Z^\Phi,Z^{\mathsf A}\in\mathbb R^{|\mathcal M_K|}$. As the root labels $\mathcal I_n$ and the mode set
$\mathcal M_K$ are fixed over the step, this is a fixed Euclidean space of dimension
$2d\,|\mathcal I_n|+2\,|\mathcal M_K|$ (with $|\mathcal M_K|=(2K+1)^d$); the active positions
$\{x_\alpha\}$, whose number varies under branching, are not arguments of $g$. For such $g$ and
$\Xi_\theta=\xi$,
\begin{equation}\label{eq:L4-gendef}
\mathcal G^{\mathrm{aux}}g(\xi)
:=\lim_{h\downarrow0}\frac1h\Bigl(\mathbb E\bigl[g(\Xi_{\theta+h})\mid\Xi_\theta=\xi\bigr]-g(\xi)\Bigr),
\end{equation}
the limit existing because $\Xi_\theta$ is a finite-dimensional jump--diffusion and $g\in C^2$.
We split the generator
\begin{equation}\label{eq:L4-split}
\mathcal G^{\mathrm{aux}}=\mathcal L_Y+\mathcal L_Z+\mathcal L_{YZ}+\mathcal J,
\end{equation}
with
\begin{equation}\label{eq:L4-genops}
\begin{aligned}
\mathcal L_Yg&=\sum_{i\in\mathcal I_n}\Bigl[
b_{Z^\Phi}(\widehat Y_i)\!\cdot\!\partial_{\widehat Y_i}
+b_{Z^{\mathsf A}}(Y_i)\!\cdot\!\partial_{Y_i}
+D\bigl(\Delta_{\widehat Y_i}+\Delta_{Y_i}+2\nabla_{\widehat Y_i}\!\cdot\!\nabla_{Y_i}\bigr)\Bigr]g,\\
\mathcal L_Zg&=\sum_{\mathbf k}\Bigl(b^{Z^\Phi}_{\mathbf k}\,\partial_{Z^\Phi_{\mathbf k}}
+b^{Z^{\mathsf A}}_{\mathbf k}\,\partial_{Z^{\mathsf A}_{\mathbf k}}\Bigr)g
+\tfrac12\sum_{\mathbf k,\ell}\frac{\mathrm d\langle Z^\Phi_{\mathbf k},Z^\Phi_\ell\rangle}{\mathrm d\theta}\,\partial^2_{Z^\Phi_{\mathbf k}Z^\Phi_\ell}g\\
&\quad+\tfrac12\sum_{\mathbf k,\ell}\frac{\mathrm d\langle Z^{\mathsf A}_{\mathbf k},Z^{\mathsf A}_\ell\rangle}{\mathrm d\theta}\,\partial^2_{Z^{\mathsf A}_{\mathbf k}Z^{\mathsf A}_\ell}g
+\sum_{\mathbf k,\ell}\frac{\mathrm d\langle Z^\Phi_{\mathbf k},Z^{\mathsf A}_\ell\rangle}{\mathrm d\theta}\,\partial^2_{Z^\Phi_{\mathbf k}Z^{\mathsf A}_\ell}g,\\
\mathcal L_{YZ}g&=\frac{2D}{N}\sum_{\mathbf k}\sum_{i\in\mathcal I_n}\Bigl[
\nabla\psi_{\mathbf k}(\widehat Y_i)\!\cdot\!\partial_{Z^\Phi_{\mathbf k}}\nabla_{\widehat Y_i}g
+\nabla\psi_{\mathbf k}(\widehat Y_i)\!\cdot\!\partial_{Z^\Phi_{\mathbf k}}\nabla_{Y_i}g\\
&\qquad\qquad+\nabla\psi_{\mathbf k}(Y_i)\!\cdot\!\partial_{Z^{\mathsf A}_{\mathbf k}}\nabla_{\widehat Y_i}g
+\nabla\psi_{\mathbf k}(Y_i)\!\cdot\!\partial_{Z^{\mathsf A}_{\mathbf k}}\nabla_{Y_i}g\Bigr],\\
\mathcal Jg&=\sum_{\alpha}\sum_{\sigma=\pm}c^\sigma_{\nu^\Phi_\theta}(x_\alpha)\,
\bigl[g(Z^\Phi+h^\sigma_\alpha,\cdot)-g(Z^\Phi,\cdot)\bigr],
\end{aligned}
\end{equation}
where $\mathbf k,\ell\in\mathcal M_K$, and in $\mathcal L_{YZ}$ the mode--root
cross-variation coefficients arise from the shared Brownian stream, the $Z^\Phi$-modes
contributing $\tfrac{2D}N\nabla\psi_{\mathbf k}(\widehat Y_i)$ and the $Z^{\mathsf A}$-modes
$\tfrac{2D}N\nabla\psi_{\mathbf k}(Y_i)$. In $\mathcal J$, for a type-$\sigma$ event
($\sigma\in\{+,-\}$) at $x_\alpha$ the vector $h^\sigma_\alpha\in\mathbb R^{|\mathcal M_K|}$
is the resulting increment of the mode vector
$Z^\Phi=(Z^\Phi_{\mathbf k})_{\mathbf k\in\mathcal M_K}$, namely
$(h^\sigma_\alpha)_{\mathbf k}=\tfrac{\sigma}N\psi_{\mathbf k}(x_\alpha)$ (a birth $\sigma=+$
adds the atom $\tfrac1N\delta_{x_\alpha}$, a death $\sigma=-$ removes it), so the branching
jumps move only the $Z^\Phi$ coordinates. 

Since $A_f(0)=0$, Dynkin's formula applied to the cylinder functional $A_f$ gives
\begin{equation}\label{eq:L4-dynkin}
\mathbb E_nA_f(r)=\int_0^r\mathbb E_n\bigl[\mathcal G^{\mathrm{aux}}A_f(q)\bigr]\,\mathrm dq,
\end{equation}
the martingale part of $A_f(\Xi_\theta)$ having zero $\mathbb E_n$-mean; thus a pointwise bound
on the generator controls the conditional mean. We claim
\begin{equation}\label{eq:L4-gen}
|\mathcal G^{\mathrm{aux}}A_f(\theta)|\le C_K\|f\|_{C^4}.
\end{equation}

We now bound the derivatives of $A_f$ entering \eqref{eq:L4-genops}. Recall
$b_z(x)=-a\bigl(P_z(x)\bigr)$ with $P_z(x):=\sum_{\mathbf k\in\mathcal M_K}z_{\mathbf k}\psi_{\mathbf k}(x)$,
so that $b_{Z^\Phi}=b_z|_{z=Z^\Phi}$ and $b_{Z^{\mathsf A}}=b_z|_{z=Z^{\mathsf A}}$. In \eqref{eq:L4-Af}
the $i$-th hatted summand $b_{Z^\Phi}(\widehat Y_i)\!\cdot\!\nabla f(\widehat Y_i)+D\Delta f(\widehat Y_i)$
depends only on $(Z^\Phi,\widehat Y_i)$ and the $i$-th plain summand only on $(Z^{\mathsf A},Y_i)$; no summand
depends on two distinct roots, nor on both $\widehat Y_i$ and $Y_i$. It thus suffices to differentiate the
bracket $b_z(x)\!\cdot\!\nabla f(x)+D\Delta f(x)$ in $x$ and in $z$.

For the position derivatives, by the product rule
$\nabla_x[b_z\!\cdot\!\nabla f+D\Delta f]=(\nabla_x b_z)^{\!\top}\nabla f+(\nabla^2 f)b_z+D\nabla\Delta f$;
differentiating once more and grouping by the order of the $f$-factor, with
$\|b_z\|_{C^0}+\|b_z\|_{C^1}+\|b_z\|_{C^2}\le C_K$ by \eqref{eq:Bb},
\[
\bigl\|\nabla_x[b_z\!\cdot\!\nabla f+D\Delta f]\bigr\|_\infty
\le\|b_z\|_{C^1}\|f\|_{C^1}+\|b_z\|_{C^0}\|f\|_{C^2}+D\|f\|_{C^3}\le C_K\|f\|_{C^3},
\]
\[
\bigl\|\nabla_x^2[b_z\!\cdot\!\nabla f+D\Delta f]\bigr\|_\infty
\le\|b_z\|_{C^2}\|f\|_{C^1}+2\|b_z\|_{C^1}\|f\|_{C^2}+\|b_z\|_{C^0}\|f\|_{C^3}+D\|f\|_{C^4}\le C_K\|f\|_{C^4}.
\]
For the mode derivatives, since $D\Delta f$ is $z$-independent and $\partial_{z_{\mathbf k}}P_z=\psi_{\mathbf k}$,
\[
\partial_{z_{\mathbf k}}b_z=-a'(P_z)\psi_{\mathbf k},\qquad
\partial^2_{z_{\mathbf k}z_\ell}b_z=-a''(P_z)\psi_{\mathbf k}\psi_\ell,
\]
so $\partial_{z_{\mathbf k}}[b_z\!\cdot\!\nabla f+D\Delta f]=(\partial_{z_{\mathbf k}}b_z)\!\cdot\!\nabla f$ and
$\partial^2_{z_{\mathbf k}z_\ell}[b_z\!\cdot\!\nabla f+D\Delta f]=(\partial^2_{z_{\mathbf k}z_\ell}b_z)\!\cdot\!\nabla f$;
by (B1), $\|a'\|_\infty+\|a''\|_\infty\le C$, and $\|\psi_{\mathbf k}\|_{C^2}\le C_K$, whence
$\|\partial_{z_{\mathbf k}}b_z\|_\infty+\|\partial^2_{z_{\mathbf k}z_\ell}b_z\|_\infty\le C_K$ and
\[
\bigl|\partial_{z_{\mathbf k}}[b_z\!\cdot\!\nabla f+D\Delta f]\bigr|
+\bigl|\partial^2_{z_{\mathbf k}z_\ell}[b_z\!\cdot\!\nabla f+D\Delta f]\bigr|\le C_K\|f\|_{C^1}.
\]
For the mixed mode--position derivative, with $\beta_{\mathbf k}:=\partial_{z_{\mathbf k}}b_z$,
$\nabla_x\beta_{\mathbf k}=-a''(P_z)(\nabla P_z)\psi_{\mathbf k}-a'(P_z)\nabla\psi_{\mathbf k}$, and
$\|\beta_{\mathbf k}\|_\infty+\|\nabla_x\beta_{\mathbf k}\|_\infty\le C_K$ using
$\|\nabla P_z\|_\infty\le K\|P_z\|_\infty\le C_K$ at the mode vectors arising along the dynamics
(Bernstein \eqref{equ:Bernstein_inequality} and (B2)); hence
$\partial_{z_{\mathbf k}}\nabla_x[b_z\!\cdot\!\nabla f+D\Delta f]
=(\nabla_x\beta_{\mathbf k})^{\!\top}\nabla f+(\nabla^2 f)\beta_{\mathbf k}$ has sup norm $\le C_K\|f\|_{C^2}$.

We pass to $A_f=\tfrac1N\sum_{i\in\mathcal I_n}[\,\cdots\,]$. A root $\widehat Y_i$ (resp.\ $Y_i$) lies in a
single summand, so $\partial_{\widehat Y_i}A_f=\tfrac1N\nabla_x[\,\cdots\,](\widehat Y_i)$, which gives the
$\tfrac1N$ in the position bounds; as no summand contains both $\widehat Y_i$ and $Y_i$, the off-diagonal
block $\nabla_{\widehat Y_i}\!\cdot\!\nabla_{Y_i}A_f=0$, leaving only the diagonal Laplacian blocks in
$\nabla^2_{(\widehat Y_i,Y_i)}A_f$. A mode $Z_{\mathbf k}\in\{Z^\Phi_{\mathbf k},Z^{\mathsf A}_{\mathbf k}\}$ lies in
all $|\mathcal I_n|$ summands of one sign, so
$\partial_{Z_{\mathbf k}}A_f=\tfrac1N\sum_{i\in\mathcal I_n}\partial_{z_{\mathbf k}}[\,\cdots\,]$ and, by
$|\mathcal I_n|\le C_NN$ (with $C_N$ absorbed into $C_K$),
$|\partial_{Z_{\mathbf k}}A_f|\le\tfrac{|\mathcal I_n|}N C_K\|f\|_{C^1}\le C_K\|f\|_{C^1}$, and likewise for
$\partial^2_{Z_{\mathbf k}Z_\ell}A_f$; the mixed derivative
$\partial_{Z_{\mathbf k}}\nabla_{\widehat Y_i}A_f=\tfrac1N\partial_{z_{\mathbf k}}\nabla_x[\,\cdots\,](\widehat Y_i)$
touches one root and keeps the $\tfrac1N$. Collecting these,
\begin{equation}\label{eq:L4-derivs}
\begin{gathered}
|\partial_{\widehat Y_i}A_f|+|\partial_{Y_i}A_f|\le\tfrac{C_K}N\|f\|_{C^3},\qquad
|\nabla^2_{(\widehat Y_i,Y_i)}A_f|\le\tfrac{C_K}N\|f\|_{C^4},\\[2pt]
|\partial_{Z_{\mathbf k}}A_f|+|\partial^2_{Z_{\mathbf k}Z_\ell}A_f|\le C_K\|f\|_{C^1},\qquad
|\partial_{Z_{\mathbf k}}\nabla_{(\widehat Y_i,Y_i)}A_f|\le\tfrac{C_K}N\|f\|_{C^2}.
\end{gathered}
\end{equation}
The position bounds carry $\tfrac1N$ because each root sits in a single summand; the mode bounds carry
none because the $\tfrac1N$ is cancelled by the $|\mathcal I_n|\le C_NN$ summands sharing each mode.

By \eqref{eq:L4-genops}, $\|b_\nu\|_\infty\le C_K$, and \eqref{eq:L4-derivs}, the root
transport--diffusion is bounded, the shared-noise cross term $\nabla_{\widehat Y_i}\!\cdot\!\nabla_{Y_i}A_f$
vanishing, by
\[
|\mathcal L_YA_f|
\le\sum_{i\in\mathcal I_n}\Bigl[\|b_\nu\|_\infty\bigl(|\partial_{\widehat Y_i}A_f|+|\partial_{Y_i}A_f|\bigr)
+2D\,|\nabla^2_{(\widehat Y_i,Y_i)}A_f|\Bigr]
\le|\mathcal I_n|\,\frac{C_K}N\|f\|_{C^4}\le C_K\|f\|_{C^4}.
\]
Using the coefficient bounds stated with \eqref{eq:L4-genops} together with \eqref{eq:L4-derivs},
the mode drift--diffusion and the mode--root cross-variation satisfy
\[
|\mathcal L_ZA_f|\le\Bigl(|\mathcal M_K|+\frac{|\mathcal M_K|^2}N\Bigr)C_K\|f\|_{C^1}\le C_K\|f\|_{C^1},
\qquad
|\mathcal L_{YZ}A_f|
\le|\mathcal M_K|\,|\mathcal I_n|\,\frac{C_K}{N}\,\frac{C_K}N\|f\|_{C^2}\le C_K\|f\|_{C^2}.
\]
For the branching jumps, a second-order Taylor expansion in $Z^\Phi$ with \eqref{eq:L4-derivs} gives
$|A_f(Z^\Phi+h^\sigma_\alpha,\cdot)-A_f(Z^\Phi,\cdot)|\le C_K\|f\|_{C^1}(\tfrac1N+\tfrac1{N^2})$,
while $\sum_\alpha\sum_{\sigma=\pm}c^\sigma_{\nu^\Phi_\theta}(x_\alpha)\le2C_K|\mathcal I(N,t_n+\theta)|\le C_KN$, so
\[
|\mathcal JA_f|
\le\sum_{\alpha}\sum_{\sigma=\pm}c^\sigma_{\nu^\Phi_\theta}(x_\alpha)\bigl|A_f(Z^\Phi+h^\sigma_\alpha,\cdot)-A_f(Z^\Phi,\cdot)\bigr|
\le C_KN\cdot C_K\|f\|_{C^1}\Bigl(\frac1N+\frac1{N^2}\Bigr)\le C_K\|f\|_{C^1}.
\]
The split intensities $c^\pm$ enter only as bounded rates and are never differentiated. Summing the
four bounds proves \eqref{eq:L4-gen}.

By \eqref{eq:L4-dynkin}, the bound \eqref{eq:L4-gen} gives $|\mathbb E_nA_f(r)|\le C_Kr\|f\|_{C^4}$.
Using \eqref{eq:L4-ito}
and $\mathbb E_nM_f(\tau)=0$,
\[
\bigl|\mathbb E_n\langle f,\ell^{\mathrm{drift}}_{n+1}\rangle\bigr|
\le\int_0^\tau|\mathbb E_nA_f(r)|\,\mathrm dr\le C_K\tau^2\|f\|_{C^4}.
\]
Since $s>d/2+4$, $H^s(\mathbb T^d)\hookrightarrow C^4(\mathbb T^d)$ by (B4); taking the
supremum over $\|f\|_{H^s}\le1$,
\[
\|\mathbb E_n\ell^{\mathrm{drift}}_{n+1}\|_{H^{-s}}\le C_K\tau^2.
\]

\medskip
\noindent\textbf{Step 5: centered second moment.}
From \eqref{eq:L4-ito},
\begin{equation}\label{eq:L4-center}
H_f(\tau)-\mathbb E_nH_f(\tau)
=\int_0^\tau\bigl(A_f(r)-\mathbb E_nA_f(r)\bigr)\,\mathrm dr+M_f(\tau).
\end{equation}
Using Cauchy--Schwarz in time and \eqref{eq:L4-Af2},
\begin{equation}\label{eq:L4-driftvar}
\begin{aligned}
\mathbb E_n\Bigl|\int_0^\tau\bigl(A_f(r)-\mathbb E_nA_f(r)\bigr)\,\mathrm dr\Bigr|^2
&\le\tau\int_0^\tau\mathbb E_n|A_f(r)-\mathbb E_nA_f(r)|^2\,\mathrm dr\\
&\le4\tau\int_0^\tau\mathbb E_n|A_f(r)|^2\,\mathrm dr\\
&\le C_K\|f\|_{C^4}^2\Bigl(\tau^4+\frac{\tau^3}{N}\Bigr).
\end{aligned}
\end{equation}
For the Brownian martingale, It\^o's isometry and \eqref{eq:L4-disp2} yield
\begin{equation}\label{eq:L4-Mvar}
\begin{aligned}
\mathbb E_n|M_f(\tau)|^2
&=\frac{2D}{N^2}\sum_{i\in\mathcal I_n}\int_0^\tau\mathbb E_n\bigl|\nabla f(\widehat Y_i(r))-\nabla f(Y_i(r))\bigr|^2\,\mathrm dr\\
&\le\frac{C\|f\|_{C^2}^2}{N}\int_0^\tau\mathbb E_n\Gamma_r^2\,\mathrm dr\\
&\le C_K\|f\|_{C^2}^2\Bigl(\frac{\tau^5}{N}+\frac{\tau^4}{N^2}\Bigr)
\le C_K\|f\|_{C^2}^2\,\frac{\tau^3}{N},\qquad 0<\tau\le1,\ N\ge1 .
\end{aligned}
\end{equation}
Combining \eqref{eq:L4-center}--\eqref{eq:L4-Mvar},
\[
\mathbb E_n\bigl|\langle f,\ell^{\mathrm{drift}}_{n+1}-\mathbb E_n\ell^{\mathrm{drift}}_{n+1}\rangle\bigr|^2
\le C_K\|f\|_{C^4}^2\Bigl(\tau^4+\frac{\tau^3}{N}\Bigr).
\]

Finally, let $\{e_{\mathbf k}\}_{\mathbf k\in\mathbb Z^d}$ be the complex Fourier basis.
Since $\|e_{\mathbf k}\|_{C^4}^2\le C(1+|\mathbf k|^2)^4$,
\[
\begin{aligned}
\mathbb E_n\bigl\|\ell^{\mathrm{drift}}_{n+1}-\mathbb E_n\ell^{\mathrm{drift}}_{n+1}\bigr\|_{H^{-s}}^2
&=\sum_{\mathbf k\in\mathbb Z^d}(1+|\mathbf k|^2)^{-s}\,
\mathbb E_n\bigl|\langle e_{\mathbf k},\ell^{\mathrm{drift}}_{n+1}-\mathbb E_n\ell^{\mathrm{drift}}_{n+1}\rangle\bigr|^2\\
&\le C_K\Bigl(\tau^4+\frac{\tau^3}{N}\Bigr)\sum_{\mathbf k\in\mathbb Z^d}(1+|\mathbf k|^2)^{-s+4}.
\end{aligned}
\]
The last sum is finite because $s>d/2+4$ by (B4). Therefore
\[
\mathbb E_n\bigl\|\ell^{\mathrm{drift}}_{n+1}-\mathbb E_n\ell^{\mathrm{drift}}_{n+1}\bigr\|_{H^{-s}}^2
\le C_K\Bigl(\tau^4+\frac{\tau^3}{N}\Bigr).
\]
Since $0<\tau\le1$, this implies the weaker but more convenient form
\[
\mathbb E_n\bigl\|\ell^{\mathrm{drift}}_{n+1}-\mathbb E_n\ell^{\mathrm{drift}}_{n+1}\bigr\|_{H^{-s}}^2
\le C_K\Bigl(\tau^4+\frac{\tau^2}{N}\Bigr).
\]
This proves the lemma.
\end{proof}

\begin{lemma}[Lie-splitting defect]\label{lem:split}
Under Assumption~\ref{ass:time}, couple the full simultaneous flow
$\Phi^{N,\omega}_\tau$ and the exact Lie-split flow
$\widetilde\Phi^{N,\omega}_\tau$ by a Harris graphical construction: every
genealogical label carries a common Brownian stream and a common Poisson random
measure, with $c_+:=\max(c,0)$, $c_-:=\max(-c,0)$ as birth and death intensities.
Then the Lie-splitting defect $\ell^{\mathrm{split}}_{n+1}
=\Phi^{N,\omega}_\tau(\mu^N_n)-\widetilde\Phi^{N,\omega}_\tau(\mu^N_n)$ satisfies
\begin{equation}\label{eq:split}
\|\mathbb E_n\ell^{\mathrm{split}}_{n+1}\|_{H^{-s}}\le C_K\,\tau^2,
\qquad
\mathbb E_n\|\ell^{\mathrm{split}}_{n+1}-\mathbb E_n\ell^{\mathrm{split}}_{n+1}\|_{H^{-s}}^2
\le C_K\Bigl(\frac{\tau^{3/2}}{N}+\tau^4\Bigr),
\quad 0<\tau\le1.
\end{equation}
\end{lemma}

\begin{proof}
Define the weak generators on measures
\[
\langle f,\mathsf L^{\mathsf A}\nu\rangle:=\langle b_\nu\!\cdot\!\nabla f+D\Delta f,\nu\rangle,
\qquad
\langle f,\mathsf L^{\mathsf B}\nu\rangle:=\langle c_\nu f,\nu\rangle,
\]
so that the full simultaneous flow $\Phi^{N,\omega}_\tau$ has generator
$\mathsf L^{\mathsf A}+\mathsf L^{\mathsf B}$ and the exact Lie-split flow
$\widetilde\Phi^{N,\omega}_\tau=\mathsf B^{N,\omega}_\tau\circ\mathsf A^{N,\omega}_\tau$
applies $\mathsf L^{\mathsf A}$ then $\mathsf L^{\mathsf B}$ in succession. Write
$\mathsf L:=\mathsf L^{\mathsf A}+\mathsf L^{\mathsf B}$ for brevity, and let
$\nu^\Phi_\theta$ ($0\le\theta\le\tau$) denote the empirical measure of
$\Phi^{N,\omega}_\theta(\mu^N_n)$.

\emph{Full simultaneous flow expansion.} For $f\in C^4$ define the drift
functional
\[
F_f(\nu):=\langle b_\nu\!\cdot\!\nabla f+c_\nu f+D\Delta f,\,\nu\rangle,
\]
which depends on $\nu$ both linearly (through the pairing) and nonlinearly
(through $b_\nu,c_\nu$, which involve the finitely many Fourier modes
$|\mathbf k|_\infty\le K$ of $\nu$). By the Dynkin formula for the empirical
measure process,
\[
\mathbb E_nF_f(\nu^\Phi_\theta)-F_f(\mu^N_n)
=\int_0^\theta\mathbb E_n[\mathcal G^N F_f(\nu^\Phi_r)]\,\mathrm dr,
\]
where $\mathcal G^N$ is the generator of $\Phi^{N,\omega}$ on functionals of the
empirical measure. Computing $\mathcal G^N F_f$, the generator acts both on the
linear $f$-pairing (producing terms up to $\Delta^2 f$) and on the finite-mode
coefficient dependence of $b_\nu,c_\nu$ (smooth in the modes $|\mathbf k|_\infty
\le K$ by (B1)); each coefficient factor and its first two derivatives is
$\le C_K$ by \eqref{eq:Bb}, and the It\^o quadratic-variation correction is of
order $1/N$. Hence
\begin{equation}\label{eq:GF-bound}
|\mathcal G^N F_f(\nu)|\le C_K\|f\|_{C^4},
\qquad\text{so}\qquad
|\mathbb E_nF_f(\nu^\Phi_\theta)-F_f(\mu^N_n)|\le C_K\theta\|f\|_{C^4}.
\end{equation}
The weak formulation of $\Phi^{N,\omega}$ reads
$\mathbb E_n\langle f,\Phi^{N,\omega}_\tau(\mu^N_n)\rangle-\langle f,\mu^N_n\rangle
=\int_0^\tau\mathbb E_nF_f(\nu^\Phi_\theta)\,\mathrm d\theta$, so subtracting the
left-endpoint value $\tau F_f(\mu^N_n)$ and using \eqref{eq:GF-bound},
\begin{equation}\label{eq:Phi-exp}
\mathbb E_n\langle f,\Phi^{N,\omega}_\tau(\mu^N_n)\rangle
=\langle f,\mu^N_n\rangle+\tau F_f(\mu^N_n)
+O\bigl(C_K\tau^2\|f\|_{C^4}\bigr),
\end{equation}
the remainder being $\int_0^\tau|\mathbb E_nF_f(\nu^\Phi_\theta)-F_f(\mu^N_n)|\mathrm d\theta
\le\int_0^\tau C_K\theta\|f\|_{C^4}\mathrm d\theta=\tfrac12C_K\tau^2\|f\|_{C^4}$.
Here $F_f(\mu^N_n)=\langle f,\mathsf L\mu^N_n\rangle$ with
$\mathsf L:=\mathsf L^{\mathsf A}+\mathsf L^{\mathsf B}$.

\emph{Lie-split flow expansion.} The same argument applied to $\mathsf A^{N,\omega}_\tau$
(generator $\mathsf L^{\mathsf A}$) gives
\[
\mathbb E_n\langle f,\mathsf A^{N,\omega}_\tau(\mu^N_n)\rangle
=\langle f,\mu^N_n\rangle+\tau\,\langle f,\mathsf L^{\mathsf A}\mu^N_n\rangle
+O\bigl(C_K\tau^2\|f\|_{C^4}\bigr).
\]
Applying $\mathsf B^{N,\omega}_\tau$ to the output and using the same expansion
for $\mathsf L^{\mathsf B}$ in turn,
\begin{equation}\label{eq:splitflow-exp}
\mathbb E_n\langle f,\widetilde\Phi^{N,\omega}_\tau(\mu^N_n)\rangle
=\langle f,\mu^N_n\rangle+\tau\,\langle f,\mathsf L^{\mathsf A}\mu^N_n\rangle
+\tau\,\langle f,\mathsf L^{\mathsf B}\mu^N_n\rangle
+O\bigl(C_K\tau^2\|f\|_{C^4}\bigr),
\end{equation}
the cross term $\tau^2\langle f,\mathsf L^{\mathsf B}\mathsf L^{\mathsf A}\mu^N_n\rangle$
being absorbed into the $O(C_K\tau^2\|f\|_{C^4})$ remainder.

Subtracting \eqref{eq:splitflow-exp} from \eqref{eq:Phi-exp}, the first-order
terms in $\tau$ agree (both equal $\tau\langle f,\mathsf L\mu^N_n\rangle$), so
$|\mathbb E_n\langle f,\ell^{\mathrm{split}}_{n+1}\rangle|\le C_K\tau^2\|f\|_{C^4}$;
by $H^s\hookrightarrow C^4$ and the supremum over $\|f\|_{H^s}\le1$,
$\|\mathbb E_n\ell^{\mathrm{split}}_{n+1}\|_{H^{-s}}\le C_K\tau^2$.

\emph{Pathwise representation and exact decomposition.} We make the decomposition
precise as an identity of stochastic integrals on the graphical construction, so the
four pieces below are explicit functionals of the driving noise and \emph{partition}
$\ell^{\mathrm{split}}_{n+1}$ without overlap. Write
$\widetilde{\mathcal N}^{\sigma}_\alpha(\mathrm d\theta,\mathrm du):=\mathcal N^{\sigma}_\alpha(\mathrm d\theta,\mathrm du)-\mathrm d\theta\,\mathrm du$
($\sigma\in\{+,-\}$) for the compensated Poisson random measures, $X^\Phi_\alpha(\theta)$
for the positions in the simultaneous flow and $X^{\mathsf A}_i(\theta)$ for those in
the transport substep, and $\nu^\Phi_\theta,\nu^{\mathsf A}_\theta,\eta_\theta$ for the
empirical measures of $\Phi^{N,\omega}$, of $\mathsf A^{N,\omega}$, and of the
pure-reaction substep $\mathsf B^{N,\omega}$ started from $\nu^{\mathsf A}_\tau$, so
that $\widetilde\Phi^{N,\omega}_\tau(\mu^N_n)=\eta_\tau$. For $f\in C^2(\mathbb T^d)$,
It\^o's formula for these branching diffusions gives the pathwise weak identities
\begin{align}
\langle f,\nu^\Phi_\tau\rangle
&=\langle f,\mu^N_n\rangle
+\int_0^\tau\!\bigl\langle (b_{\nu^\Phi_\theta}\!\cdot\!\nabla+D\Delta+c_{\nu^\Phi_\theta})f,\nu^\Phi_\theta\bigr\rangle\mathrm d\theta
+\Sigma^\Phi_f+J^\Phi_f,\label{eq:Phi-path}\\
\langle f,\eta_\tau\rangle
&=\langle f,\mu^N_n\rangle
+\int_0^\tau\!\bigl\langle (b_{\nu^{\mathsf A}_\theta}\!\cdot\!\nabla+D\Delta)f,\nu^{\mathsf A}_\theta\bigr\rangle\mathrm d\theta
+\int_0^\tau\!\bigl\langle c_{\eta_\theta}f,\eta_\theta\bigr\rangle\mathrm d\theta
+\Sigma^{\mathsf A}_f+J^{\mathsf B}_f,\label{eq:split-path}
\end{align}
where $\Sigma$ are the Brownian and $J$ the compensated-jump martingales,
\[
\Sigma^\Phi_f=\tfrac1N\!\sum_{\alpha}\!\int_0^\tau\!\!\sqrt{2D}\,\nabla f(X^\Phi_\alpha)\!\cdot\!\mathrm d\mathbf W_\alpha,
\quad
J^\Phi_f=\tfrac1N\!\sum_{\alpha,\sigma}\!\sigma\!\int_0^\tau\!\!\!\int_0^\infty\!\! f(X^\Phi_\alpha(\theta^-))\mathbf 1_{\{u\le c^{\sigma}_{\nu^\Phi_\theta}(X^\Phi_\alpha(\theta^-))\}}\widetilde{\mathcal N}^{\sigma}_\alpha,
\]
and $\Sigma^{\mathsf A}_f,J^{\mathsf B}_f$ are the analogous sums over the root labels
$i\in\mathcal I_n$ with positions $X^{\mathsf A}_i(\theta)$ (resp.\ frozen at
$X^{\mathsf A}_i(\tau)$ in $J^{\mathsf B}$) and intensities $c^{\sigma}_{\eta_\theta}$.
Subtracting \eqref{eq:split-path} from \eqref{eq:Phi-path},
$\langle f,\ell^{\mathrm{split}}_{n+1}\rangle$ is the sum of a drift difference (the
$\mathrm d\theta$-integrals), a Brownian-martingale difference
$\Sigma^\Phi_f-\Sigma^{\mathsf A}_f$, and a compensated-jump difference
$J^\Phi_f-J^{\mathsf B}_f$. Along the shared clocks the points of $\mathcal N^\sigma_\alpha$
split disjointly into \emph{matched}
($u\le\min(c^\sigma_{\nu^\Phi_\theta},c^\sigma_{\eta_\theta})$) and \emph{surplus}
($\min<u\le\max$) points; with this we define, as measure-valued random variables,
\begin{itemize}
\item $\ell^{\mathrm{drift}}$: the displacement of the \emph{root} particles caused
by the two flows transporting them under the different running drifts
$b_{\nu^\Phi}$ and $b_{\nu^{\mathsf A}}$ (same roots, same Brownian streams, drift
only)—equivalently the transport drift difference together with the root part of
$\Sigma^\Phi_f-\Sigma^{\mathsf A}_f$. It is the reaction-to-drift feedback, treated
by the finite-dimensional generator estimate of
Lemma~\ref{lem:drift-feedback-generator};
\item $M_{\mathrm{disp}}:=\tfrac1N\sum_{i,\sigma}\sigma\!\int_0^\tau\!\!\int_0^\infty\!
\bigl[f(X^\Phi_i(\theta^-))-f(X^{\mathsf A}_i(\tau))\bigr]\mathbf 1_{\{u\le\min(c^\sigma_{\nu^\Phi_\theta}(X^\Phi_i(\theta^-)),\,c^\sigma_{\eta_\theta}(X^{\mathsf A}_i(\tau)))\}}\widetilde{\mathcal N}^\sigma_i$,
the matched-point part of $J^\Phi_f-J^{\mathsf B}_f$: events firing in both flows but
at the displaced positions $X^\Phi_i(\theta^-)$ versus $X^{\mathsf A}_i(\tau)$;
\item $M_{\mathrm{surplus}}$, the surplus-point part of $J^\Phi_f-J^{\mathsf B}_f$
(the compensated integral over $\{\min<u\le\max\}$), present in exactly one flow;
\item $R_{\mathrm{rem}}:=\ell^{\mathrm{split}}_{n+1}-\ell^{\mathrm{drift}}-M_{\mathrm{disp}}-M_{\mathrm{surplus}}$,
defined by exact subtraction. It collects all terms not included in the root-level
matched/surplus single-event contributions: the post-birth transport of descendant
labels in $\Phi$ over the remaining part of the step --- \emph{including both the drift
displacement and the Brownian motion} $\Sigma^\Phi_f|_{\text{descendants}}$, absent in
the motionless split branching --- and all labels carrying two or more branching events
on $[0,\tau]$.
\end{itemize}
The four terms sum to $\ell^{\mathrm{split}}_{n+1}$ by definition; the matched and
surplus point sets are disjoint, and (since $M_{\mathrm{disp}},M_{\mathrm{surplus}}$
carry only the single-event/root jumps) they do not overlap the multiple-event,
descendant remainder $R_{\mathrm{rem}}$, so there is no double counting. The drift--jump
interactions are not separated into a Fr\'echet response remainder. They are
included in the finite-dimensional cylinder functional $A_f$ in
Lemma~\ref{lem:drift-feedback-generator} and are controlled by Dynkin's formula
and the associated carr\'e-du-champ bounds. Only the centered part of each term enters the
fluctuation bound; $M_{\mathrm{disp}},M_{\mathrm{surplus}}$ are compensated (mean
zero), while $\mathbb E_n\ell^{\mathrm{drift}}$ and $\mathbb E_nR_{\mathrm{rem}}$ are
deterministic biases already counted in
$\|\mathbb E_n\ell^{\mathrm{split}}_{n+1}\|_{H^{-s}}\le C_K\tau^2$. We bound each
centered piece in $H^{-s}$.

\emph{Reaction-to-drift feedback $\ell^{\mathrm{drift}}$.}
In the simultaneous flow, branching events during the step change the running
empirical field $\nu^\Phi_\theta$, and hence the drift $b_{\nu^\Phi_\theta}$
experienced by the root particles. In the split transport substep, the root
particles instead see the branching-free field $\nu^{\mathsf A}_\theta$. We
compare these two root transports with the same Brownian streams and define
$\ell^{\mathrm{drift}}$ as in Lemma~\ref{lem:drift-feedback-generator}. That
lemma gives directly
\[
\|\mathbb E_n\ell^{\mathrm{drift}}\|_{H^{-s}}\le C_K\tau^2,
\]
and
\[
\mathbb E_n\|\ell^{\mathrm{drift}}-\mathbb E_n\ell^{\mathrm{drift}}\|_{H^{-s}}^2
\le C_K\Bigl(\frac{\tau^3}{N}+\tau^4\Bigr)
\le C_K\Bigl(\frac{\tau^2}{N}+\tau^4\Bigr).
\]
The proof uses only the finite-dimensional Fourier-mode dependence of
$b_\nu=-a(\mathcal P_K\nu)$, Dynkin's formula for cylinder functionals,
and martingale quadratic-variation estimates. In particular no Fr\'echet
differentiability of the transport flow on $H^{-s}$, and no differentiability of
the hard split $c^\pm$, is invoked. The common-field correlation across root
particles --- a single branching event perturbs the drift of \emph{all} root
particles simultaneously --- is precisely what the centered second-moment step of
Lemma~\ref{lem:drift-feedback-generator} controls, without any appeal to
per-particle independence.

\emph{Matched events with displaced positions $M_{\mathrm{disp}}$.} An event
common to both flows fires at the same clock time $\theta$, but the particle sits
at $Y_i(\tau;\theta)$ (simultaneous, transported the full step) versus $X_i(\tau)$
(split, transported before branching); these differ by a transport displacement
over time $\le\tau$. Writing the compensated jump martingale $M_{\mathrm{disp}}$
and using $\|\delta_x-\delta_y\|_{H^{-s}}^2\le C_s|x-y|^2$ with the predictable
event intensity $\le C_K$,
\[
\mathbb E_n\|M_{\mathrm{disp}}\|_{H^{-s}}^2
\le\frac{C_s}{N^2}\sum_{i\in\mathcal I_n}\int_0^\tau
\mathbb E_n\|\delta_{Y_i(\tau;\theta)}-\delta_{X_i(\tau)}\|_{H^{-s}}^2\,\mathrm d\theta
\le C_K\,\frac{\tau^2}{N},
\]
since each integrand is $\le C_s\,\mathbb E_n|Y_i(\tau;\theta)-X_i(\tau)|^2\le C_K\tau$
and the time integral contributes one more power of $\tau$.

\emph{Surplus events from intensity mismatch $M_{\mathrm{surplus}}$.} Following the
graphical construction we treat birth and death separately; for $\sigma\in\{+,-\}$
the cumulative $\sigma$-intensities seen by particle $i$ are
\[
\Lambda^{\mathrm{full}}_{i,\sigma}=\int_0^\tau c^{\sigma}_{\nu^\Phi_\theta}(X_i(\theta))\,\mathrm d\theta,
\qquad
\Lambda^{\mathrm{split}}_{i,\sigma}=\int_0^\tau c^{\sigma}_{\eta_\theta}(X_i(\tau))\,\mathrm d\theta,
\qquad c^{+}=\max(c,0),\ \ c^{-}=\max(-c,0).
\]
Since $x\mapsto x_\pm$ is $1$-Lipschitz,
$|c^{\sigma}_{\mu}(x)-c^{\sigma}_{\nu}(y)|\le|c_{\mu}(x)-c_{\nu}(y)|$, so it suffices
to bound the signed difference.
Their difference has a measure part and a spatial part: since
$c_\nu(x)=r(\mathcal P_K\nu(x))$,
\[
|c_{\nu^\Phi_\theta}(X_i(\theta))-c_{\eta_\theta}(X_i(\tau))|
\le C_K\|\nu^\Phi_\theta-\eta_\theta\|_{H^{-s}}+C_K|X_i(\theta)-X_i(\tau)|,
\]
by \eqref{eq:Blip} with $q=0$ and $\|\nabla c\|_{L^\infty}\le C_K$. The measure
part $\|\nu^\Phi_\theta-\eta_\theta\|_{H^{-s}}$ is split by the triangle
inequality into
\[
\|\nu^\Phi_\theta-\eta_\theta\|_{H^{-s}}
\le\|\nu^\Phi_\theta-\nu^{\mathsf A}_\theta\|_{H^{-s}}
+\|\nu^{\mathsf A}_\theta-\nu^{\mathsf A}_\tau\|_{H^{-s}}
+\|\nu^{\mathsf A}_\tau-\eta_\theta\|_{H^{-s}},
\]
where $\nu^{\mathsf A}_\theta$ is the branching-free transport flow and
$\eta_\theta$ the pure reaction flow started from $\nu^{\mathsf A}_\tau$. The
first and third terms are reaction mismatches, separating at rate $O(C_K)$ in
conditional mean and contributing $\le C_K\theta$, hence $\le C_K\tau^2$ after
integration. The middle term is a transport time-shift: in fixed-$K$ $H^{-s}$ it
is controlled by the Brownian displacement of the transported particles,
$\mathbb E_n\|\nu^{\mathsf A}_\theta-\nu^{\mathsf A}_\tau\|_{H^{-s}}\le
\frac{C_s}{N}\sum_j\mathbb E_n|X_j(\theta)-X_j(\tau)|\le C_K\sqrt{\tau-\theta}$,
integrating to $C_K\tau^{3/2}$. The spatial part of the rate difference uses the
same Brownian-scale bound $\mathbb E_n|X_i(\theta)-X_i(\tau)|\le C\sqrt{\tau-\theta}$,
integrating to $C_K\int_0^\tau\sqrt{\tau-\theta}\,\mathrm d\theta=C_K\tau^{3/2}$.
Hence the signed intensity-length difference has
$\mathbb E_n|\Lambda^{\mathrm{full}}_i-\Lambda^{\mathrm{split}}_i|\le C_K\tau^{3/2}$,
the transport time-shift and spatial parts dominating, and the same holds for each
$\sigma$-component $L_{i,\sigma}:=|\Lambda^{\mathrm{full}}_{i,\sigma}-\Lambda^{\mathrm{split}}_{i,\sigma}|$,
so $\mathbb E_n L_{i,\sigma}\le C_K\tau^{3/2}$. We bound $M_{\mathrm{surplus}}$ through the
compensated Poisson random measure (PRM) isometry rather than through any independent
Cox-count statement, since the predictable surplus intensities depend on the same
empirical field $c_{\mu_\theta}$ and are correlated across labels. The surplus part of
$J^\Phi_f-J^{\mathsf B}_f$ is the compensated integral
\[
M_{\mathrm{surplus}}
=\frac1N\sum_{i,\sigma}\sigma\!\int_{[0,\tau]\times[0,\infty)}
\delta_{X_i}\,\mathbf 1_{\{\min<u\le\max\}}\,\widetilde{\mathcal N}^\sigma_i(\mathrm d\theta,\mathrm du),
\]
a sum of stochastic integrals against the \emph{independent} PRMs
$\{\widetilde{\mathcal N}^\sigma_i\}$ of the graphical construction. By the It\^o
isometry for compensated PRMs (orthogonal across disjoint label/sign families,
irrespective of intensity correlations) and $\|\delta_x\|_{H^{-s}}\le C_s$,
\[
\mathbb E_n\|M_{\mathrm{surplus}}\|_{H^{-s}}^2
\le\frac{C_s}{N^2}\sum_{i,\sigma}
\mathbb E_n\!\int_0^\tau\!\!\bigl|c^\sigma_{\nu^\Phi_\theta}(X_i(\theta))
-c^\sigma_{\eta_\theta}(X_i(\tau))\bigr|\,\mathrm d\theta,
\]
the inner integral being precisely the predictable quadratic-variation density of the
$(i,\sigma)$ surplus channel, i.e.\ the expected directing length
$\mathbb E_n L_{i,\sigma}$. (Equivalently, conditional on the predictable intensities
the surplus count is a Cox count with mean $\mathbb E_n L_{i,\sigma}$; we use this only
to identify the compensator, not as an independence hypothesis.) Substituting
$\mathbb E_n L_{i,\sigma}\le C_K\tau^{3/2}$ and the population bound,
\[
\mathbb E_n\|M_{\mathrm{surplus}}\|_{H^{-s}}^2
\le\frac{C_s}{N^2}\sum_{i\in\mathcal I_n}\sum_{\sigma\in\{+,-\}}\mathbb E_n L_{i,\sigma}
\le C_K\,\frac{\tau^{3/2}}{N}.
\]

\emph{Multiple-event and descendant remainder $R_{\mathrm{rem}}$.} A given genealogical
label produces two or more events (a second-generation birth, or a birth followed by a
death) within one step with predictable probability $O_K(\tau^2)$, since each event has
predictable intensity $\le C_K$. We do \emph{not} condition on an endogenous field path
to claim cross-label independence of the counts; instead we write the multiple-event
jump part as $R^{\mathrm{jump}}_{\mathrm{rem}}=\mathcal M+A$, where $\mathcal M$ is a sum
of compensated stochastic integrals over the disjoint label/sign PRM families
$\{\widetilde{\mathcal N}^\sigma_\alpha\}$ of the graphical construction (restricted to
the second-and-later events on each label) and $A$ is the predictable compensator. For
the compensated part, the It\^o isometry for compensated PRMs --- orthogonal across
disjoint label/sign families, irrespective of intensity correlations --- gives
\[
\mathbb E_n\|\mathcal M\|_{H^{-s}}^2
=\frac1{N^2}\sum_{\alpha}\mathbb E_n\!\int_0^\tau\!\|\delta_{x_\alpha}\|_{H^{-s}}^2\,\lambda^{(2)}_\alpha(\theta)\,\mathrm d\theta
\le\frac{C_s}{N^2}\,C_NN\,C_K\tau^2=C_K\frac{\tau^2}{N},
\]
where $\lambda^{(2)}_\alpha$ is the predictable second-and-later-event intensity, of
expected integral $O_K(\tau^2)$ per label (a prior event of probability $\le C_K\tau$
followed by rate $\le C_K$), and $\|\delta_{x_\alpha}\|_{H^{-s}}\le C_s$ for $s>d/2$. The
compensator $A$ is predictable and field-dependent; its conditional mean
$\mathbb E_nA=\mathbb E_nR^{\mathrm{jump}}_{\mathrm{rem}}$ is the deterministic bias of
size $O_K(\tau^2)$ in $H^{-s}$ (counted in the mean bound), and its centered part is
bounded crudely by its second moment,
$\mathbb E_n\|A-\mathbb E_nA\|_{H^{-s}}^2\le\mathbb E_n\|A\|_{H^{-s}}^2\le C_K\tau^4$, the
magnitude $\|A\|_{H^{-s}}\le\frac{C_K}{N}\,C_NN\,C_K\tau^2=C_K\tau^2$ following from the
$\le C_NN$ active labels and the $O_K(\tau^2)$ per-label multiple-event probability. Thus
the standalone $\tau^4$ enters through the dependence of the predictable intensities on
the common field, \emph{not} through any conditional-independence statement. Adding the
two parts,
\[
\mathbb E_n\|R^{\mathrm{jump}}_{\mathrm{rem}}-\mathbb E_nR^{\mathrm{jump}}_{\mathrm{rem}}\|_{H^{-s}}^2\le C_K\Bigl(\frac{\tau^2}{N}+\tau^4\Bigr),
\]
where $R^{\mathrm{jump}}_{\mathrm{rem}}$ is the multiple-event jump part. The
remaining contribution to $R_{\mathrm{rem}}$ is the post-birth transport of descendant
labels in $\Phi$ (absent in the motionless split branching $\mathsf B$), comprising both
a drift displacement and the extra Brownian motion $\Sigma^\Phi_f|_{\mathrm{descendants}}$:
a label produces at most one descendant on $[0,\tau]$ with predictable probability
$\le C_K\tau$, and that descendant then moves for time $\le\tau$. The Brownian part
contributes a martingale of quadratic variation
\[
\le\frac{2DC_s'}{N^2}\sum_{i\in\mathcal I_n}\mathbb P_n(\text{label }i\text{ births})\,\mathbb E_n[\text{descendant lifetime}]
\le\frac{2DC_s'}{N^2}\,C_NN\,C_K\tau\cdot\tau=C_K\frac{\tau^2}{N},
\]
where $C_s'$ is as in Proposition~\ref{prop:u_K_mu}. The drift
part displaces each descendant by $\le C_K\tau$, hence by an $H^{-s}$-atom of size
$\le C_sC_K\tau$ (using $\|\delta_x-\delta_y\|_{H^{-s}}\le C_s|x-y|$ for $s>d/2+1$). It is
likewise a sum of compensated integrals over the independent birth PRM families plus a
predictable compensator. The compensated part has, by the PRM isometry,
\[
\frac{C_s^2C_K^2\,\tau^2}{N^2}\sum_{i\in\mathcal I_n}\mathbb P_n(\text{label }i\text{ births})
\le\frac{C_s^2C_K^2\,\tau^2}{N^2}\,C_NN\,C_K\tau=C_K\frac{\tau^3}{N}\le C_K\frac{\tau^2}{N},
\]
while the predictable compensator is the deterministic $O_K(\tau^2)$ bias (counted in the
mean bound) together with a field-dependent fluctuation bounded crudely by its second
moment $\le C_K\tau^4$. Hence the drift part contributes $\le C_K(\tau^2/N+\tau^4)$, again
with the standalone $\tau^4$ coming from the field-dependence of the predictable
intensities rather than from any conditional-independence statement. All contributions to
$R_{\mathrm{rem}}$ are therefore $\le C_K(\tau^2/N+\tau^4)$; equivalently
\[
\|\mathbb E_nR_{\mathrm{rem}}\|_{H^{-s}}\le C_K\tau^2,\qquad
\mathbb E_n\|R_{\mathrm{rem}}-\mathbb E_nR_{\mathrm{rem}}\|_{H^{-s}}^2\le C_K\Bigl(\frac{\tau^2}{N}+\tau^4\Bigr),
\]
the standalone $\tau^4$ being the across-field multiplicity and descendant-transport
variance.

Combining the four centered contributions via $\|\sum\|^2\le4\sum\|\cdot\|^2$
($M_{\mathrm{disp}}$, $M_{\mathrm{surplus}}$, the drift feedback $\ell^{\mathrm{drift}}$,
and the descendant/multiple-event remainder $R_{\mathrm{rem}}$),
\[
\mathbb E_n\|\ell^{\mathrm{split}}_{n+1}-\mathbb E_n\ell^{\mathrm{split}}_{n+1}\|_{H^{-s}}^2
\le C_K\frac{\tau^2}{N}+C_K\frac{\tau^{3/2}}{N}+C_K\Bigl(\frac{\tau^2}{N}+\tau^4\Bigr)+C_K\Bigl(\frac{\tau^2}{N}+\tau^4\Bigr)
\le C_K\Bigl(\frac{\tau^{3/2}}{N}+\tau^4\Bigr)
\]
for $\tau\le1$, the surplus term dominating the $1/N$ part and the drift feedback
contributing the standalone $\tau^4$; this is the second bound in
\eqref{eq:split}. The Brownian scaling
$\mathbb E_n|X_i(\theta)-X_i(\tau)|\sim\sqrt{\tau-\theta}$ in
$M_{\mathrm{surplus}}$ is responsible for the exponent $3/2$ rather than $2$, and
is the limiting factor in the $1/N$ part of the splitting fluctuation. In the
drift-feedback part, the standalone $O(\tau^4)$ term comes from the safe
non-orthogonal estimate
$\tau\int_0^\tau\mathbb E_n|A_f(r)-\mathbb E_nA_f(r)|^2\,\mathrm dr$ of
Lemma~\ref{lem:drift-feedback-generator}, where the common empirical field
creates correlations across root particles; it is carried as a standalone
higher-order variance contribution and is globally harmless in
Lemma~\ref{lem:accum}.
\end{proof}

\begin{lemma}[Net reaction drift and frozen-map estimates]\label{lem:sens}
Under Assumption~\ref{ass:time}, define the \emph{net reaction drift}
$\mathsf L^{\mathsf B}(\rho):=c_\rho\,\rho$ (with $c_\rho=r(\mathcal P_K\rho)$),
in weak form by
\[
\langle f,\mathsf L^{\mathsf B}(\rho)\rangle:=\langle c_\rho f,\rho\rangle,
\qquad f\in H^s,
\]
and the \emph{frozen-rate expected map} of the one-step branching algorithm,
\[
\bar S_\tau(\rho):=e^{\tau c_\rho}\rho
\qquad\text{(pointwise multiplier acting on the measure }\rho).
\]
Then, for all admissible positive empirical measures
$\rho,\eta\in\mathcal A_{K,R,N}$ and $0<\tau\le1$,
\begin{align}
\|\mathsf L^{\mathsf B}(\rho)\|_{H^{-s}}\le C_K,
\qquad
\|\mathsf L^{\mathsf B}(\rho)-\mathsf L^{\mathsf B}(\eta)\|_{H^{-s}}
&\le C_K\|\rho-\eta\|_{H^{-s}},
\label{eq:LB-lip}\\
\|\bar S_\tau(\rho)-\rho\|_{H^{-s}}&\le C_K\tau,
\label{eq:Sbar-0}\\
\|\bar S_\tau(\rho)-\rho-\tau\,\mathsf L^{\mathsf B}(\rho)\|_{H^{-s}}&\le C_K\tau^2,
\label{eq:Sbar-1}
\end{align}
and the finite-difference Taylor estimate
\begin{equation}\label{eq:Sbar-taylor}
\bigl\|\bar S_\tau(\rho)-\bar S_\tau(\eta)-(\rho-\eta)
-\tau\bigl(\mathsf L^{\mathsf B}(\rho)-\mathsf L^{\mathsf B}(\eta)\bigr)\bigr\|_{H^{-s}}
\le C_K\tau^2\,\|\rho-\eta\|_{H^{-s}}.
\end{equation}
All constants are independent of $N,\tau$, and the estimates use only the smooth net
rate $c_\rho=r(\mathcal P_K\rho)$ through the multiplier and
coefficient-difference bounds of Lemma~\ref{lem:coeff}; \emph{no} derivative of the
hard split $c^\pm=\max(\pm c,0)$ is taken. In particular we do \emph{not} assert any
Lipschitz or second-order sensitivity of the \emph{exact} expected reaction map
$\rho\mapsto\mathbb E[\mathsf B^{N,\omega}_\theta\rho]$ on signed or arbitrary inputs;
that stronger property is avoided throughout, the only reaction sensitivity used
downstream being the paired genealogical estimate of Lemma~\ref{lem:Bstab}.
\end{lemma}

\begin{proof}
For $f\in H^s$, split
\[
\langle f,\mathsf L^{\mathsf B}(\rho)-\mathsf L^{\mathsf B}(\eta)\rangle
=\langle c_\eta f,\rho-\eta\rangle+\langle (c_\rho-c_\eta)f,\rho\rangle .
\]
By \eqref{eq:Bb}, $\|c_\eta\|_{C^{\lceil s\rceil}}\le C_K$, so $c_\eta$ is a
multiplier on $H^s$ (as $H^{s-1}$ is an algebra by (B4)), giving
$|\langle c_\eta f,\rho-\eta\rangle|\le C_K\|f\|_{H^s}\|\rho-\eta\|_{H^{-s}}$. By
\eqref{eq:Blip}, $\|c_\rho-c_\eta\|_{C^{\lceil s\rceil}}\le C_K\|\rho-\eta\|_{H^{-s}}$,
and the mass $\rho(\mathbb T^d)\le C_N$ is bounded on $\mathcal A_{K,R,N}$, so
$|\langle (c_\rho-c_\eta)f,\rho\rangle|\le\|(c_\rho-c_\eta)f\|_{C^0}\,\rho(\mathbb T^d)
\le C_K\|f\|_{H^s}\|\rho-\eta\|_{H^{-s}}$. Taking the supremum over $\|f\|_{H^s}\le1$
proves the Lipschitz bound in \eqref{eq:LB-lip}; the bound
$\|\mathsf L^{\mathsf B}(\rho)\|_{H^{-s}}\le C_K$ follows from
$|\langle c_\rho f,\rho\rangle|\le\|c_\rho\|_{C^0}\|f\|_{C^0}\,\rho(\mathbb T^d)\le C_K\|f\|_{H^s}$.

For the frozen map, expand the smooth multiplier pointwise:
\[
e^{\tau c_\rho}=1+\tau c_\rho+\tau^2\,q_\tau(c_\rho),
\qquad
q_\tau(z):=z^2\!\int_0^1(1-\lambda)\,e^{\lambda\tau z}\,\mathrm d\lambda,
\]
so that, since $\|c_\rho\|_{C^{\lceil s\rceil}}\le C_K$ and the smooth functional
calculus is bounded on $C^{\lceil s\rceil}$, one has
$\|q_\tau(c_\rho)\|_{C^{\lceil s\rceil}}\le C_K$ uniformly for $0<\tau\le1$. Testing
$\bar S_\tau(\rho)-\rho=(e^{\tau c_\rho}-1)\rho$ against $f$ and using the multiplier
bound with bounded mass gives \eqref{eq:Sbar-0}; testing
$\bar S_\tau(\rho)-\rho-\tau\,\mathsf L^{\mathsf B}(\rho)=\tau^2 q_\tau(c_\rho)\rho$
gives \eqref{eq:Sbar-1}. For \eqref{eq:Sbar-taylor}, the same identity yields
\[
\bar S_\tau(\rho)-\bar S_\tau(\eta)-(\rho-\eta)
-\tau\bigl(\mathsf L^{\mathsf B}(\rho)-\mathsf L^{\mathsf B}(\eta)\bigr)
=\tau^2\bigl(q_\tau(c_\rho)\rho-q_\tau(c_\eta)\eta\bigr).
\]
Writing $q_\tau(c_\rho)\rho-q_\tau(c_\eta)\eta
=q_\tau(c_\eta)(\rho-\eta)+\bigl(q_\tau(c_\rho)-q_\tau(c_\eta)\bigr)\rho$ and using
$\|q_\tau(c_\eta)\|_{C^{\lceil s\rceil}}\le C_K$ and
$\|q_\tau(c_\rho)-q_\tau(c_\eta)\|_{C^{\lceil s\rceil}}\le C_K\|\rho-\eta\|_{H^{-s}}$
(smooth functional calculus applied to \eqref{eq:Blip}), together with the bounded
mass of $\rho$, both terms are $\le C_K\|\rho-\eta\|_{H^{-s}}$ in $H^{-s}$ after
testing against $\|f\|_{H^s}\le1$. This proves \eqref{eq:Sbar-taylor}. No derivative
of $c^\pm$ is used at any step.
\end{proof}
\begin{lemma}[Stability of the exact reaction subflow under genealogical coupling]\label{lem:Bstab}
Let
\[
\nu=\frac1N\sum_{i\in I}\delta_{x_i},
\qquad
\bar\nu=\frac1N\sum_{i\in I}\delta_{\bar x_i}
\]
be two paired empirical measures with the same label set $I$ and the population
bound $|I|\le C_NN$, and let $\mathsf B^{N,\omega}_\tau\nu$,
$\mathsf B^{N,\omega}_\tau\bar\nu$ be the exact pure reaction subflows driven by a
common Harris graphical construction on a genealogical label tree, so that every
potential descendant label carries common birth and death Poisson random measures
($c_+:=\max(c,0)$, $c_-:=\max(-c,0)$ as intensities) in the two coupled systems.
Set
\[
\Delta:=\nu-\bar\nu,\qquad
\Delta^c:=\Delta-\mathbb E_n\Delta,
\qquad
\Gamma_x:=\Bigl(\frac1N\sum_{i\in I}\mathbb E_n|x_i-\bar x_i|^2\Bigr)^{1/2}.
\]
Then there is $C_K$ such that, for $\tau\le\tau_0$, the following hold. Let
$\mathcal H\supseteq\sigma(\nu,\bar\nu)$ be any $\sigma$-algebra under which the
paired inputs are fixed, and set $Z_\tau:=\mathsf B^{N,\omega}_\tau\nu
-\mathsf B^{N,\omega}_\tau\bar\nu$. The conditional mean and variance satisfy
\begin{align}
\bigl\|\mathbb E[Z_\tau\mid\mathcal H]\bigr\|_{H^{-s}}
&\le (1+C_K\tau)\|\Delta\|_{H^{-s}}+C_K\tau\,\Gamma_x^{\mathcal H},
\label{eq:Bstab-cond-mean}\\
\mathbb E\bigl[\|Z_\tau-\mathbb E[Z_\tau\mid\mathcal H]\|_{H^{-s}}^2\mid\mathcal H\bigr]
&\le\frac{C_K\tau}{N}\bigl(\|\Delta\|_{H^{-s}}+\Gamma_x^{\mathcal H}+(\Gamma_x^{\mathcal H})^2\bigr),
\label{eq:Bstab-cond-var}
\end{align}
where $\Gamma_x^{\mathcal H}=(\frac1N\sum_{i\in I}|x_i-\bar x_i|^2)^{1/2}$ is the
($\mathcal H$-measurable) labelwise displacement of the fixed inputs. Moreover, the
conditional mean is a magnitude perturbation of the identity,
$\mathbb E[Z_\tau\mid\mathcal H]=\Delta+R_\tau$ with
$\|R_\tau\|_{H^{-s}}\le C_K\tau(\|\Delta\|_{H^{-s}}+\Gamma_x^{\mathcal H})$, although
$R_\tau$ depends on the pair $(\nu,\bar\nu)$ and not on $\Delta$ alone. Specializing to
$\mathcal H:=\mathcal F_n\vee\sigma(\nu,\bar\nu)$ --- under which the paired inputs are
automatically fixed, whether or not they are themselves $\mathcal F_n$-measurable ---
and towering back to $\mathcal F_n$ via
$\mathbb E_n(\cdot)=\mathbb E_n[\mathbb E(\cdot\mid\mathcal H)]$, with
$\Delta^c=\Delta-\mathbb E_n\Delta$, then gives
\begin{equation}\label{eq:Bstab-mean}
\bigl\|\mathbb E_n(\mathsf B^{N,\omega}_\tau\nu-\mathsf B^{N,\omega}_\tau\bar\nu)\bigr\|_{H^{-s}}
\le(1+C_K\tau)\|\mathbb E_n\Delta\|_{H^{-s}}
+C_K\tau\bigl(\mathbb E_n\|\Delta^c\|_{H^{-s}}^2\bigr)^{1/2}
+C_K\tau\,\Gamma_x,
\end{equation}
and
\begin{equation}\label{eq:Bstab-var}
\begin{aligned}
&\mathbb E_n\bigl\|(\mathsf B^{N,\omega}_\tau\nu-\mathsf B^{N,\omega}_\tau\bar\nu)
-\mathbb E_n(\mathsf B^{N,\omega}_\tau\nu-\mathsf B^{N,\omega}_\tau\bar\nu)\bigr\|_{H^{-s}}^2\\
&\qquad\le(1+C_K\tau)\mathbb E_n\|\Delta^c\|_{H^{-s}}^2
+C_K\tau\|\mathbb E_n\Delta\|_{H^{-s}}^2
+C_K\tau\,\mathbb E_n(\Gamma_x^{\mathcal H})^2\\
&\qquad\quad
+\frac{C_K\tau}{N}\Bigl[\|\mathbb E_n\Delta\|_{H^{-s}}
+\bigl(\mathbb E_n\|\Delta^c\|_{H^{-s}}^2\bigr)^{1/2}+\Gamma_x\Bigr].
\end{aligned}
\end{equation}
In particular, if $x_i=\bar x_i$ for all $i\in I$ (so $\nu=\bar\nu$ labelwise),
then $\mathsf B^{N,\omega}_\tau\nu=\mathsf B^{N,\omega}_\tau\bar\nu$ pathwise and
both sides of \eqref{eq:Bstab-mean}--\eqref{eq:Bstab-var} vanish.
\end{lemma}

\begin{proof}
We first prove the conditional estimates
\eqref{eq:Bstab-cond-mean}--\eqref{eq:Bstab-cond-var}, working throughout under
$\mathcal H$, so that $\nu,\bar\nu$ (hence $\Delta$ and $\Gamma_x^{\mathcal H}$)
are fixed and the only randomness is the genealogical branching. The non-centered
estimates \eqref{eq:Bstab-mean}--\eqref{eq:Bstab-var} then follow by the tower
property as stated.

Let $\eta_\theta:=\mathsf B^{N,\omega}_\theta\nu$,
$\bar\eta_\theta:=\mathsf B^{N,\omega}_\theta\bar\nu$ for $0\le\theta\le\tau$, and
$\Delta_\theta:=\eta_\theta-\bar\eta_\theta$, $\phi_\theta:=R\Delta_\theta$, so
$\|\phi_\theta\|_{H^s}^2=\|\Delta_\theta\|_{H^{-s}}^2$. The exact reaction subflow
creates descendants, so the active label set is time-dependent: write
$\mathcal I^\eta_\theta,\mathcal I^{\bar\eta}_\theta$ for the alive genealogical
labels of the two systems at time $\theta$, with the population bound
$|\mathcal I^\eta_\theta|+|\mathcal I^{\bar\eta}_\theta|\le C_NN$, which holds on the
stopped admissible set (equivalently, under the capped population mechanism). In each
subflow particle positions are fixed and the only randomness is the genealogical
Poisson branching with rate $c_\eta$ at the current configuration; the weak
formulation reads
\[
\mathrm d\langle f,\eta_\theta\rangle
=\langle c_{\eta_\theta}f,\eta_\theta\rangle\,\mathrm d\theta
+\mathrm dM^\eta_f(\theta),
\]
and identically for $\bar\eta_\theta$ on the same genealogical clocks. Subtracting,
\[
\mathrm d\langle f,\Delta_\theta\rangle
=\langle c_{\eta_\theta}\eta_\theta-c_{\bar\eta_\theta}\bar\eta_\theta,f\rangle\,\mathrm d\theta
+\mathrm dM^\Delta_f(\theta),
\qquad M^\Delta_f=M^\eta_f-M^{\bar\eta}_f.
\]
The drift splits as
$c_{\eta_\theta}\eta_\theta-c_{\bar\eta_\theta}\bar\eta_\theta
=c_{\bar\eta_\theta}\Delta_\theta+(c_{\eta_\theta}-c_{\bar\eta_\theta})\eta_\theta$.
By \eqref{eq:Bb} and \eqref{eq:Blip} (with $q=\lceil s\rceil\le m-1$), both
contributions are bounded by $C_K\|\phi_\theta\|_{H^s}\|\Delta_\theta\|_{H^{-s}}
=C_K\|\Delta_\theta\|_{H^{-s}}^2$ when paired against $\phi_\theta$ in $H^s$, so
the drift in the It\^o energy identity for $\|\Delta_\theta\|_{H^{-s}}^2$ is
$\le C_K\|\Delta_\theta\|_{H^{-s}}^2$ (conditionally on $\mathcal H$).

The crucial point is the quadratic variation of the difference $M^\Delta_f$.
Under the genealogical Harris coupling the two flows read the same birth/death
Poisson random measures along every label, so the events split into
\emph{matched} arrivals (present in both flows) and \emph{surplus} arrivals
(present in only one, on the interval where the two intensities differ); we treat
the two QV contributions separately.

\emph{Matched-clock QV.} A matched event fires simultaneously in both flows, but
at the (possibly distinct) positions $x_\alpha$ and $\bar x_\alpha$ of a common
genealogical label $\alpha$, mapping
$\Delta_\theta\mapsto\Delta_\theta+\frac1N(\delta_{x_\alpha}-\delta_{\bar x_\alpha})$.
Since $s>d/2+1$ gives $\|\delta_x-\delta_y\|_{H^{-s}}\le C_s|x-y|$, the
matched-clock QV over $[0,\tau]$, summed over the active labels
$\alpha\in\mathcal I^\eta_\theta\cap\mathcal I^{\bar\eta}_\theta$, is in conditional
expectation
\[
\frac{C_K}{N^2}\,\mathbb E_n\!\sum_{\alpha}\int_0^\tau\|\delta_{x_\alpha}-\delta_{\bar x_\alpha}\|_{H^{-s}}^2\,\mathrm d\theta
\le\frac{C_K\tau}{N^2}\,\mathbb E_n\!\sum_{\alpha}|x_\alpha-\bar x_\alpha|^2 .
\]
Here the sum runs over \emph{descendant} labels, which copy the displacement of
their root: in the pure-reaction subflow positions are fixed, so a descendant
$\alpha$ of a matched pair $i$ sits at $x_i$ (resp.\ $\bar x_i$) and
$|x_\alpha-\bar x_\alpha|=|x_i-\bar x_i|$. The bounded reaction rate gives each root
an expected descendant multiplicity $\le e^{C_K\tau}\le C_K$, hence
\[
\mathbb E_n\!\!\sum_{\alpha\in\mathcal I^\eta_\theta\cap\mathcal I^{\bar\eta}_\theta}\!\!|x_\alpha-\bar x_\alpha|^2
\le C_K\sum_{i\in I}|x_i-\bar x_i|^2=C_KN\,(\Gamma_x^{\mathcal H})^2,
\]
so the matched-clock QV is
$\le\frac{C_K\tau}{N}(\Gamma_x^{\mathcal H})^2\le\frac{C_K\tau}{N}\Gamma_x^{\mathcal H}$,
the last step using the uniform bound $\Gamma_x^{\mathcal H}\le C$ on the torus.

\emph{Surplus-clock QV.} The per-label surplus intensity-length over
$[0,\tau]$ is
\[
\int_0^\tau\bigl|c_{\eta_\theta}(x_\alpha)-c_{\bar\eta_\theta}(\bar x_\alpha)\bigr|\,\mathrm d\theta
\le C_K\tau\bigl(\|\Delta_\theta\|_{H^{-s}}+|x_\alpha-\bar x_\alpha|\bigr),
\]
by \eqref{eq:Blip} with $q=0$ for the $\eta$-dependence and
$\|\nabla c\|_{L^\infty}\le C_K$ for the spatial dependence. Each surplus event
changes $\Delta_\theta$ by a signed atom of mass $1/N$; since
$\|\delta_x\|_{H^{-s}}\le C_s$ for $s>d/2$, each event contributes at most
$C_s/N^2$ to the QV. Therefore the surplus-clock QV, summed over the active labels
$\alpha\in\mathcal I^\eta_\theta\cup\mathcal I^{\bar\eta}_\theta$, is bounded by
\[
\frac{C_s}{N^2}\sum_{\alpha}\int_0^\theta
\bigl|c_{\eta_r}(x_\alpha)-c_{\bar\eta_r}(\bar x_\alpha)\bigr|\,\mathrm dr
\le\frac{C_K\theta}{N}\Bigl[\sup_{r\le\theta}\|\Delta_r\|_{H^{-s}}+\Gamma_x^{\mathcal H}\Bigr],
\]
using $|\mathcal I^\eta_\theta|+|\mathcal I^{\bar\eta}_\theta|\le C_NN$ and
Cauchy--Schwarz in the spatial sum. Adding the matched-clock contribution
$\frac{C_K\theta}{N}(\Gamma_x^{\mathcal H})^2$ and the surplus-clock contribution,
the total conditional QV of $\Delta_\theta$ over $[0,\theta]$ is
\begin{equation}\label{eq:QV-total}
\le\frac{C_K\theta}{N}\Bigl[\sup_{r\le\theta}\|\Delta_r\|_{H^{-s}}
+\Gamma_x^{\mathcal H}+(\Gamma_x^{\mathcal H})^2\Bigr],
\end{equation}
where we keep the squared matched-clock term explicit (it may afterwards be
relaxed via $(\Gamma_x^{\mathcal H})^2\le C\,\Gamma_x^{\mathcal H}$ on the torus,
$\Gamma_x^{\mathcal H}\le C$).
When the input matches labelwise ($x_\alpha=\bar x_\alpha$, $\nu=\bar\nu$),
$\Delta_r\equiv0$ and $\Gamma_x^{\mathcal H}=0$, so all three QV terms vanish and the
two paths coincide; the remainder is mismatch-proportional.

\emph{Cemetery-state coupling and unmatched labels.} The matched/surplus split above
prices each surplus event by the rate \emph{difference}
$|c_{\eta_\theta}(x_\alpha)-c_{\bar\eta_\theta}(\bar x_\alpha)|$, which is legitimate
only while $\alpha$ is alive in both flows. Once a surplus event leaves $\alpha$ present
in exactly one flow, its later births and deaths fire at the \emph{full} rate $O(C_K)$,
not at a difference, and $H^{-s}$ by itself does not control the number of such
unmatched labels. We therefore adjoin a cemetery state: each genealogical label carries
$x_\alpha(\theta)\in\mathbb T^d\cup\{\partial\}$ in the $\eta$-flow and
$\bar x_\alpha(\theta)\in\mathbb T^d\cup\{\partial\}$ in the $\bar\eta$-flow, with
$\partial$ marking ``dead/absent'', and we strengthen the mismatch functional to
\[
D_\theta:=\|\Delta_\theta\|_{H^{-s}}+\Gamma_\theta+U_\theta,\qquad
\Gamma_\theta:=\Bigl(\frac1N\!\!\sum_{\alpha\text{ paired}}\!\!|x_\alpha(\theta)-\bar x_\alpha(\theta)|^2\Bigr)^{1/2},
\]
\[
U_\theta:=\frac1N\,\#\bigl\{\alpha:\text{exactly one of }x_\alpha(\theta),\bar x_\alpha(\theta)\text{ is }\partial\bigr\},
\]
where $\Gamma_\theta$ runs over labels alive in both flows and $U_\theta$ counts the
unmatched ones. At $\theta=0$ the two systems share the label set $I$ with every label
alive in both, so $U_0=0$ and $D_0=\|\Delta\|_{H^{-s}}+\Gamma_x^{\mathcal H}$. The
bounded descendant multiplicity (each root has expected progeny $\le e^{C_K\tau}\le C_K$)
gives the conditional second-moment bound
$\mathbb E[\Gamma_\theta^2\mid\mathcal H]\le C_K(\Gamma_x^{\mathcal H})^2$, and hence
$\mathbb E[\Gamma_\theta\mid\mathcal H]\le C_K\Gamma_x^{\mathcal H}$ by Cauchy--Schwarz,
for $0\le\theta\le\tau$; we use these in conditional expectation rather than any
pathwise bound.

New unmatched labels are created only by surplus events, at the per-label surplus rate
$\le C_K(\|\Delta_\theta\|_{H^{-s}}+|x_\alpha-\bar x_\alpha|)$ bounded above; summed over
the $\le C_NN$ active labels with the $1/N$ weight and taking
$\mathbb E[\,\cdot\mid\mathcal H]$, this adds
$\le C_K(\mathbb E[\|\Delta_\theta\|_{H^{-s}}\mid\mathcal H]+\mathbb E[\Gamma_\theta\mid\mathcal H])$
to the growth of $\mathbb E[D_\theta\mid\mathcal H]$, using
$\frac1N\sum_\alpha|x_\alpha-\bar x_\alpha|\le\mathbb E[\Gamma_\theta\mid\mathcal H]$ in
conditional expectation by Cauchy--Schwarz. Once unmatched, a label and its descendants
branch at full rate $O(C_K)$, but each resulting label is again unmatched, so $U_\theta$
obeys a \emph{linear} birth--death inequality
$\frac{\mathrm d}{\mathrm d\theta}\mathbb E[U_\theta\mid\mathcal H]
\le C_K\,\mathbb E[U_\theta\mid\mathcal H]+C_K(\mathbb E[\|\Delta_\theta\|_{H^{-s}}\mid\mathcal H]+\mathbb E[\Gamma_\theta\mid\mathcal H])$.
Together with the mismatch-Lipschitz drift on $\mathbb E[\|\Delta_\theta\|_{H^{-s}}\mid\mathcal H]$
and $\mathbb E[\Gamma_\theta\mid\mathcal H]\le C_K\Gamma_x^{\mathcal H}$, every forcing
term is $\le C_K\,\mathbb E[D_\theta\mid\mathcal H]$, so Gr\"onwall gives the closed linear
bound
\begin{equation}\label{eq:Bstab-Dbound}
\mathbb E[D_\theta\mid\mathcal H]\le(1+C_K\theta)D_0,\qquad 0\le\theta\le\tau,
\end{equation}
which vanishes when the inputs coincide labelwise.

Every matched, surplus, or unmatched event moves $\Delta_\theta$ by a signed atom of
mass $1/N$ (matched events by a displacement atom of size $|x_\alpha-\bar x_\alpha|$,
surplus and unmatched events by a full atom), contributing $\le C_s/N^2$ per event to
the quadratic variation. Counting the three families with the $1/N$ normalization and
the full unmatched rate $O(C_K)$ yields the pointwise QV density
\begin{equation}\label{eq:Bstab-QVdensity}
\frac{\mathrm d}{\mathrm d\theta}\langle M^\Delta\rangle_\theta\le\frac{C_K}{N}\,D_\theta,
\end{equation}
which includes the unmatched contribution absent from the cumulative bound
\eqref{eq:QV-total}. Since $\|\Delta_\theta\|_{H^{-s}}\le D_\theta$, inserting
\eqref{eq:Bstab-Dbound} into \eqref{eq:Bstab-QVdensity} closes the energy identities
below with the unmatched labels now explicitly controlled, and all bounds remain
mismatch-proportional, vanishing when $\nu=\bar\nu$ labelwise.

\emph{Running mismatch in conditional mean.} We first dominate the running
disagreement in conditional mean, using only the drift split and the matched/surplus
event structure above --- not any sensitivity of an exact expected-measure map. Under
the genealogical coupling $\Delta_\theta$ is the signed sum, over active labels, of
matched displacement atoms $\frac1N(\delta_{x_\alpha}-\delta_{\bar x_\alpha})$ and
surplus full atoms $\pm\frac1N\delta$. Its conditional-mean growth is driven by the
mismatch-Lipschitz drift ($\le C_K\|\Delta_\theta\|_{H^{-s}}$ via the split
$c_{\eta_\theta}\eta_\theta-c_{\bar\eta_\theta}\bar\eta_\theta
=c_{\bar\eta_\theta}\Delta_\theta+(c_{\eta_\theta}-c_{\bar\eta_\theta})\eta_\theta$ and
\eqref{eq:Blip}), by matched-atom creation ($\le C_K\Gamma_x^{\mathcal H}$ in
expectation, using descendant multiplicity $\le e^{C_K\tau}$ and
$\frac1N\sum_\alpha|x_\alpha-\bar x_\alpha|\le C_K\Gamma_x^{\mathcal H}$), and by the
surplus rate ($\le C_K(\|\Delta_\theta\|_{H^{-s}}+\Gamma_x^{\mathcal H})$, since the
instantaneous surplus intensity is
$|c_{\eta_\theta}(x_\alpha)-c_{\bar\eta_\theta}(\bar x_\alpha)|
\le C_K(\|\Delta_\theta\|_{H^{-s}}+|x_\alpha-\bar x_\alpha|)$ using only the
$1$-Lipschitzness of $x\mapsto x^\pm$ and \eqref{eq:Blip} for the smooth net rate).
This gives the \emph{linear} integral inequality
\[
w(\theta):=\mathbb E[\|\Delta_\theta\|_{H^{-s}}\mid\mathcal H]
\le\|\Delta\|_{H^{-s}}+C_K\!\int_0^\theta\bigl(w(r)+\Gamma_x^{\mathcal H}\bigr)\,\mathrm dr
\]
(no square root), and Gr\"onwall yields the mismatch-proportional running bound
\begin{equation}\label{eq:run-mismatch}
\mathbb E[\|\Delta_\theta\|_{H^{-s}}\mid\mathcal H]\le(1+C_K\theta)\|\Delta\|_{H^{-s}}+C_K\theta\,\Gamma_x^{\mathcal H},
\qquad 0\le\theta\le\tau,
\end{equation}
which vanishes when $\Delta=0$ and $\Gamma_x^{\mathcal H}=0$.

\emph{Conditional mean via the weak drift identity.} We bound the conditional mean
through the weak drift identity for the paired flows, using the \emph{safe} net
reaction drift of Lemma~\ref{lem:sens} --- and \emph{not} any Lipschitz property of an
exact expected-measure map $S_\theta$. Writing
$m_\theta:=\mathbb E[\Delta_\theta\mid\mathcal H]$, taking conditional expectations of
the weak difference equation (the martingale part vanishing in $\mathbb E[\,\cdot\mid\mathcal H]$),
\begin{equation}\label{eq:cond-drift}
\mathbb E[Z_\tau\mid\mathcal H]=\Delta+\int_0^\tau\mathbb E\bigl[\mathsf L^{\mathsf B}(\eta_r)-\mathsf L^{\mathsf B}(\bar\eta_r)\,\big|\,\mathcal H\bigr]\,\mathrm dr,
\end{equation}
where $\mathsf L^{\mathsf B}(\rho)=c_\rho\rho$ is the net reaction drift of
Lemma~\ref{lem:sens}. By the drift Lipschitz bound \eqref{eq:LB-lip},
$\|\mathsf L^{\mathsf B}(\eta_r)-\mathsf L^{\mathsf B}(\bar\eta_r)\|_{H^{-s}}\le C_K\|\Delta_r\|_{H^{-s}}$;
taking $\mathbb E[\,\cdot\mid\mathcal H]$ and inserting the running-mismatch bound
\eqref{eq:run-mismatch},
\[
\bigl\|\mathbb E[Z_\tau\mid\mathcal H]\bigr\|_{H^{-s}}
\le\|\Delta\|_{H^{-s}}+C_K\!\int_0^\tau\bigl[(1+C_Kr)\|\Delta\|_{H^{-s}}+C_Kr\,\Gamma_x^{\mathcal H}\bigr]\,\mathrm dr
\le(1+C_K\tau)\|\Delta\|_{H^{-s}}+C_K\tau\,\Gamma_x^{\mathcal H},
\]
which is \eqref{eq:Bstab-cond-mean}. Equivalently $\mathbb E[Z_\tau\mid\mathcal H]=\Delta+R_\tau$
with the magnitude bound
$\|R_\tau\|_{H^{-s}}\le C_K\tau(\|\Delta\|_{H^{-s}}+\Gamma_x^{\mathcal H})$; this is a
genuine finite-difference estimate for the \emph{paired} flows, valid only for inputs
sharing a label set, and it does not assert that $R_\tau$ is a function of $\Delta$
alone. The same computation on $[0,\theta]$ gives the pointwise bound
\begin{equation}\label{eq:mu-theta}
\|m_\theta\|_{H^{-s}}\le(1+C_K\theta)\|\Delta\|_{H^{-s}}+C_K\theta\,\Gamma_x^{\mathcal H},
\qquad 0\le\theta\le\tau,
\qquad m_\theta:=\mathbb E[\Delta_\theta\mid\mathcal H].
\end{equation}


\emph{Conditional variance (rate-based closure).} Let
$h(\theta):=\mathbb E[\|\Delta_\theta-m_\theta\|_{H^{-s}}^2\mid\mathcal H]$, so
$h(0)=0$. The centered process
$\Delta_\theta-m_\theta$ is driven by the martingale $M^\Delta$ and the centered
drift; pairing against $R\Delta_\theta$, the drift contributes $\le C_Kh(\theta)$ to
the It\^o energy identity. The crucial point, which removes the running supremum,
is that the quadratic variation enters as a \emph{rate} (predictable density) rather
than through the cumulative bound \eqref{eq:QV-total}: the instantaneous surplus
intensity of a label $\alpha$ at time $\theta$ is
$|c_{\eta_\theta}(x_\alpha)-c_{\bar\eta_\theta}(\bar x_\alpha)|
\le C_K(\|\Delta_\theta\|_{H^{-s}}+|x_\alpha-\bar x_\alpha|)$ by \eqref{eq:Blip}, so,
summing over the $\le C_NN$ active labels with surplus full-atom size $1/N$ and
matched displacement size $|x_\alpha-\bar x_\alpha|$ and using Cauchy--Schwarz in
the spatial sum,
\[
\frac{\mathrm d}{\mathrm d\theta}\langle M^\Delta\rangle_\theta
\le\frac{C_K}{N}\Bigl(\|\Delta_\theta\|_{H^{-s}}+\Gamma_x^{\mathcal H}+(\Gamma_x^{\mathcal H})^2\Bigr)
\qquad\text{pointwise in }\theta.
\]
Thus no running $\sup_{r\le\theta}\|\Delta_r\|_{H^{-s}}$ and no maximal inequality are
needed: the supremum in \eqref{eq:QV-total} was only an artifact of integrating the
intensity-length crudely, and is unnecessary once the QV density is used pointwise.
Inserting the running-mismatch bound \eqref{eq:run-mismatch} already established into
the surplus rate, the conditional QV forcing becomes purely mismatch-proportional
with no running self-reference,
\[
\mathbb E\Bigl[\tfrac{\mathrm d}{\mathrm d\theta}\langle M^\Delta\rangle_\theta\,\Big|\,\mathcal H\Bigr]
\le\frac{C_K}{N}\bigl(\|\Delta\|_{H^{-s}}+\Gamma_x^{\mathcal H}+(\Gamma_x^{\mathcal H})^2\bigr),
\]
so the energy identity for $h$ reads, with $h(0)=0$ and a forcing independent of $h$,
\[
h'(\theta)\le C_Kh(\theta)+\frac{C_K}{N}\bigl(\|\Delta\|_{H^{-s}}+\Gamma_x^{\mathcal H}+(\Gamma_x^{\mathcal H})^2\bigr).
\]
Gr\"onwall over $[0,\tau]$ gives
\[
h(\tau)=\mathbb E\bigl[\|Z_\tau-\mathbb E[Z_\tau\mid\mathcal H]\|_{H^{-s}}^2\mid\mathcal H\bigr]
\le\frac{C_K\tau}{N}\bigl(\|\Delta\|_{H^{-s}}+\Gamma_x^{\mathcal H}+(\Gamma_x^{\mathcal H})^2\bigr),
\]
which is \eqref{eq:Bstab-cond-var}. No standalone $O(\tau/N^2)$ term arises: by
\eqref{eq:run-mismatch} the QV is driven throughout by the initial mismatch, so the
right-hand side vanishes exactly when the inputs coincide labelwise, consistent with
the pathwise vanishing asserted in the lemma.

\emph{Tower property.} Set $\mathcal H:=\mathcal F_n\vee\sigma(\nu,\bar\nu)$ (so the
inputs are $\mathcal H$-fixed, whether or not they are $\mathcal F_n$-measurable) and
integrate back to $\mathcal F_n$: the conditional mean \eqref{eq:Bstab-cond-mean}
together with $\mathbb E[Z_\tau\mid\mathcal H]=\Delta+R_\tau$ gives, by the triangle
inequality, $\mathbb E_n\|\Delta\|_{H^{-s}}\le\|\mathbb E_n\Delta\|_{H^{-s}}
+(\mathbb E_n\|\Delta^c\|_{H^{-s}}^2)^{1/2}$, and $\mathbb E_n\Gamma_x^{\mathcal H}=\Gamma_x$,
\[
\|\mathbb E_nZ_\tau\|_{H^{-s}}=\|\mathbb E_nm_\tau\|_{H^{-s}}
\le\|\mathbb E_n\Delta\|_{H^{-s}}+C_K\tau\,\mathbb E_n\|\Delta\|_{H^{-s}}+C_K\tau\,\Gamma_x
\le(1+C_K\tau)\|\mathbb E_n\Delta\|_{H^{-s}}+C_K\tau(\mathbb E_n\|\Delta^c\|_{H^{-s}}^2)^{1/2}+C_K\tau\,\Gamma_x,
\]
which is \eqref{eq:Bstab-mean}. For the variance we use the law of total variance over
the nested pair $\mathcal F_n\subseteq\mathcal H$,
\[
\mathbb E_n\|Z_\tau-\mathbb E_nZ_\tau\|_{H^{-s}}^2
=\mathbb E_n\operatorname{Var}(Z_\tau\mid\mathcal H)
+\operatorname{Var}_n\bigl(\mathbb E[Z_\tau\mid\mathcal H]\bigr).
\]
The first term is bounded by \eqref{eq:Bstab-cond-var} after taking $\mathbb E_n$,
with $\mathbb E_n\|\Delta\|_{H^{-s}}\le\|\mathbb E_n\Delta\|_{H^{-s}}
+(\mathbb E_n\|\Delta^c\|_{H^{-s}}^2)^{1/2}$, $\mathbb E_n\Gamma_x^{\mathcal H}=\Gamma_x$,
and the torus relaxation $\mathbb E_n(\Gamma_x^{\mathcal H})^2\le C\,\mathbb E_n\Gamma_x^{\mathcal H}=C\Gamma_x$;
it is $\le\frac{C_K\tau}{N}[\|\mathbb E_n\Delta\|_{H^{-s}}
+(\mathbb E_n\|\Delta^c\|_{H^{-s}}^2)^{1/2}+\Gamma_x]$. For the second term we do
\emph{not} treat $T_\tau:=\mathbb E[Z_\tau\mid\mathcal H]$ as a Lipschitz function
of $\Delta$. Instead we use only the magnitude decomposition
$T_\tau=\Delta+R_\tau$ with
$\|R_\tau\|_{H^{-s}}\le C_K\tau(\|\Delta\|_{H^{-s}}+\Gamma_x^{\mathcal H})$ from
\eqref{eq:cond-drift}--\eqref{eq:mu-theta}: writing
$T_\tau-\mathbb E_nT_\tau=\Delta^c+(R_\tau-\mathbb E_nR_\tau)$,
Cauchy--Schwarz and Young give
\[
\operatorname{Var}_n(T_\tau)
\le\mathbb E_n\|\Delta^c\|_{H^{-s}}^2
+2\bigl(\mathbb E_n\|\Delta^c\|_{H^{-s}}^2\bigr)^{1/2}\bigl(\mathbb E_n\|R_\tau\|_{H^{-s}}^2\bigr)^{1/2}
+\mathbb E_n\|R_\tau\|_{H^{-s}}^2,
\]
and with $\mathbb E_n\|R_\tau\|_{H^{-s}}^2\le C_K^2\tau^2\,\mathbb E_n(\|\Delta\|_{H^{-s}}+\Gamma_x^{\mathcal H})^2
\le 2C_K^2\tau^2(\|\mathbb E_n\Delta\|_{H^{-s}}^2+\mathbb E_n\|\Delta^c\|_{H^{-s}}^2+\mathbb E_n(\Gamma_x^{\mathcal H})^2)$,
\begin{equation}\label{eq:Ttau-vardecomp}
\operatorname{Var}_n(T_\tau)
\le(1+C_K\tau)\,\mathbb E_n\|\Delta^c\|_{H^{-s}}^2+C_K\tau\,\|\mathbb E_n\Delta\|_{H^{-s}}^2
+C_K\tau\,\mathbb E_n(\Gamma_x^{\mathcal H})^2 .
\end{equation}
Here the genuinely base-measure-dependent part of $T_\tau$ has been confined to the
higher-order terms $C_K\tau\|\mathbb E_n\Delta\|_{H^{-s}}^2$ and
$C_K\tau\,\mathbb E_n(\Gamma_x^{\mathcal H})^2$. Adding the two
contributions gives \eqref{eq:Bstab-var}.
\end{proof}

In the transport application of Lemma~\ref{lem:transport} below, the paired
inputs are $\nu_\tau:=\mathsf A^{N,\omega}_\tau(\mu^N_n)$ and
$\bar\nu_\tau:=\mathsf A^{N,\omega}_{\tau,h}(\mu^N_n)$, which depend on the
transport Brownian noise. To apply this lemma cleanly, one first conditions on
the transport $\sigma$-algebra $\mathcal G_n$ of \eqref{eq:filtrations}, under
which the paired inputs are fixed and the only
remaining randomness in $\mathsf B^{N,\omega}_\tau$ comes from the genealogical
branching clocks; applying \eqref{eq:Bstab-mean}--\eqref{eq:Bstab-var} conditionally
on $\mathcal G_n$ and then taking $\mathbb E_n$, together with the
pre-branching transport estimates and the labelwise displacement bound
\[
\Gamma_x=\Bigl(\frac1N\sum_{i\in\mathcal I_n}\mathbb E_n|Y_i(\tau)-\bar Y_i(\tau)|^2\Bigr)^{1/2}\le C_K\tau^{3/2},
\]
yields the post-branching bounds $\|\mathbb E_n\ell^{\mathrm{tr}}_{n+1}\|_{H^{-s}}\le C_K\tau^2$
(the extra mean contribution $C_K\tau\Gamma_x\le C_K\tau^{5/2}$ being absorbed) and
$\mathbb E_n\|\ell^{\mathrm{tr}}_{n+1}-\mathbb E_n\ell^{\mathrm{tr}}_{n+1}\|_{H^{-s}}^2
\le C_K(\tau^2/N+\tau^4)$, the standalone $\tau^4$ arising from the
$C_K\tau\,\mathbb E_n(\Gamma_x^{\mathcal H})^2\le C_K\tau^4$ term of
\eqref{eq:Bstab-var}. Thus the $1/N$ part of the remainder in
\eqref{eq:Bstab-mean}--\eqref{eq:Bstab-var} is $O(C_K\tau^2/N)$ in that
application, plus a standalone higher-order $O(C_K\tau^4)$ variance contribution;
both are mismatch-proportional in general.

\begin{lemma}[Centered fluctuation of a stopped finite-particle cylinder functional]\label{lem:cyl-centered}
Let $X_\theta$ be the stopped finite-particle configuration process of one substep (the
active labeled positions together with the finitely many projected modes
$|\mathbf k|_\infty\le K$ that enter the coefficients), and let $G$ be a real cylinder
functional of $X_\theta$ that is $L_K$-Lipschitz in the cylinder coordinates and whose
continuous-martingale part has quadratic-variation density $\le C_K/N$. Then, for
$0\le\theta\le\tau$,
\begin{equation}\label{eq:cyl-centered}
\mathbb E_n\bigl|G(X_\theta)-\mathbb E_nG(X_\theta)\bigr|^2\le C_K\,\frac\theta N,
\end{equation}
with $C_K$ depending on $L_K$, the coefficient bounds and the stopped admissibility
constants, but not on $N$.
\end{lemma}

\begin{proof}
Let $\Xi_\theta:=\sum_j|c_j(\theta)-\mathbb E_nc_j(\theta)|^2$ be the centered energy of
the cylinder coordinates $c_j(\theta)$ of $X_\theta$ (the finitely many projected modes
and the per-label pairings). Each $c_j$ is a one-dimensional semimartingale whose
centered drift is Lipschitz in the centered coordinates with constant $C_K$ and whose
martingale part has QV density $\le C_K/N$; applying It\^o to $|c_j-\mathbb E_nc_j|^2$ and
Gr\"onwall from the vanishing initial value (exactly as for a single mode, cf.\
\eqref{eq:b-centered}), then summing the finitely many coordinates,
\[
\mathbb E_n\Xi_\theta\le C_K\,\frac\theta N,\qquad 0\le\theta\le\tau .
\]
Since $G$ is $L_K$-Lipschitz in the cylinder coordinates, the independent-copy identity
$\operatorname{Var}_n\!\bigl(G(X_\theta)\bigr)
=\tfrac12\mathbb E_n\bigl|G(X_\theta)-G(\widetilde X_\theta)\bigr|^2$, with
$\widetilde X_\theta$ an independent copy of $X_\theta$ given $\mathcal F_n$, gives
\[
\mathbb E_n\bigl|G(X_\theta)-\mathbb E_nG(X_\theta)\bigr|^2
\le L_K^2\,\mathbb E_n\Xi_\theta\le C_K\,\frac\theta N,
\]
which is \eqref{eq:cyl-centered}. The deterministic $O(\theta)$ drift of $G(X_\theta)$ is
removed by centering, so no $\theta^2$ term appears.
\end{proof}

\begin{lemma}[Pre-branching transport defect]\label{lem:transport-pre}
Under Assumption~\ref{ass:time}, let
$\nu_\tau:=\mathsf A^{N,\omega}_\tau(\mu^N_n)$ and
$\bar\nu_\tau:=\mathsf A^{N,\omega}_{\tau,h}(\mu^N_n)$ be the exact transport--diffusion
and Euler--Maruyama outputs over one step, with exact and EM particle trajectories
$Y_i,\bar Y_i$, and set $\widetilde\ell^{\mathrm{tr}}_{n+1}:=\nu_\tau-\bar\nu_\tau$.
Then
\begin{equation}\label{eq:transport-pre}
\|\mathbb E_n\widetilde\ell^{\mathrm{tr}}_{n+1}\|_{H^{-s}}\le C_K\,\tau^2,
\qquad
\mathbb E_n\|\widetilde\ell^{\mathrm{tr}}_{n+1}-\mathbb E_n\widetilde\ell^{\mathrm{tr}}_{n+1}\|_{H^{-s}}^2
\le C_K\,\frac{\tau^2}{N},
\end{equation}
and the labelwise exact/EM displacement obeys
\begin{equation}\label{eq:transport-disp}
\mathbb E_n\frac1N\sum_{i\in\mathcal I_n}|Y_i(\tau)-\bar Y_i(\tau)|^2\le C_K\,\tau^3 .
\end{equation}
\end{lemma}

\begin{proof}
This is exactly the content of Steps~1--3 in the proof of Lemma~\ref{lem:transport}
below: the mean bound is established in Steps~1--2, the centered-variance bound in
Step~3 via the cylinder-functional centered estimate of Lemma~\ref{lem:cyl-centered}
(the semimartingale bound \eqref{eq:b-centered}), and the
displacement bound \eqref{eq:transport-disp} in \eqref{eq:Gamma-bound}. These three
steps involve only the transport substep and the common empirical field $\nu_\theta$,
and do not invoke the reaction subflow; we therefore state them separately and apply
Lemma~\ref{lem:Bstab} to the two paired post-transport inputs in Step~4.
\end{proof}

\begin{lemma}[Transport discretization defect]\label{lem:transport}
Under Assumption~\ref{ass:time}, the transport defect
$\ell^{\mathrm{tr}}_{n+1}=\mathsf B^{N,\omega}_\tau(\mathsf A^{N,\omega}_\tau(\mu^N_n))
-\mathsf B^{N,\omega}_\tau(\mathsf A^{N,\omega}_{\tau,h}(\mu^N_n))$ (the exact
nonlinear transport--diffusion subflow against its Euler--Maruyama discretization,
followed in either case by the exact pure reaction subflow) satisfies
\begin{equation}\label{eq:transport}
\|\mathbb E_n\ell^{\mathrm{tr}}_{n+1}\|_{H^{-s}}\le C_K\,\tau^2,
\qquad
\mathbb E_n\|\ell^{\mathrm{tr}}_{n+1}-\mathbb E_n\ell^{\mathrm{tr}}_{n+1}\|_{H^{-s}}^2
\le C_K\Bigl(\frac{\tau^2}{N}+\tau^4\Bigr).
\end{equation}
\end{lemma}

\begin{proof}
Write $\nu_\tau:=\mathsf A^{N,\omega}_\tau(\mu^N_n)$ for the exact nonlinear
transport--diffusion output and $\bar\nu_\tau:=\mathsf A^{N,\omega}_{\tau,h}(\mu^N_n)$
for the EM output, and set the \emph{pre-branching} transport defect
$\tilde\ell^{\mathrm{tr}}_{n+1}:=\nu_\tau-\bar\nu_\tau$. The post-branching defect
$\ell^{\mathrm{tr}}_{n+1}=\mathsf B^{N,\omega}_\tau\nu_\tau
-\mathsf B^{N,\omega}_\tau\bar\nu_\tau$ is the action of the common exact reaction
subflow on the two paired inputs; by Lemma~\ref{lem:Bstab} the post-branching
estimates reduce to estimates of $\tilde\ell^{\mathrm{tr}}_{n+1}$ together with the
labelwise displacement
$\Gamma_x=\bigl(\frac1N\sum_{i\in\mathcal I_n}\mathbb E_n|Y_i(\tau)-\bar Y_i(\tau)|^2\bigr)^{1/2}$,
where $Y_i,\bar Y_i$ are the exact and EM particle trajectories. We therefore
proceed in four steps: (1) pre-branching mean via weak generator expansion;
(2) time-Lipschitz of the cylinder drift; (3) pre-branching centered variance
via Fourier basis --- Steps~1--3 together establishing the pre-branching
Lemma~\ref{lem:transport-pre} --- and (4) post-branching reduction using
\eqref{eq:Bstab-mean}--\eqref{eq:Bstab-var}.

\emph{Step 1: pre-branching weak mean.} For $f\in H^s\hookrightarrow C^4$ (by
(B4), $s>d/2+4$), introduce the cylinder drift functional
\[
F_f(\nu):=\langle b_\nu\!\cdot\!\nabla f+D\Delta f,\,\nu\rangle.
\]
Its nonlinear coefficient dependence is finite-dimensional: $b_\nu=-a(\mathcal P_K\nu)$
depends on $\nu$ only through the finitely many Fourier coefficients
$\{\hat\nu(\mathbf k):|\mathbf k|_\infty\le K\}$; the remaining dependence on
$\nu$ is linear through the pairing against smooth functions involving $f$. By
the weak formulation of the exact nonlinear transport subflow,
\[
\mathbb E_n\langle f,\nu_\tau\rangle=\langle f,\nu_0\rangle
+\int_0^\tau\mathbb E_nF_f(\nu_\theta)\,\mathrm d\theta,
\qquad\nu_0=\mu^N_n,
\]
and by a Taylor expansion of $f(\mathbf X^i_n+b_{\nu_0}(\mathbf X^i_n)\tau
+\sqrt{2D}\,\Delta\mathbf W^N_i)$ to second order, averaged over the Brownian
increment $\Delta\mathbf W^N_i$ and the population (the first-order term
reproduces $\tau F_f(\nu_0)$; the second-order term involves
$b_{\nu_0}\!\cdot\!\nabla(b_{\nu_0}\!\cdot\!\nabla f)$ and $D^2\Delta^2 f$,
controlled by $\|f\|_{C^4}$ with coefficient factors $\|b_{\nu_0}\|_{C^2}\le C_K$
from Lemma~\ref{lem:coeff}, the diffusion term carrying no power of $K$; the EM
second-order term need not match the exact $\tfrac12\mathcal L^2f$, but is of the
same order $O(C_K\tau^2\|f\|_{C^4})$),
\[
\mathbb E_n\langle f,\bar\nu_\tau\rangle=\langle f,\nu_0\rangle+\tau\,F_f(\nu_0)
+O\bigl(C_K\tau^2\|f\|_{C^4}\bigr).
\]
Subtracting,
\[
\mathbb E_n\langle f,\tilde\ell^{\mathrm{tr}}_{n+1}\rangle
=\int_0^\tau\bigl(\mathbb E_nF_f(\nu_\theta)-F_f(\nu_0)\bigr)\,\mathrm d\theta
+O\bigl(C_K\tau^2\|f\|_{C^4}\bigr).
\]

\emph{Step 2: time-Lipschitz of $F_f$.} We claim
\begin{equation}\label{eq:F-Lip}
|\mathbb E_nF_f(\nu_\theta)-F_f(\nu_0)|\le C_K\,\theta\,\|f\|_{C^4},
\qquad 0\le\theta\le\tau.
\end{equation}
Apply It\^o (in measure form) to $F_f$ along the exact transport subflow,
\[
\mathbb E_nF_f(\nu_\theta)=F_f(\nu_0)+\int_0^\theta\mathbb E_n[\mathsf L^{\mathsf A}_NF_f(\nu_r)]\,\mathrm dr,
\]
where $\mathsf L^{\mathsf A}_N$ is the generator of $\mathsf A^{N,\omega}_\tau$ in
measure form. Writing $F_f(\nu)=\langle g_f^\nu,\nu\rangle$ with
$g_f^\nu:=b_\nu\!\cdot\!\nabla f+D\Delta f$,
\[
\mathsf L^{\mathsf A}_N F_f(\nu)
=\underbrace{\langle b_\nu\!\cdot\!\nabla g_f^\nu+D\Delta g_f^\nu,\nu\rangle}_{\text{(I): linear pairing}}
+\underbrace{\bigl\langle(\partial_\nu g_f^\nu)\cdot(b_\nu\!\cdot\!\nabla\!+\!D\Delta)\nu,\,\nu\bigr\rangle}_{\text{(II): measure-derivative of }b_\nu}
+\underbrace{\frac{D}{N}\,\mathrm{QV-term}}_{\text{(III): order $1/N$}}.
\]
Each group is bounded by $C_K\|f\|_{C^4}$:
\begin{itemize}
\item (I) substitutes $g_f^\nu=b_\nu\!\cdot\!\nabla f+D\Delta f$, producing
$b\!\cdot\!\nabla(b\!\cdot\!\nabla f)+b\!\cdot\!\nabla(D\Delta f)+D\Delta(b\!\cdot\!\nabla f)+D^2\Delta^2f$,
which involves up to $\Delta^2f$ and depends on $b,\nabla b,\Delta b$ only;
each factor of $b$ and its derivatives is $\le C_K$ by \eqref{eq:Bb},
giving $|\langle\cdot,\nu\rangle|\le C_K C_N\|f\|_{C^4}\le C_K\|f\|_{C^4}$;
\item (II) uses that $\partial_\nu b_\nu=-a'(\mathcal P_K\nu)\nabla\mathcal P_K(\cdot)$
is finite-rank in Fourier modes $|\mathbf k|_\infty\le K$ and smooth in those modes
by (B1); paired against $(b_\nu\!\cdot\!\nabla+D\Delta)\nu$ this contributes at
most $\|f\|_{C^3}$ in $f$ and $C_K$ in coefficients;
\item (III) is the quadratic-variation correction in It\^o's formula for a
functional of the empirical measure, of order $1/N$ times $\|f\|_{C^4}$ and
fixed-$K$ factors, hence $\le C_K\|f\|_{C^4}/N$, absorbed into $C_K\|f\|_{C^4}$.
\end{itemize}
Therefore $|\mathsf L^{\mathsf A}_N F_f(\nu)|\le C_K\|f\|_{C^4}$ uniformly in
$\nu$, and integrating gives \eqref{eq:F-Lip}. Inserting into Step 1 and using
$H^s\hookrightarrow C^4$,
\[
|\mathbb E_n\langle f,\tilde\ell^{\mathrm{tr}}_{n+1}\rangle|\le C_K\tau^2\|f\|_{C^4}\le C_K\tau^2\|f\|_{H^s},
\]
and taking the supremum over $\|f\|_{H^s}\le1$,
$\|\mathbb E_n\tilde\ell^{\mathrm{tr}}_{n+1}\|_{H^{-s}}\le C_K\tau^2$.

\emph{Step 3: pre-branching centered variance via Fourier basis.} Let
$\{e_\ell\}_{\ell\in\mathbb Z^d}$ be the Fourier basis on $\mathbb T^d$ with
weights $(1+|\ell|^2)^{-s/2}$, so that for any signed measure $m$,
\[
\|m\|_{H^{-s}}^2=\sum_{\ell\in\mathbb Z^d}(1+|\ell|^2)^{-s}|\hat m(\ell)|^2.
\]
Apply this to $m=\tilde\ell^{\mathrm{tr}}_{n+1}-\mathbb E_n\tilde\ell^{\mathrm{tr}}_{n+1}$
and take $\mathbb E_n$:
\[
\mathbb E_n\|\tilde\ell^{\mathrm{tr}}_{n+1}-\mathbb E_n\tilde\ell^{\mathrm{tr}}_{n+1}\|_{H^{-s}}^2
=\sum_{\ell}(1+|\ell|^2)^{-s}\,\mathbb E_n\bigl|\hat m(\ell)\bigr|^2,
\qquad
\hat m(\ell)=\langle e_\ell,m\rangle.
\]
For each mode $\ell$, the per-particle deviation
$Y_i(\tau)-\bar Y_i(\tau)=\int_0^\tau[b_{\nu_r}(Y_i(r))-b_{\nu_0}(\mathbf X^i_n)]\,\mathrm dr$
splits as
\begin{equation}\label{eq:disp-split}
b_{\nu_r}(Y_i(r))-b_{\nu_0}(\mathbf X^i_n)
=\underbrace{b_{\nu_0}(Y_i(r))-b_{\nu_0}(\mathbf X^i_n)}_{\text{spatial}}
+\underbrace{b_{\nu_r}(Y_i(r))-b_{\nu_0}(Y_i(r))}_{\text{measure}},
\end{equation}
with the spatial part bounded by $\|\nabla b\|_{L^\infty}\le C_K$ times
$|Y_i(r)-\mathbf X^i_n|$ (a Brownian-scale displacement,
$\mathbb E_n|Y_i(r)-\mathbf X^i_n|^2\le Cr$), and the measure part bounded by
\eqref{eq:Blip} with $q=0$ times $\|\nu_r-\nu_0\|_{H^{-s}}$. The finite-mode
Fourier coefficients $\hat\nu_\theta(\mathbf k)$ ($|\mathbf k|_\infty\le K$)
satisfy, by It\^o on each $\hat\nu_\theta(\mathbf k)=\langle\varphi_{\mathbf k},\nu_\theta\rangle$,
\[
\mathrm d\hat\nu_\theta(\mathbf k)
=\langle b_{\nu_\theta}\!\cdot\!\nabla\varphi_{\mathbf k}+D\Delta\varphi_{\mathbf k},\nu_\theta\rangle\,\mathrm d\theta
+\frac{\sqrt{2D}}N\sum_{i}\nabla\varphi_{\mathbf k}(Y_i(\theta))\!\cdot\!\mathrm d\mathbf W^N_i,
\]
with drift bounded by $C_K$ (Lemma~\ref{lem:coeff} and population bound) and QV
density $\le C_K/N$. This yields two distinct estimates, which we keep separate.
The \emph{total increment} (drift plus martingale) is
\begin{equation}\label{eq:b-fluct}
\mathbb E_n|\hat\nu_\theta(\mathbf k)-\hat\nu_0(\mathbf k)|^2\le C_K\Bigl(\theta^2+\frac\theta N\Bigr),
\qquad\text{hence}\qquad
\mathbb E_n\|b_{\nu_\theta}-b_{\nu_0}\|_{C^1}^2\le C_K\Bigl(\theta^2+\frac\theta N\Bigr),
\end{equation}
the $\theta^2$ coming from the deterministic drift and the $\theta/N$ from the
martingale. For the centered fluctuation we argue directly. Subtracting its
conditional mean, the centered coefficient
$\xi_\theta(\mathbf k):=\hat\nu_\theta(\mathbf k)-\mathbb E_n\hat\nu_\theta(\mathbf k)$
satisfies $\xi_0(\mathbf k)=0$ and, from the It\^o equation above,
\[
\mathrm d\xi_\theta(\mathbf k)
=\bigl(\langle\cdot,\nu_\theta\rangle-\mathbb E_n\langle\cdot,\nu_\theta\rangle\bigr)\,\mathrm d\theta
+\frac{\sqrt{2D}}N\sum_i\nabla\varphi_{\mathbf k}(Y_i(\theta))\!\cdot\!\mathrm d\mathbf W^N_i,
\]
where the centered drift is a cylinder functional of the stopped finite-particle
configuration --- the coefficient $b_{\nu_\theta}$ entering through its finitely many
projected modes, the pairing through the particle positions --- and is Lipschitz in the
configuration's cylinder coordinates with constant
$C_K$ and the martingale has QV density $\le C_K/N$. By It\^o's formula for
$|\xi_\theta(\mathbf k)|^2$, $\frac{\mathrm d}{\mathrm d\theta}
\mathbb E_n|\xi_\theta(\mathbf k)|^2\le C_K\mathbb E_n|\xi_\theta(\mathbf k)|^2+C_K/N$,
and Gr\"onwall from $\xi_0=0$ gives
\begin{equation}\label{eq:b-centered}
\sup_{|\mathbf k|_\infty\le K}\mathbb E_n|\hat\nu_\theta(\mathbf k)-\mathbb E_n\hat\nu_\theta(\mathbf k)|^2
\le C_K\frac\theta N,
\qquad 0\le\theta\le\tau.
\end{equation}
The deterministic $O(\theta)$ drift of $\hat\nu_\theta(\mathbf k)$ has been
removed by centering, so no $\theta^2$ term appears. The same argument applies to
the EM interpolation $\bar\nu_\theta$ and to the joint empirical process
$\Xi^N_\theta=\frac1N\sum_i\delta_{(Y_i(\theta),\bar Y_i(\theta),\mathbf X^i_n)}$,
whose cylinder functionals of the stopped finite-particle configuration obey the same centered bound
\eqref{eq:b-centered} by the identical It\^o argument; this is the instance of
Lemma~\ref{lem:cyl-centered} for those functionals.
Combining the two contributions of \eqref{eq:disp-split} via Cauchy--Schwarz in
time,
\[
\frac1N\sum_{i\in\mathcal I_n}\mathbb E_n|Y_i(\tau)-\bar Y_i(\tau)|^2
\le C_K\tau^3+C_K\frac{\tau^3}{N}\le C_K\tau^3.
\]
This is the labelwise strong defect; the displacement is coupled across particles
through the common empirical field $\nu_r$, so we do \emph{not} treat the centered
modes as an independent per-particle sum. Instead we estimate each Fourier mode
$e_\ell$ directly. Writing $\Delta_i(t):=Y_i(t)-\bar Y_i(t)$, the mode evolves as
\[
\mathrm d\langle e_\ell,\nu_t-\bar\nu_t\rangle
=B_\ell(t)\,\mathrm dt
+\frac{\sqrt{2D}}N\sum_{i}\bigl(\nabla e_\ell(Y_i(t))-\nabla e_\ell(\bar Y_i(t))\bigr)\!\cdot\!\mathrm d\mathbf W^N_i(t),
\]
where $B_\ell$ is the drift difference. For the martingale part $M_\ell$, It\^o
isometry with $|\nabla e_\ell(Y)-\nabla e_\ell(\bar Y)|\le C_d|\ell|^2|Y-\bar Y|$
(two derivatives of $e_\ell$, giving the factor $|\ell|^2$ before squaring) yields
\[
\mathbb E_n|M_\ell(\tau)|^2
\le\frac{C_d(1+|\ell|^2)^2}{N^2}\sum_i\int_0^\tau\mathbb E_n|\Delta_i(t)|^2\,\mathrm dt
\le C_K(1+|\ell|^2)^2\,\frac{\tau^4}{N}
\le C_K(1+|\ell|^2)^2\,\frac{\tau^3}{N},
\]
using $\frac1N\sum_i\mathbb E_n|\Delta_i(t)|^2\le C_Kt^3\le C_K\tau^3$ and $\tau\le1$.
For the drift centered part, $B_\ell$ pairs $e_\ell$ against the drift
difference $b_{\nu_t}\!\cdot\!\nabla(\cdot)+D\Delta(\cdot)$; the diffusion term
contributes $D(\Delta e_\ell(Y_i)-\Delta e_\ell(\bar Y_i))$, whose Lipschitz
bound involves a \emph{third} derivative of $e_\ell$, so the relevant weight is
$(1+|\ell|^2)^3$. More precisely, $B_\ell(t)$ is a cylinder functional of the
joint empirical process
$\Xi^N_t=\frac1N\sum_i\delta_{(Y_i(t),\bar Y_i(t),\mathbf X^i_n)}$ through the
finitely many modes $|\mathbf k|_\infty\le K$ of $\nu_t$ (entering $b_{\nu_t}$) and
the per-particle pairings of $e_\ell$ and its derivatives; differentiating the
pairing shows it is Lipschitz in those cylinder coordinates with constant
$\le C_K(1+|\ell|^2)^{3/2}$ (one factor $(1+|\ell|^2)^{1/2}$ per derivative of
$e_\ell$, up to the third). Consequently the centered fluctuation
$B_\ell(t)-\mathbb E_nB_\ell(t)$ is driven by the finite-particle centered estimate
\eqref{eq:b-centered} for $\Xi^N_t$ (not by the total increment
\eqref{eq:b-fluct}, which contains the $O(t)$ drift): subtracting
$\mathbb E_nB_\ell$ removes the deterministic part and leaves, by the Lipschitz
bound applied to the centered cylinder coordinates of QV density $\le C_K/N$, a
martingale of QV density $\le C_K(1+|\ell|^2)^3/N$, giving
\[
\mathbb E_n|B_\ell(t)-\mathbb E_nB_\ell(t)|^2\le C_K(1+|\ell|^2)^3\,\frac tN,
\qquad
\mathbb E_n\Bigl|\int_0^\tau(B_\ell(t)-\mathbb E_nB_\ell(t))\,\mathrm dt\Bigr|^2
\le C_K(1+|\ell|^2)^3\,\frac{\tau^3}{N}.
\]
Summing the two contributions against $(1+|\ell|^2)^{-s}$ (the larger weight
$(1+|\ell|^2)^3$ controlling),
\[
\mathbb E_n\|\tilde\ell^{\mathrm{tr}}_{n+1}-\mathbb E_n\tilde\ell^{\mathrm{tr}}_{n+1}\|_{H^{-s}}^2
\le C_K\frac{\tau^3}{N}\sum_{\ell\in\mathbb Z^d}(1+|\ell|^2)^{-s}(1+|\ell|^2)^3
\le C_K\frac{\tau^3}{N}\le C_K\frac{\tau^2}{N},
\]
the $\ell$-sum $\sum_\ell(1+|\ell|^2)^{-(s-3)}$ being finite since $s-3>d/2$ by
(B4) ($s>d/2+4$), and the last step using $\tau\le1$.

\emph{Step 4: post-branching reduction.} The pre-branching inputs
$\nu_\tau=\mathsf A^{N,\omega}_\tau(\mu^N_n)$ and
$\bar\nu_\tau=\mathsf A^{N,\omega}_{\tau,h}(\mu^N_n)$ depend on the transport
Brownian noise. Let $\mathcal G_n$ be the transport $\sigma$-algebra of
Lemma~\ref{lem:Bstab}, set $Z:=\mathsf B^{N,\omega}_\tau\nu_\tau
-\mathsf B^{N,\omega}_\tau\bar\nu_\tau=\ell^{\mathrm{tr}}_{n+1}$ and
$\mu_{\mathcal G}:=\mathbb E[Z\mid\mathcal G_n]$, the conditional-mean map
\[
T_\tau(\nu_\tau,\bar\nu_\tau):=\mathbb E\bigl[\mathsf B^{N,\omega}_\tau\nu_\tau
-\mathsf B^{N,\omega}_\tau\bar\nu_\tau\mid\mathcal G_n\bigr]=\mu_{\mathcal G}.
\]
Conditionally on $\mathcal G_n$ the inputs $\nu_\tau,\bar\nu_\tau$ are \emph{fixed},
so the conditional input difference $\Delta=\nu_\tau-\bar\nu_\tau$ is
$\mathcal G_n$-measurable and its conditional centered part vanishes; the only
randomness left in $Z$ given $\mathcal G_n$ is the genealogical branching. The
labelwise displacement $\Gamma_x=(\frac1N\sum_i|Y_i(\tau)-\bar Y_i(\tau)|^2)^{1/2}$
is itself $\mathcal G_n$-measurable and random, with
\begin{equation}\label{eq:Gamma-bound}
\mathbb E_n\Gamma_x^2=\mathbb E_n\frac1N\sum_i|Y_i(\tau)-\bar Y_i(\tau)|^2\le C_K\tau^3,
\qquad\text{hence}\qquad \mathbb E_n\Gamma_x\le C_K\tau^{3/2}.
\end{equation}
Lemma~\ref{lem:Bstab} applied conditionally (with $\mathcal H=\mathcal G_n$, under
which $\Delta=\nu_\tau-\bar\nu_\tau$ and $\Gamma_x^{\mathcal H}$ are fixed) gives
\[
\|\mu_{\mathcal G}\|_{H^{-s}}\le(1+C_K\tau)\|\Delta\|_{H^{-s}}+C_K\tau\,\Gamma_x^{\mathcal H},
\qquad
\operatorname{Var}(Z\mid\mathcal G_n)\le\frac{C_K\tau}{N}\bigl(\|\Delta\|_{H^{-s}}+\Gamma_x^{\mathcal H}+(\Gamma_x^{\mathcal H})^2\bigr),
\]
the conditional variance being driven only by $\Gamma_x^{\mathcal H}$ (and its square)
and the (now $\mathcal G_n$-fixed) $\|\Delta\|_{H^{-s}}$, not by any $\Delta^c$; no
sensitivity of an exact expected reaction map is used.

For the mean, $\mathbb E_n Z=\mathbb E_n\mu_{\mathcal G}$, so the tower bound
\eqref{eq:Bstab-mean} gives
$\|\mathbb E_nZ\|_{H^{-s}}\le(1+C_K\tau)\|\mathbb E_n\Delta\|_{H^{-s}}
+C_K\tau(\mathbb E_n\|\Delta^c\|_{H^{-s}}^2)^{1/2}+C_K\tau\,\Gamma_x
\le C_K\tau^2+C_K\tau\,(C_K\tau^2/N)^{1/2}+C_K\tau\cdot C_K\tau^{3/2}\le C_K\tau^2$,
using the Step~1--3 pre-branching estimates $\|\mathbb E_n\Delta\|_{H^{-s}}\le C_K\tau^2$,
$\mathbb E_n\|\Delta^c\|_{H^{-s}}^2\le C_K\tau^2/N$, the displacement bound
$\Gamma_x=(\mathbb E_n(\Gamma_x^{\mathcal H})^2)^{1/2}\le C_K\tau^{3/2}$ from
\eqref{eq:Gamma-bound}, and $N\ge1$ (the extra term $C_K\tau\Gamma_x\le C_K\tau^{5/2}$
being absorbed into $C_K\tau^2$).

For the variance we use the total variance decomposition over $\mathcal G_n$,
\[
\mathbb E_n\|Z-\mathbb E_nZ\|_{H^{-s}}^2
=\underbrace{\mathbb E_n\operatorname{Var}(Z\mid\mathcal G_n)}_{\text{branching}}
+\underbrace{\mathbb E_n\|\mu_{\mathcal G}-\mathbb E_n\mu_{\mathcal G}\|_{H^{-s}}^2}_{\text{transport}}.
\]
The branching term, by the conditional variance above with
$\mathbb E_n\|\Delta\|_{H^{-s}}\le(\mathbb E_n\|\Delta\|_{H^{-s}}^2)^{1/2}\le C_K\tau$,
$\mathbb E_n\Gamma_x^{\mathcal H}=\Gamma_x\le C_K\tau^{3/2}$ and
$\mathbb E_n(\Gamma_x^{\mathcal H})^2\le C_K\tau^3$, is
\[
\mathbb E_n\operatorname{Var}(Z\mid\mathcal G_n)
\le\frac{C_K\tau}{N}\bigl(C_K\tau+C_K\tau^{3/2}+C_K\tau^3\bigr)\le C_K\frac{\tau^2}{N}.
\]
The transport term is the variance of the conditional-mean map
$T_\tau:=\mu_{\mathcal G}$ (here $\mathcal H=\mathcal G_n$). By Lemma~\ref{lem:Bstab},
$T_\tau=\Delta+R_\tau$ with the \emph{magnitude} bound
$\|R_\tau\|_{H^{-s}}\le C_K\tau(\|\Delta\|_{H^{-s}}+\Gamma_x^{\mathcal H})$ --- a paired
finite-difference bound that evaluates only the net reaction drift of
Lemma~\ref{lem:sens} at the two \emph{empirical} measures $\nu_\tau,\bar\nu_\tau$, and
never evaluates an exact expected reaction map at a non-empirical (possibly signed)
argument. Hence, by the magnitude decomposition \eqref{eq:Ttau-vardecomp} established
there (with $A:=\|\mathbb E_n\Delta\|_{H^{-s}}$, $B:=\mathbb E_n\|\Delta^c\|_{H^{-s}}^2$),
\begin{equation}\label{eq:Ttau-var}
\mathbb E_n\bigl\|T_\tau-\mathbb E_nT_\tau\bigr\|_{H^{-s}}^2
=\operatorname{Var}_n(T_\tau)
\le(1+C_K\tau)\,B+C_K\tau\,A^2+C_K\tau\,\mathbb E_n(\Gamma_x^{\mathcal H})^2 .
\end{equation}
The last two terms are \emph{standalone} (non-$1/N$) higher-order contributions: $A$
is the deterministic mean direction and does not fluctuate, and the displacement
variance $\mathbb E_n(\Gamma_x^{\mathcal H})^2$ is a pure $\tau$-fluctuation. We do
\emph{not} attempt to extract a further factor $1/N$ from either, as that would
require a Lipschitz/second-order sensitivity of the exact expected reaction map
$\mathbb E[\mathsf B^{N,\omega}_\tau\,\cdot\,]$, which we deliberately avoid; as shown
below both are globally higher order. By the Step~1--3 pre-branching bounds
$A\le C_K\tau^2$, $B\le C_K\tau^2/N$ and $\mathbb E_n(\Gamma_x^{\mathcal H})^2\le C_K\tau^3$,
\[
\mathbb E_n\|\mu_{\mathcal G}-\mathbb E_n\mu_{\mathcal G}\|_{H^{-s}}^2
\le(1+C_K\tau)\,C_K\frac{\tau^2}{N}+C_K\tau\,(C_K\tau^2)^2+C_K\tau\cdot C_K\tau^3
\le C_K\Bigl(\frac{\tau^2}{N}+\tau^4\Bigr)
\qquad\text{for every }N,
\]
the standalone terms being $C_K\tau A^2\le C_K\tau^5$ and
$C_K\tau\,\mathbb E_n(\Gamma_x^{\mathcal H})^2\le C_K\tau^4$, with $\tau^5\le\tau^4$.
This $\tau^4$ is a genuine $\tau$-fluctuation, carried as a separate contribution and
shown to be globally negligible in
Proposition~\ref{prop:time_discrete_particle} (it accumulates to $O(\tau^2)$ in the
mean-square error, dominated by the leading $O(\tau^2)$).
Adding the two, $\mathbb E_n\|\ell^{\mathrm{tr}}_{n+1}-\mathbb E_n\ell^{\mathrm{tr}}_{n+1}\|_{H^{-s}}^2
\le C_K(\tau^2/N+\tau^4)$. This proves \eqref{eq:transport}.
\end{proof}

\begin{lemma}[Branching discretization defect]\label{lem:branch}
Under Assumption~\ref{ass:time}, the branching defect
$\ell^{\mathrm{br}}_{n+1}=\mathsf B^{N,\omega}_\tau(\mathsf A^{N,\omega}_{\tau,h}(\mu^N_n))
-\mathsf B^{N,\omega}_{\tau,h}(\mathsf A^{N,\omega}_{\tau,h}(\mu^N_n))$ (the exact
pure reaction subflow against the frozen-rate Poisson branching, both starting
from the post-EM measure $\mu^{N,*}_n$) satisfies
\begin{equation}\label{eq:branch}
\|\mathbb E_n\ell^{\mathrm{br}}_{n+1}\|_{H^{-s}}\le C_K\,\tau^2,
\qquad
\mathbb E_n\|\ell^{\mathrm{br}}_{n+1}-\mathbb E_n\ell^{\mathrm{br}}_{n+1}\|_{H^{-s}}^2
\le C_K\,\Bigl(\frac{\tau^2}{N}+\tau^4\Bigr).
\end{equation}
\end{lemma}

\begin{proof}
In both flows the particle positions are fixed at the post-EM locations
$\mathbf X^*_i$ and the branching starts from the post-EM measure
$\eta_0:=\mu^{N,*}_n$. Crucially, $\eta_0$, the positions $\mathbf X^*_i$ and the
frozen rates $c_i:=c_{\eta_0}(\mathbf X^*_i)$ are produced by the EM transport step
and therefore depend on the step's Brownian increments; they are \emph{not}
$\mathcal F_n$-measurable. We argue conditionally on the transport
$\sigma$-algebra $\mathcal G_n$ of Lemma~\ref{lem:Bstab}, under which $\eta_0$,
$\mathbf X^*_i$ and $c_i$ are fixed and the only randomness left in either flow is
the genealogical branching driven by the independent Poisson random measures
$\{\mathcal N^\pm_\alpha\}$ of the graphical construction; we then tower back to
$\mathcal F_n$ using $\mathcal F_n\subseteq\mathcal G_n$. Write $c^+:=\max(c,0)$,
$c^-:=\max(-c,0)$, let $\eta_\theta$ be the running measure of the exact
pure-reaction subflow, and abbreviate
$\mathbb E_{\mathcal G}[\,\cdot\,]:=\mathbb E[\,\cdot\mid\mathcal G_n]$,
$\operatorname{Var}_{\mathcal G}:=\operatorname{Var}(\,\cdot\mid\mathcal G_n)$.

\emph{Intermediate frozen continuous-time reaction flow.} To separate the two error
mechanisms cleanly --- and to rule out any double counting between the algorithmic
surplus $e^{c_i\tau}-1-c_i\tau$ and the descendant/multiple-event remainder --- we
insert an intermediate flow $\mathsf B^{N,\omega}_{\tau,0}$ between the exact and the
implemented branching: it starts from $\eta_0:=\mu^{N,*}_n$ with positions fixed at
$\mathbf X^*_i$, and runs a continuous-time birth--death tree with the rates
\emph{frozen} at the root values $c_i:=c_{\eta_0}(\mathbf X^*_i)$ (so the per-root
birth/death intensities are the constants $c^\pm_i$ over $[0,\tau]$). Writing
$\eta_0$ for the common starting measure, we split
\[
\ell^{\mathrm{br}}_{n+1}
=\underbrace{\mathsf B^{N,\omega}_\tau(\eta_0)-\mathsf B^{N,\omega}_{\tau,0}(\eta_0)}_{=:\ \ell^{\mathrm{rate}}}
+\underbrace{\mathsf B^{N,\omega}_{\tau,0}(\eta_0)-\mathsf B^{N,\omega}_{\tau,h}(\eta_0)}_{=:\ \ell^{\mathrm{law}}}.
\]
Here $\ell^{\mathrm{rate}}$ compares the exact \emph{running} rate $c_{\eta_\theta}$ with
the frozen \emph{constant} rate $c_{\eta_0}$ on the same continuous-time tree (driven by
the same clocks), and is exactly the temporal contribution~(T) of Step~1; $\ell^{\mathrm{law}}$
compares the frozen-rate continuous-time tree with the one-shot frozen-rate algorithm
$\mathsf B^{N,\omega}_{\tau,h}$, both having the \emph{same} per-root first moment
$e^{c_i\tau}$ but different laws, and collects the algorithmic contribution~(A)
together with the descendant/multiple-event mismatch~(D). Because the two
comparisons are along disjoint legs of the telescope, (T) on the one hand and (A)+(D)
on the other are accounted for \emph{once each}; the directing-length split and the
surplus/descendant estimates of Steps~1--2 below are organized accordingly. We have
$\|\mathbb E_{\mathcal G}\ell^{\mathrm{rate}}\|_{H^{-s}}\le C_K\tau^2$ and
$\mathbb E_{\mathcal G}\|\ell^{\mathrm{rate}}-\mathbb E_{\mathcal G}\ell^{\mathrm{rate}}\|_{H^{-s}}^2\le C_K\tau^2/N$
from the temporal directing length (T) below, and likewise
$\|\mathbb E_{\mathcal G}\ell^{\mathrm{law}}\|_{H^{-s}}\le C_K\tau^2$,
$\mathbb E_{\mathcal G}\|\ell^{\mathrm{law}}-\mathbb E_{\mathcal G}\ell^{\mathrm{law}}\|_{H^{-s}}^2\le C_K\tau^2/N$
from (A)+(D); summing the two via $\|a+b\|^2\le2\|a\|^2+2\|b\|^2$ gives
\eqref{eq:branch}. The remaining steps establish these constituent bounds.

\emph{Step 1: directing length and its conditional moments.} On the shared clocks
the two branchings of label $i$ differ only through surplus events. For each
$\sigma\in\{+,-\}$ set the surplus \emph{directing length}
\[
L_{i,\sigma}:=\underbrace{\int_0^\tau\bigl|c^\sigma_{\eta_\theta}(\mathbf X^*_i)
-c^\sigma_{\eta_0}(\mathbf X^*_i)\bigr|\,\mathrm d\theta}_{\text{(T) temporal}}
+\underbrace{\mathbf 1_{\{c_i>0\}}\,\bigl|e^{c_i\tau}-1-c_i\tau\bigr|}_{\text{(A) algorithmic}},
\]
and let $R^{\mathrm{desc}}_i$ be the \emph{descendant/multiple-event remainder}: the
events on which the one-shot Poisson birth rule $1+\mathrm{Poisson}(e^{c_i\tau}-1)$
and the exact continuous-time pure-birth process disagree beyond their common
first moment.

\emph{(A)} is $\mathcal G_n$-measurable, and since $0\le c_i\le C_K$,
$\bigl|e^{c_i\tau}-1-c_i\tau\bigr|\le C_K\tau^2$.

\emph{(D)} The two birth laws share the mean $e^{c_i\tau}$ but differ in
distribution; they disagree only when root $i$ undergoes a second-generation birth
or any multiple event within the step, of $\mathcal G_n$-conditional probability
$O_K(\tau^2)$ because each event has predictable intensity $\le C_K$. Thus
$R^{\mathrm{desc}}_i$ is a nonnegative count with
\[
\mathbb E_{\mathcal G}R^{\mathrm{desc}}_i\le C_K\tau^2,
\qquad
\mathbb E_{\mathcal G}\bigl(R^{\mathrm{desc}}_i\bigr)^2\le C_K\tau^2,
\]
exactly as the multiple-event remainder $R_{\mathrm{rem}}$ in the proof of
Lemma~\ref{lem:split}.

\emph{(T)} cannot be controlled pathwise: $\eta_\theta$ is a c\`adl\`ag jump
empirical measure carrying a branching martingale part, so $\partial_\theta\eta_\theta$
does not exist along paths and the deterministic reaction-PDE estimate
``$\partial_\theta\eta_\theta=c_{\eta_\theta}\eta_\theta\Rightarrow
\|\partial_\theta\mathcal P_K\eta_\theta\|_{L^\infty}\le C_K$'' is false here.
Instead we use the semimartingale (Doob--Meyer) decomposition of the projected
measure under $\mathcal G_n$,
\[
\mathcal P_K(\eta_\theta-\eta_0)
=\int_0^\theta\mathcal P_K\bigl(c_{\eta_r}\eta_r\bigr)\,\mathrm dr+M^K_\theta,
\]
where $M^K$ is a $\mathcal G_n$-martingale whose jumps are $\frac1N\mathcal P_K\delta_x$,
each of size $\|\mathcal P_K\delta_x\|_{L^\infty}\le C_K$ (a projection onto
$\le(2K+1)^d$ Fourier modes). The compensator obeys, by $\|c\|_{L^\infty}\le C_K$
and the population bound of Assumption~\ref{ass:time}(B3)
($|\mathcal I(N,t)|\le C_NN$, hence $\eta_r(\mathbb T^d)\le C_N$),
\[
\Bigl\|\int_0^\theta\mathcal P_K\bigl(c_{\eta_r}\eta_r\bigr)\,\mathrm dr\Bigr\|_{L^\infty}
\le C_K\int_0^\theta\eta_r(\mathbb T^d)\,\mathrm dr\le C_KC_N\theta,
\]
a \emph{pathwise} bound that does not rely on any exponential mass estimate for the
random branching population. For the martingale part, the $\mathcal G_n$-expected
number of branching jumps up to $\theta$ is
$\le(\#\text{active labels})\cdot C_K\theta\le C_NN\cdot C_K\theta$ (again by (B3)),
so the expected total jump mass in $L^\infty$ is
$\le C_NN\cdot C_K\theta\cdot\frac{C_K}{N}=C_KC_N\theta$. Hence, absorbing the
fixed population constant $C_N$ into $C_K$,
\[
\mathbb E_{\mathcal G}\bigl\|\mathcal P_K(\eta_\theta-\eta_0)\bigr\|_{L^\infty}\le C_K\theta,
\qquad 0\le\theta\le\tau.
\]
By the Lipschitz continuity of $r$ and $\|\nabla c\|_{L^\infty}\le C_K$
(via \eqref{eq:Blip}), together with the $1$-Lipschitzness of $x\mapsto x_\pm$,
\[
\mathbb E_{\mathcal G}\bigl[\text{(T)}\bigr]
\le C_K\int_0^\tau\mathbb E_{\mathcal G}\bigl\|\mathcal P_K(\eta_\theta-\eta_0)\bigr\|_{L^\infty}\,\mathrm d\theta
\le C_K\tau^2 .
\]
This expected (not pathwise) bound is the correct replacement for the
intensity-gap estimate: pathwise one only has $\text{(T)}\le C_K\tau$. Combining
the three sources, $\mathbb E_{\mathcal G}L_{i,\sigma}\le C_K\tau^2$, while the
pathwise bound $L_{i,\sigma}\le C_K\tau$ gives the crude
$\mathbb E_{\mathcal G}L_{i,\sigma}^2\le C_K\tau^2$.

\emph{Step 2: surplus counts.} On the common Poisson random measure
$\mathcal N^\sigma_i$ the surplus count $S_{i,\sigma}$ is, conditionally on the
predictable intensities, a Cox (doubly stochastic Poisson) count with directing
length $L_{i,\sigma}$; by the law of total variance,
\[
\mathbb E_{\mathcal G}S_{i,\sigma}=\mathbb E_{\mathcal G}L_{i,\sigma}\le C_K\tau^2,
\qquad
\operatorname{Var}_{\mathcal G}(S_{i,\sigma})
=\mathbb E_{\mathcal G}L_{i,\sigma}+\operatorname{Var}_{\mathcal G}(L_{i,\sigma})
\le C_K\tau^2+\mathbb E_{\mathcal G}L_{i,\sigma}^2\le C_K\tau^2 .
\]
Each surplus or descendant event moves $\ell^{\mathrm{br}}_{n+1}$ by
$\pm\frac1N\delta_{\mathbf X^*_i}$.

\emph{Step 3: conditional mean.} Testing against $f\in H^s$ and summing over the
$\le C_NN$ active labels and the two signs,
\[
\bigl|\mathbb E_{\mathcal G}\langle f,\ell^{\mathrm{br}}_{n+1}\rangle\bigr|
\le\frac1N\sum_{i,\sigma}\mathbb E_{\mathcal G}S_{i,\sigma}\,\|f\|_{L^\infty}
+\frac1N\sum_i\mathbb E_{\mathcal G}R^{\mathrm{desc}}_i\,\|f\|_{L^\infty}
\le C_K\tau^2\|f\|_{L^\infty},
\]
so by $H^s\hookrightarrow L^\infty$,
$\|\mathbb E_{\mathcal G}\ell^{\mathrm{br}}_{n+1}\|_{H^{-s}}\le C_K\tau^2$. Since
$\mathbb E_n=\mathbb E_n\mathbb E_{\mathcal G}$ and the $H^{-s}$ norm is convex,
\[
\|\mathbb E_n\ell^{\mathrm{br}}_{n+1}\|_{H^{-s}}
\le\mathbb E_n\|\mathbb E_{\mathcal G}\ell^{\mathrm{br}}_{n+1}\|_{H^{-s}}\le C_K\tau^2,
\]
the first bound in \eqref{eq:branch}.

\emph{Step 4: centered part by martingale orthogonality.} Decompose the total
$\mathcal F_n$-variance by conditioning on $\mathcal G_n$ (the cross term vanishes,
its inner factor being $\mathcal G_n$-measurable):
\[
\mathbb E_n\bigl\|\ell^{\mathrm{br}}_{n+1}-\mathbb E_n\ell^{\mathrm{br}}_{n+1}\bigr\|_{H^{-s}}^2
=\underbrace{\mathbb E_n\,\mathbb E_{\mathcal G}\bigl\|\ell^{\mathrm{br}}_{n+1}
-\mathbb E_{\mathcal G}\ell^{\mathrm{br}}_{n+1}\bigr\|_{H^{-s}}^2}_{\text{within }\mathcal G_n}
+\underbrace{\mathbb E_n\bigl\|\mathbb E_{\mathcal G}\ell^{\mathrm{br}}_{n+1}
-\mathbb E_n\ell^{\mathrm{br}}_{n+1}\bigr\|_{H^{-s}}^2}_{(\star)\ \text{across }\mathcal G_n}.
\]
For the within-$\mathcal G_n$ term, write the $\mathcal G_n$-centered defect as
compensated surplus-jump integrals against the graphical Poisson measures plus the
centered descendant remainder,
\[
\ell^{\mathrm{br}}_{n+1}-\mathbb E_{\mathcal G}\ell^{\mathrm{br}}_{n+1}
=\frac1N\sum_{i,\sigma}\int_{[0,\tau]\times[0,\infty)}\delta_{\mathbf X^*_i}\,
\widetilde{\mathcal N}_{i,\sigma}(\mathrm d\theta,\mathrm du)
+\bigl(R^{\mathrm{desc}}-\mathbb E_{\mathcal G}R^{\mathrm{desc}}\bigr),
\]
where $\widetilde{\mathcal N}_{i,\sigma}$ is the compensated surplus point measure
carved out of $\mathcal N^\sigma_i$ by the matched/surplus thinning. Because the
running rate $c_{\eta_\theta}$ is a \emph{single global random field} shared by all
labels, the surplus counts are \emph{not} conditionally independent across
particles; we therefore do not invoke independence but the orthogonality of
stochastic integrals against the \emph{independent} Poisson random measures
$\{\mathcal N^\sigma_\alpha\}$ of the graphical construction: for
$(i,\sigma)\ne(j,\sigma')$ the compensated integrals are orthogonal in
$L^2(\mathbb P\mid\mathcal G_n)$ by the martingale isometry, irrespective of any
correlation between their predictable intensities. The centered descendant
remainder is treated identically and \emph{not} by cross-particle independence:
$R^{\mathrm{desc}}-\mathbb E_{\mathcal G}R^{\mathrm{desc}}$ is a sum of compensated
stochastic integrals over \emph{disjoint} Ulam--Harris subtrees, each rooted at a
distinct label and driven by its own independent Poisson families, plus a
predictable remainder of order $O_K(\tau^2/N)$; the same martingale orthogonality
over these disjoint genealogical families validates the labelwise variance sum
below. Hence, with $\|\delta_{\mathbf X^*_i}\|_{H^{-s}}\le C_s$ ($s>d/2$),
\[
\mathbb E_{\mathcal G}\bigl\|\ell^{\mathrm{br}}_{n+1}-\mathbb E_{\mathcal G}\ell^{\mathrm{br}}_{n+1}\bigr\|_{H^{-s}}^2
\le\frac{C_s}{N^2}\sum_{i,\sigma}\operatorname{Var}_{\mathcal G}(S_{i,\sigma})
+\frac{C_s}{N^2}\sum_i\operatorname{Var}_{\mathcal G}(R^{\mathrm{desc}}_i)
\le C_K\frac{\tau^2}{N},
\]
by Step~2 and the population bound; taking $\mathbb E_n$ preserves this.

For $(\star)$ we do \emph{not} use any Lipschitz estimate for the exact expected
reaction map $\rho\mapsto\mathbb E[\mathsf B^{N,\omega}_\tau\rho]$, nor any
independent-copy/second-divided-difference argument for a difference of expected
maps. We use only the $\mathcal G_n$-conditional mean bound already established in
Step~3,
\[
\|\mathbb E_{\mathcal G}\ell^{\mathrm{br}}_{n+1}\|_{H^{-s}}\le C_K\tau^2 ,
\]
together with the elementary fact that a conditional variance is dominated by the
corresponding conditional second moment. Since $\mathbb E_{\mathcal G}\ell^{\mathrm{br}}_{n+1}$
is $\mathcal G_n$-measurable and $\mathbb E_n\ell^{\mathrm{br}}_{n+1}
=\mathbb E_n[\mathbb E_{\mathcal G}\ell^{\mathrm{br}}_{n+1}]$,
\[
(\star)
:=\mathbb E_n\bigl\|\mathbb E_{\mathcal G}\ell^{\mathrm{br}}_{n+1}-\mathbb E_n\ell^{\mathrm{br}}_{n+1}\bigr\|_{H^{-s}}^2
=\operatorname{Var}_n\bigl(\mathbb E_{\mathcal G}\ell^{\mathrm{br}}_{n+1}\bigr)
\le\mathbb E_n\bigl\|\mathbb E_{\mathcal G}\ell^{\mathrm{br}}_{n+1}\bigr\|_{H^{-s}}^2
\le C_K\tau^4 .
\]
This bound is deliberately \emph{not} divided by $N$. It avoids the (unproved, and
under the hard split $c^\pm=\max(\pm c,0)$ delicate) $C_K\tau^2$-Lipschitz property of
the exact expected reaction map, relying instead only on the safe frozen-map and
net-drift estimates of Lemma~\ref{lem:sens} that already underlie the Step~3 mean
bound. The resulting standalone $O(\tau^4)$ per-step variance is harmless in the
global accumulation: after the Young weight $\tau^{-1}$ in Lemma~\ref{lem:accum} and
summation over $O(\tau^{-1})$ steps it contributes $O(\tau^2)$ to the final
mean-square error. Adding the within-$\mathcal G_n$ term and $(\star)$,
\[
\mathbb E_n\|\ell^{\mathrm{br}}_{n+1}-\mathbb E_n\ell^{\mathrm{br}}_{n+1}\|_{H^{-s}}^2
\le C_K\Bigl(\frac{\tau^2}{N}+\tau^4\Bigr),
\]
which is the second bound in \eqref{eq:branch}.
\end{proof}

Before assembling these into the global recursion, we summarize the per-step local
defect bounds established above. Throughout, $\zeta_n:=\mu^N(t_n)-\mu^N_n$ denotes the
global error and $E_n:=\mathbb E\|\zeta_n\|_{H^{-s}}^2$ its mean square, with $\zeta_0=0$
and $E_0=0$ since the scheme is initialized at the exact state ($\mu^N(0)=\mu^N_0$).
\begin{center}
\renewcommand{\arraystretch}{1.4}
\begin{tabular}{lccc}
\hline
defect & $\ell^{\mathrm{split}}_{n+1}$ & $\ell^{\mathrm{tr}}_{n+1}$ & $\ell^{\mathrm{br}}_{n+1}$\\
\hline
$\|\mathbb E_n\ell\|_{H^{-s}}$ & $C_K\tau^2$ & $C_K\tau^2$ & $C_K\tau^2$\\
$\mathbb E_n\|\ell-\mathbb E_n\ell\|_{H^{-s}}^2$ & $C_K(\tau^{3/2}/N+\tau^4)$ & $C_K(\tau^2/N+\tau^4)$ & $C_K(\tau^2/N+\tau^4)$\\
\hline
\end{tabular}
\end{center}
The means are all $O(C_K\tau^2)$, so after accumulation over $O(1/\tau)$ steps they
contribute $O(C_K\tau)$ to the root-mean-square error. The fluctuations have a leading
$1/N$ part together with a standalone higher-order $O(C_K\tau^4)$ part, with two
features. (i) The leading $1/N$ fluctuation of the Lie-splitting term is
$O(C_K\tau^{3/2}/N)$: the Brownian displacement of the transported particles creates a
surplus branching interval of intensity-length $O(\tau^{3/2})$ rather than $O(\tau^2)$
(Lemma~\ref{lem:split} and Remark~\ref{rem:K-scaling}), whereas the transport and
branching $1/N$ fluctuations are $O(C_K\tau^2/N)$. This single $\tau^{3/2}$ is
why the global fluctuation is $\tau^{-1/2}/N$ rather than $1/N$, and hence why the
final bound is $\tau^2+\tau^{-1/2}/N$ and not $\tau^2+1/N$. (ii) All three defects also
carry a \emph{standalone} $O(C_K\tau^4)$ fluctuation (no factor $1/N$), introduced by
the safe treatment of the reaction mean maps and of the multiplicity correction: in
Lemma~\ref{lem:split} it is the safe non-orthogonal drift-feedback variance of
Lemma~\ref{lem:drift-feedback-generator} together with the across-field
multiplicity variance of the multiple-event remainder $R_{\mathrm{rem}}$; in Lemma~\ref{lem:transport} it is the
displacement-variance term $C_K\tau\,\mathbb E_n(\Gamma_x^{\mathcal H})^2\le C_K\tau^4$
of the paired reaction estimate (Lemma~\ref{lem:Bstab}); and in
Lemma~\ref{lem:branch} it is the across-$\mathcal G_n$ variance of the conditional
mean, bounded crudely by $\|\mathbb E_{\mathcal G}\ell^{\mathrm{br}}\|_{H^{-s}}^2\le C_K\tau^4$
rather than by an unproved Lipschitz estimate for the exact expected reaction map.
This standalone term is globally negligible: a per-step centered defect of variance
$O(\tau^4)$ is multiplied by the $O(\tau^{-1})$ Young weight in the
recursion--accumulation, giving $O(\tau^3)$ per step, and after summation over
$O(\tau^{-1})$ steps it contributes $O(\tau^2)$ to $\mathbb E\|\zeta_n\|_{H^{-s}}^2$ ---
of the same order as the leading deterministic time error and therefore absorbed into
it, with \emph{no} constraint on $N$.

\begin{lemma}[Cross terms and accumulation]\label{lem:accum}
Under Assumption~\ref{ass:time}, the iterates $E_n=\mathbb E\|\zeta_n\|_{H^{-s}}^2$
satisfy, for $\tau\le\tau_0$,
\[
E_{n+1}\le(1+C_K'\tau)E_n+C_K\,\tau^3+C_K\,\frac{\tau^{1/2}}{N},
\]
and consequently, for all $t_n\le T$,
\[
E_n\le e^{C_K'T}E_0+C_{T,K}\Bigl(\tau^2+\frac{\tau^{-1/2}}{N}\Bigr).
\]
\end{lemma}

\begin{proof}
Write $A_{n+1}:=\Phi^{N,\omega}_\tau(\mu^N(t_n))-\Phi^{N,\omega}_\tau(\mu^N_n)$ for
the difference of the exact flow applied to the two states at $t_n$, so that
$\zeta_{n+1}=A_{n+1}+\ell_{n+1}$ as a pathwise identity, the inserted
$\Phi^{N,\omega}_\tau(\mu^N_n)$ being a single realization under the noise $\omega$
common to $A_{n+1}$ and $\ell_{n+1}$. Applying Lemma~\ref{lem:stability} with
$\mu=\mu^N(t_n)$, $\mu'=\mu^N_n$ gives
$\mathbb E_n\|A_{n+1}\|_{H^{-s}}^2\le(1+C_K\tau)\|\zeta_n\|_{H^{-s}}^2+C_K\tau/N$.
Collect the local defect $\ell_{n+1}=\ell^{\mathrm{split}}_{n+1}
+\ell^{\mathrm{tr}}_{n+1}+\ell^{\mathrm{br}}_{n+1}$
and write its mean and centered parts $\bar\ell_{n+1}:=\mathbb E_n\ell_{n+1}
=\sum_\alpha\bar\ell^\alpha_{n+1}$, $m_{n+1}:=\ell_{n+1}-\bar\ell_{n+1}
=\sum_\alpha m^\alpha_{n+1}$, with $\alpha\in\{\mathrm{split},\mathrm{tr},\mathrm{br}\}$.
By the triangle inequality on $\bar\ell_{n+1}$ and $\|\sum_\alpha m^\alpha\|^2
\le3\sum_\alpha\|m^\alpha\|^2$ on the centered parts (used in place of any
orthogonality, since the three centered defects all depend on the
step-$(n{+}1)$ noise), together with the per-defect bounds from
Lemmas~\ref{lem:split}, \ref{lem:transport}, and \ref{lem:branch} (among the
$1/N$ parts the centered bound $C_K\tau^{3/2}/N$ of Lemma~\ref{lem:split} dominates the
$C_K\tau^2/N$ of the other two; all three defects additionally carry a standalone
$C_K\tau^4$, from the drift-feedback and multiplicity variance in Lemma~\ref{lem:split} (the
safe non-orthogonal estimate of Lemma~\ref{lem:drift-feedback-generator}) and from
the safe reaction-map treatment in
Lemmas~\ref{lem:transport}--\ref{lem:branch}),
\[
\|\bar\ell_{n+1}\|_{H^{-s}}\le C_K\,\tau^2,
\qquad
\mathbb E_n\|m_{n+1}\|_{H^{-s}}^2\le C_K\,\frac{\tau^{3/2}}{N}+C_K\,\tau^4.
\]
From $\zeta_{n+1}=A_{n+1}+\bar\ell_{n+1}+m_{n+1}$,
\[
\|\zeta_{n+1}\|_{H^{-s}}^2
=\|A_{n+1}\|_{H^{-s}}^2+\|\bar\ell_{n+1}+m_{n+1}\|_{H^{-s}}^2
+2\langle A_{n+1},\bar\ell_{n+1}\rangle+2\langle A_{n+1},m_{n+1}\rangle.
\]
By the bounds above,
$\mathbb E_n\|\bar\ell_{n+1}+m_{n+1}\|_{H^{-s}}^2\le 2\|\bar\ell_{n+1}\|^2
+2\mathbb E_n\|m_{n+1}\|^2\le C_K\tau^4+C_K\tau^{3/2}/N$ (the standalone
$C_K\tau^4$ from the centered splitting/transport/branching defects). For the drift cross term,
\[
2|\mathbb E_n\langle A_{n+1},\bar\ell_{n+1}\rangle|
\le\tau\mathbb E_n\|A_{n+1}\|_{H^{-s}}^2+\tau^{-1}\|\bar\ell_{n+1}\|_{H^{-s}}^2
\le\tau\mathbb E_n\|A_{n+1}\|_{H^{-s}}^2+C_K\tau^3,
\]
by Cauchy--Schwarz and Young with weight $\tau$. For the martingale cross term we
do not rely on orthogonality (both $A_{n+1}$ and $m_{n+1}$ depend on the
step-$(n{+}1)$ noise, so their inner product need not vanish in $\mathbb E_n$) and
again use Young with weight $\tau$,
\[
2|\mathbb E_n\langle A_{n+1},m_{n+1}\rangle|
\le\tau\mathbb E_n\|A_{n+1}\|_{H^{-s}}^2+\tau^{-1}\mathbb E_n\|m_{n+1}\|_{H^{-s}}^2
\le\tau\mathbb E_n\|A_{n+1}\|_{H^{-s}}^2+C_K\frac{\tau^{1/2}}{N}+C_K\tau^3,
\]
the last step using $\tau^{-1}(C_K\tau^{3/2}/N+C_K\tau^4)=C_K\tau^{1/2}/N+C_K\tau^3$;
the standalone $\tau^3$ is the global footprint of the splitting, transport, and
branching $\tau^4$ fluctuations, of the same order as the $C_K\tau^3$ already present
from the drift cross term and absorbed below.
Inserting these together with the stability bound \eqref{eq:stab}, using
$(1+2\tau)(1+C_K\tau)\le1+C_K'\tau$ for $\tau\le\tau_0$ and absorbing lower-order
terms via $\tau\le1$ ($\tau^4\le\tau^3$, $C_K\tau^{3/2}/N\le C_K\tau^{1/2}/N$,
$\tau/N\le\tau^{1/2}/N$), and taking total expectation,
\[
E_{n+1}\le(1+C_K'\tau)E_n+C_K\,\tau^3+C_K\,\frac{\tau^{1/2}}{N}.
\]
Iterating from a general initial value $E_0$ with $n\tau\le T$,
$(1+C_K'\tau)^n\le e^{C_K'T}$, and $\sum_{j=0}^{n-1}(1+C_K'\tau)^j
\le(e^{C_K'T}-1)/(C_K'\tau)$,
\[
E_n
\le(1+C_K'\tau)^nE_0+\Bigl(C_K\tau^3+C_K\frac{\tau^{1/2}}{N}\Bigr)\frac{e^{C_K'T}-1}{C_K'\tau}
\le e^{C_K'T}E_0+\frac{e^{C_K'T}-1}{C_K'}\Bigl(C_K\tau^2+\frac{C_K\tau^{-1/2}}{N}\Bigr)
\le e^{C_K'T}E_0+C_{T,K}\Bigl(\tau^2+\frac{\tau^{-1/2}}{N}\Bigr),
\]
the fixed-$K$ factors being absorbed into $C_{T,K}$. The exponent $-1/2$ in the
fluctuation term originates from the centered Lie-splitting bound of
Lemma~\ref{lem:split}.

\emph{Stopped form.} The local-defect bounds entering the recursion are the
\emph{stopped} estimates of Lemmas~\ref{lem:split}--\ref{lem:branch}, valid on
$\{t_n\le T\wedge\Theta_{K,R,N}\}$ (Remark~\ref{rem:admissible}). Defining the stopped
iterate
\[
E_n^\Theta:=\mathbb E\bigl\|\mu^N(t_{n\wedge n_\Theta})-\mu^N_{n\wedge n_\Theta}\bigr\|_{H^{-s}}^2,
\qquad n_\Theta:=\lfloor\Theta_{K,R,N}/\tau\rfloor,
\]
the same computation gives the recursion and its solution verbatim for $E_n^\Theta$,
\[
E_{n+1}^\Theta\le(1+C_K'\tau)E_n^\Theta+C_K\tau^3+C_K\frac{\tau^{1/2}}{N},
\qquad
E_n^\Theta\le e^{C_K'T}E_0^\Theta+C_{T,K}\Bigl(\tau^2+\frac{\tau^{-1/2}}{N}\Bigr),
\]
since beyond $n_\Theta$ the increment $\zeta_{n+1}-\zeta_n$ is frozen and the
truncated coefficients keep all per-step bounds in force. The stopping is removed
only under the auxiliary-admissibility condition $\Theta_{K,R,N}>T$ a.s.\
(Remark~\ref{rem:admissible}), in which case
$n\wedge n_\Theta=n$ and $E_n^\Theta=E_n$ on $[0,T]$.
\end{proof}

\begin{proof}[Proof of Proposition~\ref{prop:time_discrete_particle}]
Recall $\zeta_n=\mu^N(t_n)-\mu^N_n$, $E_n=\mathbb E\|\zeta_n\|_{H^{-s}}^2$, with
$\zeta_0=0$, $E_0=0$, and the exact full flow $\Phi^{N,\omega}_\tau$ and the
implemented algorithm $\Psi^{N,\omega}_\tau$ from above. Inserting the exact flow
from the discrete state,
\[
\zeta_{n+1}
=\underbrace{\Phi^{N,\omega}_\tau(\mu^N(t_n))-\Phi^{N,\omega}_\tau(\mu^N_n)}_{A_{n+1}}
+\underbrace{\Phi^{N,\omega}_\tau(\mu^N_n)-\Psi^{N,\omega}_\tau(\mu^N_n)}_{\ell_{n+1}},
\qquad
\ell_{n+1}=\ell^{\mathrm{split}}_{n+1}+\ell^{\mathrm{tr}}_{n+1}+\ell^{\mathrm{br}}_{n+1},
\]
where the inserted middle term $\Phi^{N,\omega}_\tau(\mu^N_n)$ is one and the same
realization, driven by the same noise $\omega$ (the Brownian motions and Poisson
clocks of the graphical construction) in both $A_{n+1}$ and $\ell_{n+1}$; the
decomposition is therefore a \emph{pathwise} telescoping identity, not a difference
of two independent copies of the flow. The local defect splits into a Lie-splitting
part (full simultaneous flow against the exact nonlinear Lie split), a transport
part (exact transport against EM transport), and a branching part (exact pure
reaction against frozen-rate branching). Applying Lemma~\ref{lem:stability} with $\mu=\mu^N(t_n)$ and
$\mu'=\mu^N_n$ bounds $A_{n+1}$ by
$\mathbb E_n\|A_{n+1}\|_{H^{-s}}^2\le(1+C_K\tau)\|\zeta_n\|_{H^{-s}}^2+C_K\tau/N$;
Lemmas~\ref{lem:split}, \ref{lem:transport}, and \ref{lem:branch} bound the mean
and centered parts of the three local defects; and Lemma~\ref{lem:accum} combines
them to give the stopped recursion solution
$E_n^\Theta\le e^{C_K'T}E_0^\Theta+C_{T,K}(\tau^2+\tau^{-1/2}/N)$ on
$\{t_n\le T\wedge\Theta_{K,R,N}\}$; run with the stopped
local defects this is the stopped bound \eqref{eq:P1-stopped}. Since the scheme is
initialized at the exact state, $\mu^N(0)=\mu^N_0$, we have $\zeta_0=0$ and hence
$E_0^\Theta=0$. If the auxiliary-admissibility condition $\Theta_{K,R,N}>T$ holds
a.s.\ (Remark~\ref{rem:admissible}),
then $n\wedge n_\Theta=n$ for all $t_n\le T$, $E_n^\Theta=E_n$, and the
stopped bound \eqref{eq:P1-stopped} becomes the unstopped \eqref{eq:P1}; otherwise
the conclusion is read in the stopped sense. Taking square roots of
\eqref{eq:P1-stopped} and using
$\sqrt{a+b}\le\sqrt a+\sqrt b$ gives the stopped \eqref{eq:P2} (and, under
$\Theta_{K,R,N}>T$, its unstopped form). Finally, if $N\tau^{5/2}\ge c_*$
then $N^{-1/2}\tau^{-1/4}\le c_*^{-1/2}\tau$, whence the last assertion.
\end{proof}

\begin{corollary}[Fully discrete end-to-end error]\label{cor:end-to-end}
Under the hypotheses of Propositions~\ref{prop:u_K_mu},~\ref{prop:u_u_K}
and~\ref{prop:time_discrete_particle}, suppose the scheme is initialized at an
unbiased sample of $u_0$ with $u_K(0)=u(0)$ and
$\mathbb E\|u(0)-\mu^N(0)\|_{H^{-s}}^2\le C_s/N$, and assume in addition the
auxiliary-admissibility condition $\Theta_{K,R,N}>T$ a.s.\
(Remark~\ref{rem:admissible}), so that the
unstopped form \eqref{eq:P1} of
Proposition~\ref{prop:time_discrete_particle} applies. Then for every $t_n=n\tau\le T$ and
$0<\tau\le\tau_0$, using the embedding $\|v\|_{H^{-s}}\le\|v\|_{L^2}$ to pass the
coefficient-regularization error from $L^2$ (Proposition~\ref{prop:u_u_K}) to $H^{-s}$,
\begin{equation}\label{eq:end-to-end}
\mathbb E\bigl\|u(t_n)-\mu^N_n\bigr\|_{H^{-s}}^2
\le C_T\,K^{-2m}\|u\|_{L^2(0,t_n;\mathbb H^m)}^2
+C_{T,K}\Bigl(\frac1N+\tau^2+\frac{\tau^{-1/2}}{N}\Bigr),
\end{equation}
where the coefficient-regularization constant $C_T$ is independent of $K,N,\tau$
except through the displayed $K^{-2m}$ factor (it depends on $T$, the coefficients,
and the regularity size $\|u\|_{L^\infty(0,T;\mathbb H^m)}$ of the exact solution, but
not on the particle-discretization constants), whereas $C_{T,K}$ carries the fixed-$K$
factors. Since $\tau\le1$ gives
$1/N\le\tau^{-1/2}/N$, the bound simplifies to
\[
\mathbb E\bigl\|u(t_n)-\mu^N_n\bigr\|_{H^{-s}}^2
\le C_T\,K^{-2m}\|u\|_{L^2(0,t_n;\mathbb H^m)}^2
+C_{T,K}\Bigl(\tau^2+\frac{\tau^{-1/2}}{N}\Bigr).
\]
Without the auxiliary-admissibility condition, the same bounds hold in the stopped
sense, with $u(t_n),\mu^N_n$ replaced by their stopped versions
$u(t_{n\wedge n_\Theta}),\mu^N_{n\wedge n_\Theta}$, using \eqref{eq:P1-stopped} in
place of \eqref{eq:P1}.
\end{corollary}

\begin{proof}
Decompose $u(t_n)-\mu^N_n=\bigl(u(t_n)-\mu^N(t_n)\bigr)+\bigl(\mu^N(t_n)-\mu^N_n\bigr)$
and apply $\|a+b\|_{H^{-s}}^2\le2\|a\|_{H^{-s}}^2+2\|b\|_{H^{-s}}^2$. For the first
term, a further triangle split gives
\[
\|u(t_n)-\mu^N(t_n)\|_{H^{-s}}\le\|u(t_n)-u_K(t_n)\|_{H^{-s}}+\|u_K(t_n)-\mu^N(t_n)\|_{H^{-s}}:
\]
the consistency part is controlled by Proposition~\ref{prop:u_u_K}, stated in
$L^2$, via $\|u-u_K\|_{H^{-s}}\le\|u-u_K\|_{L^2}$, and contributes
$C_T\,K^{-2m}\|u\|_{L^2(0,t_n;\mathbb H^m)}^2$ with $C_T$ independent of $K$; the
particle part is bounded directly in $H^{-s}$ by Proposition~\ref{prop:u_K_mu},
contributing $C_{T,K}\,C_s/N$. The second term is bounded by \eqref{eq:P1} of
Proposition~\ref{prop:time_discrete_particle},
$C_{T,K}\bigl(\tau^2+\tau^{-1/2}/N\bigr)$. Keeping the $K$-independent
coefficient-regularization constant $C_T$ separate from the fixed-$K$ constant
$C_{T,K}$ (and absorbing $C_s$ into the latter) gives \eqref{eq:end-to-end}; the
second form follows from $1/N\le\tau^{-1/2}/N$ for $0<\tau\le1$.
\end{proof}

\begin{remark}[Explicit $K$-dependence]\label{rem:K-scaling}
The constant $C_{T,K}$ depends on the cutoff $K$, held fixed throughout (as in the
numerical experiments, where $K$ is a fixed spectral parameter). Its
$K$-dependence enters through two channels. First, the per-step local defect: the
transport and branching defects of Lemmas~\ref{lem:transport}
and~\ref{lem:branch} are $O(C_K\tau^2)$ in the mean and have the usual
$O(C_K\tau^2/N)$ stochastic part in variance, plus a standalone higher-order
$O(C_K\tau^4)$ variance contribution (no factor $1/N$) introduced by the safe
treatment of the reaction mean maps; this contribution accumulates only to
$O(\tau^2)$ in the global mean-square error and is absorbed by the deterministic
time-discretization error. The Lie-splitting defect of Lemma~\ref{lem:split} is
$O(C_K\tau^2)$ in the mean and has a leading $O(C_K\tau^{3/2}/N)$ stochastic part (plus
the same standalone $O(C_K\tau^4)$ contribution, here from the across-field
multiplicity variance); this leading
fluctuation, stemming from the Brownian scale of the surplus branching interval,
is the dominant stochastic term and is what produces the $\tau^{-1/2}/N$ term in \eqref{eq:P1}. The $K$-growth of the transport defect originates solely from the
order-two generator expansion, which involves $b$ only through $\nabla b,\nabla^2b
\sim K,K^2$ (and never the full $\|b\|_{C^m}$); this is the mechanism behind the
sharp scaling and explains why a matching condition of the form $K^2\tau\le c_0$,
rather than the pessimistic $K^8\tau\le c_0$ suggested by a black-box weak-order
constant $(1+\|b\|_{C^4})^2$, suffices in practice. Second, the stability constant:
the Gr\"onwall exponent in Lemmas~\ref{lem:stability} and~\ref{lem:accum} carries a
constant $C_K$ that includes the higher power $K^{(m-1)+s+d/2}$ from the Lipschitz
coefficient estimate \eqref{eq:Blip}, which enters only through the fixed-$K$
factor $e^{C_K T}$ and does not affect the per-step behavior. At fixed $K$ all
these factors are absorbed into $C_{T,K}$ and \eqref{eq:P1}--\eqref{eq:P2} hold as
stated.
\end{remark}

The numerical details to implement this step is summarized in Algorithm \ref{alg:linear}.

\begin{algorithm}
    \caption{Stochastic Branching for linear scalar reaction--diffusion--advection equation} \label{alg:linear}
    \begin{algorithmic}
        \Require Initial condition of PDE $u_0(\bx)$, domain descriptor $\Omega$;
        \Require Fixed initial population size $N$
        \Require Time step $\tau>0$, horizon $T$; truncation $\mathcal M_K$; basis $\{\psi_m\}_{m\in\mathcal M_K}$
        \Require Driven term $b(\bx)$
        \Require Diffusion strength $\sigma(\bx)$
        \Require Reaction maps $r(\bx)$
        \Ensure  Solution set of PDE $\mathcal U = \{u_n\}_{n=0}^{T/\tau}$
        \Procedure{Stochastic\_Branching}{$u_0$,$N$,$\tau$,$\mathcal M_K$,$\{\psi_m\}$,$b$,$\sigma$,$r$}
        \State \textbf{Initialization}
        \State $n\gets 0$
        \State Precomputed normalization constants  $Z\gets \int_\Omega u_0(\bx)\intd \bx$  (either numerically or analytically).
        \State Generate $N$ initial particle positions set $S_0=\{\bX_{i,0}\}$ according to the distributions $\rho_0=\frac{u_0}{Z}$ and using the Metropolis--Hastings algorithm.
        \State \textbf{Interpolate (fixed-$N$):}
        \State $\tilde\rho_0 \gets \textsc{Interp}(S_0,N,\mathcal M_K,\{\psi_m\})$
        \State Append $\tilde u_0 = Z\tilde \rho_0$ into $\mathcal U$
        \While {$n\tau<T$}
            \State \textbf{SDE propagation:}
            \State $S_n^*\gets \textsc{SDE}(S_n,\tau,b,\sigma)$
            \State \textbf{Birth--Death:}
            \State $S_{n+1}\gets \textsc{Birth\_Death}(S_{n}^*,\tau,r)$
            \State \textbf{refresh fields:}
            \State $\tilde\rho_{n+1} \gets \textsc{Interpolate}(S_{n+1},N,\mathcal M_K,\{\psi_m\})$
            \State Append $\tilde u_{n+1} = Z\tilde \rho_{n+1}$ into $\mathcal U$
            \State $n\gets n+1$
        \EndWhile
        \State \Return $\mathcal U$
        \EndProcedure
    \end{algorithmic}
\end{algorithm}

\begin{remark}
Propositions \ref{prop:u_K_mu} and \ref{prop:u_u_K} present the numerical convergence analysis for a scalar ADR equation. Extension to the two-variable KS system requires additional regularity assumptions of  the chemotactic drift, at least up to a stopping time prior to blow-up. A full two-field analysis is beyond the scope of this work and is left for future study.    
\end{remark}


\section{Extension to PDE system}
\label{sec:extension}
To solve PDE systems like Equation \eqref{equ:KS}, we introduce two sets of particles $\bX_i(t)$ and $\bY_i(t)$. Therefore, empirical distribution $\rho_u(t,\intd \bx ) = \mu_u(t,\intd \bx) =  \frac{1}{N_{u,0}}\sum_{i=1}^{N_{u,n}}\delta_{\bX_n^i} $ and $\rho_v(t,\intd \bx ) = \mu_v(t,\intd \bx) =  \frac{1}{N_{v,0}}\sum_{i=1}^{N_{v,n}}\delta_{\bY_n^i}$ are used to approximate $u$ and $v$ separately, by 
\begin{equation}
    u(t,\bx) =  Z_u \rho_u(t,\bx), \quad v(t,\bx) =  Z_v \rho_v(t,\bx),
\end{equation}
where $Z_u = \int_\Omega u(0,\bx) \intd \bx$ and $Z_v = \int_\Omega v(0,\bx) \intd \bx$. Therefore, the dynamics of \eqref{equ:KS} can be written as 
% \begin{equation}
% \begin{aligned}
%     \partial_t u - \nabla \cdot ( \nabla u - \chi(u) \nabla v) &= f_u(u,v),  & x\in \mathbb R^2 ,t\in(0,T),\\
%     v_t -\Delta v & = f_v(u,v),  &x\in \mathbb R^2 , t\in(0,T),
%     \end{aligned}
% \end{equation}
\begin{equation}
    \begin{aligned}
        \partial_t \rho_u =  \mathcal A_u \rho_u + \mathcal B_u \rho_u,\\
        \partial_t \rho_v =  \mathcal A_v \rho_v + \mathcal B_v \rho_v
    \end{aligned}
\end{equation}
where $\mathcal A_u \rho = \nabla \cdot ( \nabla \rho - \frac{Z_v}{Z_u} \chi(Z_u\rho) \nabla \rho_v)$ , $\mathcal A_v \rho = \Delta \rho$, $\mathcal B_u \rho_u  = f_u(Z_u\rho_u,Z_v \rho_v)$ and $B_v \rho_v = f_v(Z_u\rho_u,Z_v \rho_v)$. Similarly, we can approximate the solution by operator splitting, i.e., 
\begin{equation}
    \begin{aligned}
        \rho_{u,n}^*  =  \exp(\mathcal A_u\tau) \rho_{u,n}, \quad 
         \rho_{v,n}^*  =  \exp(\mathcal A_v\tau) \rho_{v,n},  \\
         \rho_{u,n+1} =  \exp(\mathcal B_u\tau) \rho_{u,n}^*,  \quad 
         \rho_{v,n+1}  =  \exp(\mathcal B_v\tau) \rho_{v,n}^*,   \\
    \end{aligned}
\end{equation}
where by matching parameters, we have $a_u(\bx;\rho_u,\rho_v) =  - \frac{Z_v}{Z_u\rho_u} \chi(Z_u\rho_u) \nabla \rho_v $ and $a_v(\bx;\rho_u,\rho_v)= 0$. For the diffusion part, we have $\sigma_u(\bx;\rho_u) = \sigma_v(\bx;\rho_v) = 1$. For birth and death part, $\tilde r_u (\bx;\rho_u,\rho_v)=  \frac{f_u(Z_u\rho_u,Z_v \rho_v)}{\rho_u}  $ and $\tilde r_v (\bx;\rho_u,\rho_v )=  \frac{f_v(Z_u\rho_u,Z_v \rho_v)}{\rho_v}  $. The detailed algorithm is shown in  Algorithm \ref{alg:ks}.

\begin{algorithm}
\caption{Stochastic Branching for keller--Segel equation}\label{alg:ks}
\begin{algorithmic}
\Require Initial PDE data $u_0(\bx),\,v_0(\bx)$; domain descriptor $\Omega$
\Require Fixed initial population sizes $N_u, N_v$
\Require Time step $\tau>0$, horizon $T$; truncation $\mathcal M_K$; basis $\{\psi_m\}_{m\in\mathcal M_K}$
% \Require Diffusion strengths $\sigma_u(\bx),\,\sigma_v(\bx)$ \hfill (constant or space-dependent)
% \Require Chemotactic coefficient: $\chi$
\Require Normalization constants $Z_u=\int_\Omega u_0\,\mathrm d\bx$, $Z_v=\int_\Omega v_0\,\mathrm d\bx$
% \Require Reaction maps $f_u,f_v:\mathbb R_{\ge 0}\times\mathbb R_{\ge 0}\to\mathbb R$
\Ensure Discrete KS fields $\mathcal U=\{\tilde u_n\}_{n=0}^{T/\tau}$, $\mathcal V=\{\tilde v_n\}_{n=0}^{T/\tau}$
% \Procedure{KS\_Stochastic\_Branching}{$u_0,v_0,N_u,N_v,\tau,\mathcal M_K,\newline \phantom{\quad\quad \textbf{procesure} \texttt{KS\_Stochastic\_Branching}(}\{\psi_m\},\sigma_u,\sigma_v,\chi,f_u,f_v$}
\Procedure{KS\_Stochastic\_Branching}{$u_0,v_0,N_u,N_v,\tau,\mathcal M_K,\{\psi_m\}$}
  \State \textbf{Initialization}
  \State $n\gets 0$
  ; Generate $N_u$,$N_v$ initial particle positions $S_{u,0}=\{\mathbf X^i_0\}$, $S_{v,0}=\{\mathbf Y^i_0\}$ via MH from $\rho_{u,0}=u_0/Z_u$.
  \State \textbf{Interpolate (fixed-$N$)}
  \State $\tilde\rho_{u,0} \gets \textsc{Interp}\big(S_{u,0},N_u,\mathcal M_K,\{\psi_m\}\big)$
  ; $\tilde\rho_{v,0} \gets \textsc{Interp}\big(S_{v,0},N_v,\mathcal M_K,\{\psi_m\}\big)$
  \State Append $\tilde u_0 = Z_u \tilde\rho_{u,0}$ to $\mathcal U$;\quad Append $\tilde v_0 = Z_v \tilde\rho_{v,0}$ to $\mathcal V$
  \While{$n\tau<T$}
    \State \textbf{SDE propagation}
    \State $b_u(\bx) \gets -\dfrac{Z_v}{Z_u\,\tilde\rho_{u,n}(\bx)}\,\chi\!\big(Z_u\tilde\rho_{u,n}(\bx)\big)\,\nabla\tilde\rho_{v,n}(\bx)$; $b_v(\bx) \gets 0$
    \State $S_{u,n,*} \gets \textsc{SDE}\big(S_{u,n},\tau,b_u,\sigma_u\big)$; $S_{v,n,*} \gets \textsc{SDE}\big(S_{v,n},\tau,b_v,\sigma_v\big)$
    \State \textbf{Interpolate at starred positions (fixed-$N$)}
    \State $\tilde\rho_u^* \gets \textsc{Interp}\big(S_{u,n,*},N_u,\mathcal M_K,\{\psi_m\}\big)$
    \State $\tilde\rho_v^* \gets \textsc{Interp}\big(S_{v,n,*},N_v,\mathcal M_K,\{\psi_m\}\big)$
    \State \textbf{Birth--Death}
    \State $r_u(\bx) \gets \dfrac{f_u\!\big(Z_u\tilde\rho_u^*(\bx),\,Z_v\tilde\rho_v^*(\bx)\big)}{\tilde\rho_u^*(\bx)}$; $r_v(\bx) \gets \dfrac{f_v\!\big(Z_u\tilde\rho_u^*(\bx),\,Z_v\tilde\rho_v^*(\bx)\big)}{\tilde\rho_v^*(\bx)}$
    \State $S_{u,n+1} \gets \textsc{Birth\_Death}\big(S_{u,n,*},\tau,r_u\big)$
    ; $S_{v,n+1} \gets \textsc{Birth\_Death}\big(S_{v,n,*},\tau,r_v\big)$
    \State \textbf{Refresh fields for output (fixed-$N$)}
    \State $\tilde\rho_{u,n+1} \gets \textsc{Interp}\big(S_{u,n+1},N_u,\mathcal M_K,\{\psi_m\}\big)$
    ; $\tilde\rho_{v,n+1} \gets \textsc{Interp}\big(S_{v,n+1},N_v,\mathcal M_K,\{\psi_m\}\big)$
    \State Append $\tilde u_{n+1}=Z_u\tilde\rho_{u,n+1}$ to $\mathcal U$;\quad Append $\tilde v_{n+1}=Z_v\tilde\rho_{v,n+1}$ to $\mathcal V$
    \State $n\gets n+1$
  \EndWhile
  \State \Return $\mathcal U,\ \mathcal V$
\EndProcedure
\end{algorithmic}
\end{algorithm}



% \begin{algorithm}
% \caption{Stochastic Branching Method for the KS Equation}\label{alg:cap}
% \begin{algorithmic}
% \Require Initial population size $N_{x,0}, N_{y,0}$, time step size $\tau$, current step $n = 0$, and Fourier mode truncation $K$.
% \Require Precomputed normalization constants $Z_u$ and $Z_v$ (either numerically or analytically).
% \State Generate $N_{x,0}$ and $N_{y,0}$ initial particle positions 
% $\{\bX^i_0\}$ and $\{\bY^i_0\}$ according to the distributions 
% $\rho_u$ and $\rho_v$, respectively, using the Metropolis--Hastings algorithm.
% \State Interpolate the particle positions $\bX^i_n$ and $\bY^i_n$ onto Fourier modes up to the truncation $K$ to obtain the Fourier representations $\tilde{\rho}_u$ and $\tilde{\rho}_v$.

% \While{$n \tau < T$} \Comment{Assuming $T \mid \tau$}
%     \State Update the positions of the particles:
%     \For{$i=1$ to $N_{x,n}$}
%         \[
%         \bX^i_{n,*} = \bX^i_n - \frac{Z_v}{Z_u \tilde{\rho}_u(\bX^i_n)} \chi(Z_u \tilde{\rho}_u(\bX^i_n)) \nabla \tilde{\rho}_v(\bX^i_n) \tau + \sqrt{2\tau} \xi^{i}_{u,n}
%         \]
%         where $\xi^{i}_{u,n} \sim \mathcal{N}(0, \mathbb{I})$.
%     \EndFor
%     \For{$i=1$ to $N_{y,n}$}
%         \[
%         \bY^i_{n,*} = \bY^i_n + \sqrt{2\tau} \xi^{i}_{v,n}
%         \]
%         where $\xi^{i}_{v,n} \sim \mathcal{N}(0, \mathbb{I})$.
%     \EndFor

%     \State Interpolate the updated positions $\bX^i_{n,*}$ and $\bY^i_{n,*}$ onto Fourier modes up to truncation $K$ to obtain the Fourier representations $\tilde{\rho}_u^*$ and $\tilde{\rho}_v^*$.

%     \For{$i=1$ to $N_{x,n}$}
%         \State Compute the branching rate $\tilde{r}_u^i = \frac{f_u(Z_u \rho_u(\bX^i_{n,*}), Z_v \rho_v(\bX^i_{n,*}))}{\rho_u(\bX^i_{n,*})}$.
%         \If{$\tilde{r}_u^i > 0$}
%             \State Duplicate $\bX^i_{n,*}$ with probability $1 - \exp(-\tilde{r}_u^i \tau)$.
%         \ElsIf{$\tilde{r}_u^i < 0$}
%             \State Kill $\bX^i_{n,*}$ with probability $1 - \exp(\tilde{r}_u^i \tau)$.
%         \EndIf
%     \EndFor

%     \For{$i=1$ to $N_{y,n}$}
%         \State Compute the branching rate $\tilde{r}_v^i = \frac{f_v(Z_u \rho_u(\bY^i_{n,*}), Z_v \rho_v(\bY^i_{n,*}))}{\rho_v(\bY^i_{n,*})}$.
%         \If{$\tilde{r}_v^i > 0$}
%             \State Duplicate $\bY^i_{n,*}$ with probability $1 - \exp(-\tilde{r}_v^i \tau)$.
%         \ElsIf{$\tilde{r}_v^i < 0$}
%             \State Kill $\bY^i_{n,*}$ with probability $1 - \exp(\tilde{r}_v^i \tau)$.
%         \EndIf
%     \EndFor

%     \State Update the population sizes: $N_{x,n+1}$ and $N_{y,n+1}$.
%     \State Increment the step: $n = n + 1$.
% \EndWhile
% \end{algorithmic}
% \end{algorithm}

% \section{Kolmogorov forward equation and }
% Recall (overdamped) Langevin diffusion is the solution of the following stochastic differential equation,
% where $\bX_t \in \mathbb R^d$ and .  The probability density function $\rho(t)$ of \eqref{equ:overd_langevin} solves the Kolmogorov forward equation on $\mathbb R^d$
% \begin{equation} \label{equ:FPE}
%     \partial_t  \rho =  \nabla \cdot ( \nabla \rho +\rho \bF).
% \end{equation}
% The main idea of the work is to use the branching process to approximate the evolution of the nonlinearity locally. Before combing with the Langevin dynamics, let us consider the pure branching process given by 
% \begin{equation}
%     \partial_t \rho =  V(x) \rho. 
% \end{equation}
% It can be viewed as an infinite system of ordinary differential equations with index $x\in \mathbb R^d$ without any spatial derivative. 
% If $V(t,x) > 0$, then particles located at $x$ undergo duplication, representing local exponential growth. Conversely, if $V(t,x)<0$, particles at 
% $x$ are removed, corresponding to exponential decay. In this way, the pure branching process encodes a pointwise multiplicative effect on the density, without involving any spatial interactions.
% Combining the Langevin dynamics with the branching process yields the following equation for the evolution of the probability density: 
% \begin{equation} \label{equ:FPE_BDL}
% \partial_t \rho = \nabla \cdot ( \nabla \rho + \rho \bF ) + V(x) \rho, 
% \end{equation} 
% where the first term represents diffusion and drift due to the Langevin process, and the second term captures local growth or decay of mass driven by the branching mechanism.

% The solution of \eqref{equ:FPE_BDL} can be approximated by the empirical measure $\mu^{N}_t$ of a collection of interacting particles $\{x_t^i\}_{i=1}^{N_t}$ evolving as followings $\mu^{N}_t = \frac{1}{N} \sum_{i=1}^{N_t} \delta_{x_t^i} $:
% \begin{enumerate}
%     \item[Step 1:] Between birth/death events, each particle $\bX^i_t$ diffuses independently according to \eqref{equ:overd_langevin}.
%     \item[Step 2:] Each particle also has an independent instantaneous birth-death rate $V(\bX_t^i).$
% \end{enumerate}
% Therefore, we can approximate the solution of \eqref{equ:FPE_BDL} by a density estimation.  Unlike the classical mesh-based method, this approach does not require spatial discretization. This is the key reason why this method is promising to overcome the curse of dimensionality in high-dimensional PDEs.


% \begin{proposition}
% Let $\mu^N_t$ be the empirical measure of particles defined above. Assume that $\mu_0^N \rightharpoonup \rho_0$ as $N\to \infty$. Then for all $t\in (0,\infty)$, we have $\mu^N_t \to \rho_t$, where $\rho_t$ solves \eqref{equ:FPE_BDL} with initial condition $\rho_0$
% \end{proposition}
% \begin{proof}
%     First, we notice that the evolution of $\mu^N_t$ is Markovian since the particles $\bX^i_t$ are, and these particles are interchangeable. Following the proof in \cite{rotskoff2019global}, we want to derive the generator and the backward Kolmogorov equation of the measure-valued Markov process. Let's define for some smooth functional $\Psi:\mathcal P  (\mathbb R ^d) \to \mathbb R $, the generator $\mathcal L_N$
%     \[
%     (\mathcal L_N \Psi ) (\mu^N) := \lim_{t\to 0^+} \frac{\mathbb E_0 \Psi ( \mu_t^N) - \Psi(\mu^N)}{t},
%     \]
%     where $\mathbb E_0 $ denotes the expectation of $\Psi(\mu^N_t)$ conditioned on that $\mu_0^N=\mu^N$. To compute the limit, notice that if particle $\bX_t^i$ is killed at time $t$, the changes this indeuced on $\mu^N_t$ is 
%     \[
%     \mu^N_t =  \mu_{t_-}^N - \frac{1}{N} \delta_{\bX^i}.
%     \]
%     Similarly if particle $\bX_t^i$ is duplicated at time $t$, the changes this indeuced on $\mu^N_t$ is 
%     \[
%     \mu^N_t =  \mu_{t_-}^N + \frac{1}{N} \delta_{\bX^i}.
%     \]
%     If we account for the rate at which these events occur, as well as the effect of the diffusion process, we can explicitly compute the generator 
%     \[
%     \begin{aligned}
%     (\mathcal L _N \Psi)(\mu^N) =& \frac{1}{N} \sum_i ^N \int \left( 
%     \bF(\bX^i) \cdot \nabla D_{\mu^N} \Psi(\bX^i)  
%     + \Delta D_{\mu^N} \Psi(\bX^i)
%     \right) \delta_{\bX^i}(\intd x) \\
%     &+ \frac{1}{N} \sum_{i=1}^N \int V(\bX^i) \delta_{\bX^i}(\intd x)\left(\Psi(\mu^N + \frac{1}{N} \text{sign} V(\bX^i)\delta_{\bX^i} ) - \Psi(\mu^N)\right),
%     \end{aligned}
%     \]
%     where the functional derivative $D_\mu \Psi$ is the function from $\mathbb R^d \to \mathbb R$ defined by that for any signed measure $\nu$ with $\int \nu(\intd x)= 0  $, 
%     \[
%     \lim_{\epsilon \to 0 } \frac{\Psi(\mu+\epsilon\nu)-\Psi(\mu)}{\epsilon} = \int D_\mu\Psi(x)\nu(\intd x).
%     \]
%     By the definition of the empirical measure $\mu^N$, the generator $\mathcal L_N$ can be rewritten as 
%     \[
%     \begin{aligned}
%     (\mathcal L _N \Psi)(\mu^N) =& \int \left( 
%     \bF(x) \cdot \nabla D_{\mu^N} \Psi(x)  
%     + \Delta D_{\mu^N} \Psi(x)
%     \right) \mu^N(\intd x)  \\
%     +&  \int V(x) \mu^N(\intd x) \left(\Psi(\mu^N + \frac{1}{N} \text{sign} V(x)\delta_{x} ) - \Psi(\mu^N)\right) ,
%     \end{aligned}
%     \]
% \end{proof}
% \section{Stochastic Branching method}
% In this work, we aim to extend the above framework to fully nonlinear PDEs. A natural strategy is to adopt an iterative scheme that treats the nonlinear equation as a linear one at each iteration, conditioned on the current estimate of the solution. Specifically, at each time $t$, the equation \eqref{equ:KS} can be approximated by the following linearized system: 
% \begin{equation}
% \begin{aligned}
%     \partial_t \rho_u &=  \nabla \cdot \left ( \bF_u(t,x) \rho_u +\nabla \rho_u \right) + V_u(t,x) \rho_u,\\
%     \partial_t \rho_v &=  \nabla \cdot \left ( \bF_v(t,x) \rho_v +\nabla \rho_v \right) + V_v(t,x) \rho_v,
% \end{aligned}
% \end{equation}
% where the drift and reaction terms are given by $\bF_u =  -\frac{\chi(u)}{u}\nabla v, \bF_v = 0, V_u = \frac{f_u(u,v)}{u}$ and $V_v= \frac{f_v(u,v)}{v}$. At each iteration, the fields $u$ and $v$ are reconstructed from the particle densities as  $u = C_u \rho_u$ and $v= C_v\rho_v$, where $C_u = \int u(0,\bx)\intd \bx$ and $C_v = \int v(0,\bx)\intd \bx$ are normalization constants.

\section{Numerical Examples}
In this section, we present numerical examples in two-dimensional space to verify the accuracy of our method.

\textbf{Case 1:} We first consider the AC equation, a widely used model for phase separation dynamics, to verify the accuracy of our scheme in handling nonlinear reaction systems. The governing equation is:
\begin{equation} \label{eq:ac_pde}
    u_t = D\Delta u + u - u^3, \quad \bx \in \Omega, t > 0,
\end{equation}
with the diffusion coefficient $D=0.01$. The computational domain is set to $\Omega = [0, 2\pi]^2$ with periodic boundary conditions imposed on all boundaries. 

The initial condition is given by:
\begin{equation} \label{eq:ac_ic}
    u_0(\bx) = \sin^2(x_1) \cos^2(x_2), \quad \bx = (x_1, x_2) \in \Omega.
\end{equation}
We solve the system numerically using our proposed method with Fourier projection $K=10$ and number of particle $N=400000$ up to a final time $T=0.8$. 
Figure~\ref{fig:ac_solution} shows the snapshots of the numerical solution $\tilde{u}$ at different time instances and its comparison with the finite difference method.
\begin{figure}
    \centering
    \includegraphics[width=0.8\linewidth]{figure/case_allen_cahn.pdf}
    \caption{Numerical solution of the AC equation \eqref{eq:ac_pde} at different time snapshots ($T=0, 0.2, 0.4, 0.6, 0.8$). The top row displays the reference solution computed by a finite difference method, while the bottom row shows the results from our proposed particle method.}
    \label{fig:ac_solution}
\end{figure}

\textbf{Case 2:} Next, we consider the following KS equation in a periodic domain $\Omega = [0, 2\pi] \times [0, 2\pi]$:
\begin{equation} \label{equ:case1}
\begin{aligned}
&u_t - \nabla \cdot \left(\nabla u - u \nabla v \right) = 0  , &\quad \bx \in \Omega , t > 0, \\
&v_t - \Delta v = u - v , & \quad \bx \in \Omega, t > 0.
\end{aligned}
\end{equation}
The system is subject to periodic boundary conditions, with initial conditions given by:
$$
u_0 (\bx) =  \sin^2(\bx_1) \cos^2(\bx_2), \quad v_0(\bx) =  \cos(\bx_1) + \cos(\bx_2) + 2.
$$
The temporal evolution of the total mass of $v$ is characterized by
$$
M(t) = \int v(t,\bx) \, \intd \bx.
$$
It can be shown that the rate of change of the mass is given by
$$
\dot{M}(t) = -M(t) + \int u(x,0) \, dx = -M(t) + \pi^2.
$$
We present the numerical results for $M(t)$, comparing our method with the exact solution in Figure~\ref{fig:total_mass}.

\begin{figure}[h]
\centering
\includegraphics[width=0.8\linewidth]{figure/total_mass.pdf}
\caption{Time evolution of the total mass of $v$ for different numbers of particles. The left panel shows the total mass over time, while the right panel displays the corresponding error in the total mass.}
\label{fig:total_mass}
\end{figure}

Additionally, we compare our particle-based method with finite difference method with a grid of $100 \times 100$ points. The result for $u(t)$ and $v(t)$ at $t = 0.0, 0.1, 0.2, 0.3, 0.4$ are shown in Figures~\ref{fig:fd_comparison_u} and~\ref{fig:fd_comparison_v}, respectively. Figure~\ref{fig:convergence_error} illustrates the reduction in relative $L^2$ error as the number of samples increases, showing a clear $1/\sqrt{N}$ convergence rate.

\begin{figure}[h]
\centering
\includegraphics[width=\linewidth]{figure/case1_u.pdf}
\caption{Comparison of the numerical approximation of $u$ obtained from the finite difference method with a $100 \times 100$ grid and the present particle-based method at $t = 0.0, 0.1, 0.2, 0.3, 0.4$.}
\label{fig:fd_comparison_u}
\end{figure}

\begin{figure}[h]
\centering
\includegraphics[width=\linewidth]{figure/case1_v.pdf}
\caption{Comparison of the numerical approximation of $v$ obtained from the finite difference method with a $100 \times 100$ grid and the present particle-based method at $t = 0.0, 0.1, 0.2, 0.3, 0.4$.}
\label{fig:fd_comparison_v}
\end{figure}

\begin{figure}[h]
\centering
\includegraphics[width=0.8\linewidth]{figure/case1_err.pdf}
\caption{Relative $L^2$ error of $u$ and $v$ with respect to the number of samples $N_0$. The reference solution is obtained with $N_0 = 320,000$.}
\label{fig:convergence_error}
\end{figure}

\textbf{Case 3:} Furthermore, we solve the system \eqref{equ:case1} with the initial conditions:
$$
u_0(\bx) = 840 \exp\left(-84 (\bx_1^2 + \bx_2^2)\right), \quad v_0(\bx) = 420 \exp\left(-42(\bx_1^2 + \bx_2^2)\right).
$$
%
Figure~\ref{fig:case2_u} shows the numerical solution $u$, which undergoes a blow-up phenomenon that is generally difficult to capture using the standard grid discretization. To assess the performance of our method, we compare its results against a finite difference simulation. Three resolution levels are analyzed: the particle method with $N = 40,000$, $160,000$, and $640,000$ is compared to the finite difference method with $200 \times 200$, $400 \times 400$, and $800 \times 800$ mesh points. 
%Our method remains stable with a timestep of $1 \times 10^{-7}$, whereas the finite difference method requires a smaller timestep of $1 \times 10^{-8}$ to avoid earlier blow up of the solution. 
Figure \ref{fig:comparison} shows the numerical solution at $t=5 \times 10^{-5}$. The present particle method enables us to accurately capture the aggregation dynamics and is robust with respect to the number of particles. In contrast, the solution obtained from the finite difference method exhibits pronounced mesh size sensitivity near the aggregation. While the performance of the grid-based approach can be improved by applying more sophisticated spatial discretization such as local discontinuous Galerkin method~\cite{li2017local}, mesh refinement is often required to track the aggregation. On the other hand, the present particle method is natural for modeling such strong nonlinear behaviors without using the adaptive mesh refinement.

% At a fixed reconstruction bandwidth,  the $L^2$ norm of the numerical solution $\tilde{u}$ for different particle counts $N_0$ exhibits similar growth pattern, which indicates robustness with respect to $N_0$, i.e., increasing the particle number primarily reduces the sampling variance but does not change the effective peak shape by the reconstruction. In contrast, the solution obtained from the local discontinuous Galerkin (LDG) method ~\cite{li2017local} exhibits pronounced mesh sensitivity near aggregation, with the observed $L^2$ growth depending strongly on the grid size.

\begin{figure}
\centering
\includegraphics[width=0.8\linewidth]{figure/case2_u.pdf}
\caption{Numerical approximations of $u$ at $t = 0, 5 \times 10^{-5}, 10^{-4}$, and $1.5 \times 10^{-4}$.}
\label{fig:case2_u}
\end{figure}
\begin{figure}
\centering
\includegraphics[width=0.85\linewidth]{figure/comparison.pdf}
\caption{Comparison of the numerical solution $u$ from the finite difference method and the present particle method at $t=5\times 10^{-5}$.}
\label{fig:comparison}
\end{figure}
% \begin{figure}
% \centering
% \includegraphics[width=0.5\linewidth]{figure/case2_L2.pdf}
% \caption{Time evolution of the $L^2$-norm of the numerical solution $\tilde{u}$ using different values of $N_0$.}
% \label{fig:case2_l2}
% \end{figure}

\textbf{Case 4:} Finally, we examine a fully nonlinear case with logistic growth for the cells, where the interaction between the particles is not known a priori. The system is given by:
\begin{equation}
\label{equ:nonliear}
\begin{aligned}
&u_t - \nabla \cdot \left(\nabla u - \frac{4u}{1 + u^2} \nabla v \right) = u(1 - u), & \quad \bx \in \Omega, t > 0, \\
&v_t - \Delta v = u - v, & \quad \bx \in \Omega, t > 0.
\end{aligned}
\end{equation}
The theoretical results of this model can be found in \cite{osaki2002exponential, wang2007classical, wrzosek2004global}. The initial conditions are specified as:
$$
    \begin{aligned}
        u_0(\bx) =&  \exp\left( (\bx_1 - 0.5\pi)^2 + (\bx_2 - 0.5\pi)^2 \right) + \exp\left( (\bx_1 - \pi)^2 + (\bx_2 - \pi)^2 \right) \\
        &+ \exp\left( (\bx_1 - 1.5\pi)^2 + (\bx_2 - 1.5\pi)^2 \right), \\
        v_0(\bx) =&  \exp\left( (\bx_1 - 0.4\pi)^2 + (\bx_2 - 0.4\pi)^2 \right) + \exp\left( (\bx_1 - 0.8\pi)^2 + (\bx_2 - 0.8\pi)^2 \right) \\
        &+ \exp\left( (\bx_1 - 1.2\pi)^2 + (\bx_2 - 1.2\pi)^2 \right) + \exp\left( (\bx_1 - 1.6\pi)^2 + (\bx_2 - 1.6\pi)^2 \right).
    \end{aligned}
$$

With the non-conservative form, both the birth and death of cells, as well as the chemical concentrations $f_u$ and $f_v$, are non-zero. The particle method proposed in \cite{wang2025novel} only accounts for advection and diffusion and therefore has limited capability in addressing these effects.  In contrast, the present method can naturally handle these additional complexities in the non-conservative form. 
Figures~\ref{fig:case3_u} and~\ref{fig:case3_v} show the time evolution of the solutions $u$ and $v$ obtained using both the particle-based and the finite difference methods.

\begin{figure}
\centering
\includegraphics[width=0.8\linewidth]{figure/case3_u.pdf}
\caption{Time evolution of the numerical solution $u$ computed using the particle-based and the finite difference methods.}
\label{fig:case3_u}
\end{figure}

\begin{figure}
\centering
\includegraphics[width=0.8\linewidth]{figure/case3_v.pdf}
\caption{Time evolution of the numerical solution $v$ computed using the particle-based and the finite difference methods.}
\label{fig:case3_v}
\end{figure}

\section{Discussion}
This work presents a stochastic particle-based approach for solving nonlinear non-conservative advection–diffusion–reaction equations. 
The evolution is split into two steps: an advection–diffusion step implemented by particle dynamics, and a local reaction step represented by a stochastic branching birth–death mechanism that provides a consistent temporal discretization of the underlying reaction dynamics. 
Unlike traditional particle or Lagrangian schemes, this formulation enables the present method to naturally capture the non-conservative evolution without an explicit form of the microscopic physical law. The resulting scheme preserves nonnegativity, captures concentration and blow-up robustly, and achieves an unbiased expectation with $O(1/N)$ Monte Carlo variance. Importantly, the algorithm is inherently parallel and does not require grid-based discretization.
%
Numerical results of the AC and the KS equations show the generality and robustness. In particular, the present method accurately captures the nonlinear behaviors such as aggregation and mass evolution of the KS equation without requiring adaptive mesh refinement, and is applicable to a broad class of transport and interaction problems. 

Looking forward, the current work lays the foundation for several promising extensions. From a numerical perspective, the present particle-based formulation is particularly well-suited for high-dimensional PDE problems, where grid-based methods often show limitions. This calls for efficient interpolation of empirical particle measures. In particular, the low-rank spectral approximations and adaptive reconstruction techniques offer promising directions ~\cite{tang2025solving, Peng_Yang_tensor_arxiv_2024}. Additionally, the present method can be extended for more general boundary conditions such as Dirichlet, Neumann, and Robin types, where proper stochastic killing, reflection, or mixed rules need to be developed at the particle level. Finally, while we provide convergence results for the scalar case, theoretical analysis of fully coupled multi-species systems requires further investigation. We leave these for future studies. 

% The framework presented here offers a fresh particle-based approach to solving reaction--diffusion--advection equations. Unlike traditional particle or Lagrangian schemes, our method is not restricted to conservation laws and is able to capture aggregation and focusing phenomena in the density variable with high fidelity. This highlights its potential as a useful tool for studying a broad class of nonlinear transport and interaction problems.

% From the theoretical perspective, the current work provides heuristic intuition primarily in the linear case. Developing a rigorous convergence theory in more general nonlinear settings remains a challenging but exciting direction for future research. On the numerical side, while the interpolation of measures poses difficulties in high-dimensional problems, this limitation also points to promising avenues for methodological innovation. In particular, Fourier interpolation combined with low-rank approximations~\cite{tang2025solving} could provide a powerful mechanism for scaling the algorithm efficiently to higher dimensions.


\section*{Acknowledgments}
The work is supported in part by the National Science Foundation under Grant DMS-2143739 and the ACCESS program through allocation MTH210005. 
The authors also acknowledge the support from the Institute for Cyber-Enabled Research at Michigan State University.


\bibliographystyle{cas-model2-names}
\bibliography{main.bib}

%% else use the following coding to input the bibitems directly in the
%% TeX file.

%% Refer following link for more details about bibliography and citations.
%% https://en.wikibooks.org/wiki/LaTeX/Bibliography_Management

% \begin{thebibliography}{00}

%% For numbered reference style
%% \bibitem{label}
%% Text of bibliographic item

% \bibitem{lamport94}
%   Leslie Lamport,
%   \textit{\LaTeX: a document preparation system},
%   Addison Wesley, Massachusetts,
%   2nd edition,
%   1994.

% \end{thebibliography}
\end{document}

\endinput
%%
%% End of file `elsarticle-template-num.tex'.